\chapter*{CHAPITRE III (suite)\\[1.2ex]
\Large ÉTUDE COHOMOLOGIQUE DES FAISCEAUX COHÉRENTS}

\section{Foncteurs « Tor » locaux et globaux; formule de Künneth}
\label{section:III.6.fr}

\subsection{Introduction}
\label{subsection:III.6.1.fr}

\oldpage[III]{137}

\begin{env}[6.1.1]
Soient $f:X\to Y$ un morphisme de préschémas, $\mathscr F$ un
$\mathscr O_X$-Module quasi-cohérent. Dans l'étude des « images directes
supérieures » $\mathrm R^n f_*(\mathscr F)$, on est amené à considérer le
problème général suivant : étant donné un morphisme « changement de base »
$g:Y'\to Y$, on pose
\[
  X'=X_{(Y')}=X\times_Y Y',\qquad
  \mathscr F'=\mathscr F\otimes_Y\mathscr O_{Y'},\qquad
  f'=f_{(Y')}:X'\to Y',
\]
et l'on se propose d'avoir des renseignements sur les images directes
supérieures $\mathrm R^n f'_*(\mathscr F')$ (en supposant connus les
$\mathrm R^n f_*(\mathscr F)$). On voit aisément (cf. par exemple (7.7.2))
que l'on est ramené à étudier les variations de
$\mathrm R^n f_*(\mathscr F\otimes_{\mathscr O_Y}\mathscr G)$ pour un
$\mathscr O_Y$-Module quasi-cohérent variable $\mathscr G$, autrement dit le
foncteur
\[
  \mathscr G\longmapsto
  \mathrm R^n f_*(\mathscr F\otimes_{\mathscr O_Y}\mathscr G).
\]
Si $\mathscr F$ est plat sur $Y$, le foncteur
$\mathscr G\mapsto\mathscr F\otimes_{\mathscr O_Y}\mathscr G$ est exact
(0\textsubscript{I}, 6.7.4) et par suite le foncteur composé
$\mathscr G\mapsto\mathrm R^n f_*(\mathscr F\otimes_{\mathscr O_Y}\mathscr G)$
est encore un foncteur cohomologique. Mais il n'en est plus de même dans le
cas général; pour pouvoir appliquer les méthodes cohomologiques, on est
conduit à substituer à
$\mathscr G\mapsto\mathrm R^n f_*(\mathscr F\otimes_{\mathscr O_Y}\mathscr G)$
d'autres foncteurs, qui cette fois sont toujours des foncteurs
cohomologiques. Ces foncteurs, qui généralisent les foncteurs « Tor » de la
théorie des modules, sont définis aux n\textsuperscript{os} 6.3 à 6.7; il y a
d'ailleurs deux telles généralisations, l'une « locale » et l'autre
« globale », reliées par des suites spectrales qui seront discutées au
n\textsuperscript{o} 6.7; comme application de ces suites spectrales, on
obtient en particulier, sous certaines conditions, une « formule de Künneth »
exprimant
\[
  \mathrm R^n(f_1\times f_2)_*(\mathscr F_1\otimes_Y\mathscr F_2)
\]
à l'aide des images directes supérieures
$\mathrm R^p f_{1*}(\mathscr F_1)$ et $\mathrm R^q f_{2*}(\mathscr F_2)$.
D'autres suites spectrales (6.8) généralisent les suites spectrales
d'associativité du foncteur « Tor » de modules; enfin, le problème du
changement de base conduit lui aussi à des suites spectrales (6.9).
\end{env}

\begin{env}[6.1.2]
En outre, on constate (6.10) que les foncteurs cohomologiques
$\mathscr G\mapsto\mathscr T_\bullet(\mathscr G)$ ainsi définis sont
localement (sur $Y$) du type
\[
  \mathscr G\longmapsto
  \mathscr H_\bullet(\mathscr L_\bullet\otimes_Y\mathscr G),
\]
où $\mathscr L_\bullet$ est un complexe de $\mathscr O_Y$-modules localement
libres (défini à une homotopie près) et $\mathscr H_\bullet$ l'homologie. Il y
a alors intérêt à oublier la situation particulière qui a donné naissance à
$\mathscr T_\bullet$, et à étudier de

\oldpage[III]{138}

façon générale les foncteurs de la forme précédente
$\mathscr G\mapsto\mathscr H_\bullet(\mathscr L_\bullet\otimes_Y\mathscr G)$
(où l'on fait en outre le cas échéant des hypothèses de finitude appropriées
sur les $\mathscr L_i$ ou les $\mathscr H_i(\mathscr L_\bullet)$) : c'est ce
qui est l'objet du § 7, dont la lecture, pour l'essentiel, est indépendante du
§ 6. Les propriétés les plus importantes de ces foncteurs concernent les
propriétés d'exactitude d'une composante $\mathscr T_i$ de
$\mathscr T_\bullet$; on donnera divers critères permettant d'établir de
telles propriétés; comme application, on obtiendra des conditions permettant
d'affirmer (avec les notations de (6.1.1)) que le foncteur
$\mathscr G\mapsto\mathrm R^n f_*(\mathscr F\otimes_{\mathscr O_Y}\mathscr G)$
est exact (ce qu'on exprimera en disant que $\mathscr F$ est
cohomologiquement plat sur $Y$ en dimension $n$). Une autre propriété
importante pour les composantes $\mathscr T_i$ de
$\mathscr T_\bullet(\mathscr G)=
\mathscr H_\bullet(\mathscr L_\bullet\otimes_Y\mathscr G)$ est une propriété
de semi-continuité de la fonction
$y\mapsto\dim_{k(y)}(T_i(k(y)))$; lorsque $\mathscr T_i$ est exact, cette
propriété est remplacée par une propriété de continuité, la réciproque étant
d'ailleurs vraie d'après Grauert lorsque $Y$ est réduit (7.8.4).
\end{env}

\begin{env}[6.1.3]
Dans les §§ 6 et 7, nous avons systématiquement fait usage de
l'hypercohomologie, en prenant partout comme arguments des complexes de
faisceaux au lieu de faisceaux, bien que la nécessité de ce point de vue
n'apparaisse que dans des chapitres ultérieurs. Le formalisme cohomologique
développé à cette occasion deviendra d'ailleurs plus transparent dans le
chapitre de ce Traité qui sera consacré à la mise au point d'une algèbre des
foncteurs cohomologiques de faisceaux cohérents, incluant le formalisme de la
dualité. Mais cela demandera des développements qui sortent du cadre du
présent chapitre.
\end{env}

\begin{env}[6.1.4]
Pour abréger, étant donnés deux complexes $K^\bullet$, $K'^\bullet$ dans une
catégorie abélienne $C$, nous dirons qu'un morphisme de complexes
$f:K^\bullet\to K'^\bullet$ est un homotopisme s'il existe un morphisme
$g:K'^\bullet\to K^\bullet$ tel que les morphismes composés $f\circ g$ et
$g\circ f$ soient tous deux homotopes à l'identité (par abus de langage,
lorsqu'il existe un tel homotopisme, on dira aussi que $K^\bullet$ et
$K'^\bullet$ sont homotopes). Lorsqu'on peut définir l'hypercohomologie d'un
foncteur covariant additif $T$ de $C$ dans une catégorie abélienne $C'$, par
rapport à un complexe de $C$ (0, 11.4.3), il est immédiat qu'un homotopisme
$K^\bullet\to K'^\bullet$ de complexes de $C$ définit canoniquement un
isomorphisme
\[
  \mathrm R^\bullet T(K^\bullet)\xrightarrow{\sim}
  \mathrm R^\bullet T(K'^\bullet)
\]
pour l'hypercohomologie (loc. cit.).
\end{env}

\subsection{Hypercohomologie des complexes de modules sur un préschéma}
\label{subsection:III.6.2.fr}

\begin{env}[6.2.1]
Soient $X$ un préschéma,
$\mathscr K^\bullet=(\mathscr K^i)_{i\in\mathbf Z}$ un complexe de
$\mathscr O_X$-Modules dont l'opérateur de dérivation est de degré $+1$.
Rappelons que pour tout morphisme $f:X\to Y$ de préschémas, on a défini
(0, 12.4.1) les $\mathscr O_Y$-Modules d'hypercohomologie
$\mathscr H^n(f,\mathscr K^\bullet)$ (aussi notés
$\mathscr H_f^n(\mathscr K^\bullet)$ ou
$\mathrm R^n f_*(\mathscr K^\bullet)$) pour tout $n\in\mathbf Z$;
l'hypercohomologie $\mathscr H^\bullet(f,\mathscr K^\bullet)$ est
l'aboutissement des deux foncteurs spectraux
$'\mathscr E(f,\mathscr K^\bullet)$ et
$''\mathscr E(f,\mathscr K^\bullet)$, dont les termes $\mathscr E_2$ sont
donnés par
\begin{align}
  '\mathscr E_2^{pq}
    &=\mathscr H^p\bigl(\mathscr H^q(f,\mathscr K^\bullet)\bigr),
    \tag{6.2.1.1}\label{III.6.2.1.1.fr}\\
  ''\mathscr E_2^{pq}
    &=\mathscr H^p\bigl(f,\mathscr H^q(\mathscr K^\bullet)\bigr)
      =\mathrm R^p f_*\bigl(\mathscr H^q(\mathscr K^\bullet)\bigr).
    \tag{6.2.1.2}\label{III.6.2.1.2.fr}
\end{align}

\oldpage[III]{139}

où $\mathscr H^q(f,\mathscr K^\bullet)$ est le complexe dont le composant de
degré $i$ est
$\mathscr H^q(f,\mathscr K^i)=\mathrm R^q f_*(\mathscr K^i)$ (loc. cit.).
Rappelons aussi que lorsque $Y$ est réduit à un point, on note
l'hypercohomologie correspondante $\mathbf H^\bullet(X,\mathscr K^\bullet)$
(qui est formée de modules sur $\Gamma(X,\mathscr O_X)$ indépendants du
préschéma ponctuel $Y$ considéré); lorsque $Y=X$ et $f=1_X$, on a
$\mathscr H^n(f,\mathscr K^\bullet)=\mathscr H^n(\mathscr K^\bullet)$
(cohomologie du complexe $\mathscr K^\bullet$); lorsque
$\mathscr K^i=0$, sauf pour $i=i_0$, on a
\[
  \mathscr H^n(f,\mathscr K^\bullet)
    =\mathrm R^{n-i_0}f_*(\mathscr K^{i_0}).
\]
La suite spectrale $'\mathscr E(f,\mathscr K^\bullet)$ est toujours
régulière; les deux suites spectrales sont birégulières lorsque
$\mathscr K^\bullet$ est limité inférieurement (0, 12.4.1).

Tout homotopisme $h:\mathscr K^\bullet\to\mathscr K'^\bullet$ de complexes de
$\mathscr O_X$-Modules (6.1.4) donne un isomorphisme
$\mathscr H^\bullet(f,\mathscr K^\bullet)\xrightarrow{\sim}
\mathscr H^\bullet(f,\mathscr K'^\bullet)$ pour l'hypercohomologie. Il en est
de même lorsqu'on suppose seulement que
$\mathscr H^\bullet(h):\mathscr H^\bullet(\mathscr K^\bullet)\to
\mathscr H^\bullet(\mathscr K'^\bullet)$ est un isomorphisme et que
$\mathscr K^\bullet$ et $\mathscr K'^\bullet$ sont limités inférieurement,
comme il résulte aussitôt de (0, 11.1.5) appliqué à la suite spectrale
(6.2.1.2) et à l'analogue pour $\mathscr K'^\bullet$. Enfin, pour tout
recouvrement ouvert $\mathfrak U=(U_\alpha)$ de $X$, on a aussi défini
(0, 12.4.5) l'hypercohomologie
$\mathbf H^\bullet(\mathfrak U,\mathscr K^\bullet)$ comme la cohomologie du
bicomplexe $C^\bullet(\mathfrak U,\mathscr K^\bullet)$ (dont le composant
d'indices $(i,j)$ est par définition
$C^i(\mathfrak U,\mathscr K^j)$); les
$\mathbf H^n(\mathfrak U,\mathscr K^\bullet)$ sont encore des modules sur
$\Gamma(X,\mathscr O_X)$.
\end{env}

\begin{proposition}[6.2.2]
Soient $X$ un schéma, $\mathfrak U=(U_\alpha)$ un recouvrement de $X$ par des
ouverts affines. Pour tout complexe $\mathscr K^\bullet$ de
$\mathscr O_X$-Modules quasi-cohérents, les modules d'hypercohomologie
$\mathbf H^\bullet(X,\mathscr K^\bullet)$ et
$\mathbf H^\bullet(\mathfrak U,\mathscr K^\bullet)$ sont canoniquement
isomorphes.
\end{proposition}

\begin{proof}
En effet, toute intersection finie $V$ d'ouverts du recouvrement
$\mathfrak U$ est affine (I, 5.5.6), donc
$\mathrm H^q(V,\mathscr K^i)=0$ pour tout $i$ et tout $q>0$ (1.3.1); la
proposition est donc un cas particulier de (0, 12.4.7).
\end{proof}

\begin{proposition}[6.2.3]
Soit $f:X\to Y$ un morphisme quasi-compact et séparé de préschémas. Pour tout
complexe $\mathscr K^\bullet$ de $\mathscr O_X$-Modules quasi-cohérents, les
$\mathscr O_Y$-Modules $\mathscr H^n(f,\mathscr K^\bullet)$ sont
quasi-cohérents.
\end{proposition}

\begin{proof}
Comme les $\mathscr H^q(f,\mathscr K^i)=\mathrm R^q f_*(\mathscr K_i)$ sont
des $\mathscr O_Y$-Modules quasi-cohérents (1.4.10), il en est de même de
$'\mathscr E_2^{pq}$, qui, d'après (6.2.1.1), est quotient d'un noyau
d'homomorphisme de Modules quasi-cohérents par une image d'un tel
homomorphisme (I, 4.1.1). Pour la même raison, tous les
$\mathscr O_Y$-Modules $'\mathscr E_r^{pq}$,
$B_k('\mathscr E_2^{pq})$, $Z_k('\mathscr E_2^{pq})$ de la première suite
spectrale sont quasi-cohérents. La régularité de la suite spectrale
$'\mathscr E(f,\mathscr K^\bullet)$ entraîne que
$Z_\infty('\mathscr E_2^{pq})$ est égal à un des
$Z_k('\mathscr E_2^{pq})$, donc est quasi-cohérent, et il en est de même de
$B_\infty('\mathscr E_2^{pq})=\varinjlim_k B_k('\mathscr E_2^{pq})$
(0, 11.2.4 et I, 4.1.1); les $'\mathscr E_\infty^{pq}$ sont donc aussi
quasi-cohérents. La suite spectrale précédente étant régulière, la filtration
des $F^p(\mathscr H^n(f,\mathscr K^\bullet))$ est discrète et exhaustive;
autrement dit, le $\mathscr O_Y$-Module
$\mathscr H^n(f,\mathscr K^\bullet)$ est réunion d'une suite croissante
$(\mathscr G_k)_{k\geq0}$ de $\mathscr O_Y$-Modules telle que
$\mathscr G_0=0$ et que chaque $\mathscr G_k/\mathscr G_{k-1}$ soit égal à un
des $\mathscr O_Y$-Modules $'\mathscr E_\infty^{pq}$, donc soit
quasi-cohérent. Par récurrence sur $k$, on en déduit que les $\mathscr G_k$
sont quasi-cohérents (1.4.17), et comme
$\mathscr H^n(f,\mathscr K^\bullet)=\varinjlim_k\mathscr G_k$, la proposition
est démontrée (I, 4.1.1).
\end{proof}

\oldpage[III]{140}

\begin{corollary}[6.2.4]
Sous les hypothèses de (6.2.3), pour tout ouvert affine $V$ de $Y$,
l'homomorphisme canonique
\[
  \mathbf H^n(f^{-1}(V),\mathscr K^\bullet)
  \longrightarrow
  \Gamma(V,\mathscr H^n(f,\mathscr K^\bullet))
  \tag{6.2.4.1}\label{III.6.2.4.1.fr}
\]
est bijectif pour tout $n\in\mathbf Z$.
\end{corollary}

\begin{proof}
La démonstration est la même que celle de (1.4.11), en utilisant (6.2.2),
remplaçant $\mathscr F$ par $\mathscr K^\bullet$,
$\mathscr K^\bullet$ par
$f_*(C^\bullet(\mathfrak U,\mathscr K^\bullet))$,
$\mathscr H^\bullet(\mathscr K^\bullet)$ par
$\mathscr H^\bullet(f,\mathscr K^\bullet)$, et notant que ce dernier est un
$\mathscr O_Y$-Module quasi-cohérent par (6.2.3).
\end{proof}

\begin{proposition}[6.2.5]
Soient $Y$ un préschéma localement noethérien, $f:X\to Y$ un morphisme
propre, $\mathscr K^\bullet$ un complexe de $\mathscr O_X$-Modules tel que les
$\mathscr O_X$-Modules $\mathscr H^q(\mathscr K^\bullet)$ soient cohérents.
Alors les $\mathscr O_Y$-Modules $\mathscr H^n(f,\mathscr K^\bullet)$ sont
cohérents.
\end{proposition}

\begin{proof}
La question étant locale sur $Y$, on peut se borner au cas où $Y$ est
noethérien et affine, et il s'agit donc, en vertu de (6.2.4), de prouver que
les $\mathbf H^n(X,\mathscr K^\bullet)$ sont des
$\Gamma(Y,\mathscr O_Y)$-modules de type fini. On a alors
$\mathbf H^\bullet(X,\mathscr K^\bullet)=
\mathbf H^\bullet(\mathfrak U,\mathscr K^\bullet)$ (6.2.2), où l'on peut
supposer que $\mathfrak U$ est fini, puisque $X$ est quasi-compact. Les
cochaînes de chaque complexe $C^\bullet(\mathfrak U,\mathscr K^j)$ étant
alternées par définition, il y a un entier $r>0$ tel que
$C^i(\mathfrak U,\mathscr K^j)=0$ pour $i<0$ et $i>r$; on en conclut
(0, 11.3.3) que les deux suites spectrales du bicomplexe
$C^\bullet(\mathfrak U,\mathscr K^\bullet)$ sont birégulières. Comme les
intersections des ensembles de $\mathfrak U$ sont des ouverts affines
(I, 5.5.6), chaque foncteur
$\mathscr F\mapsto C^i(\mathfrak U,\mathscr F)$ est exact dans la catégorie
des $\mathscr O_X$-Modules quasi-cohérents; donc
$\mathrm H_I^q(C^i(\mathfrak U,\mathscr K^\bullet))
=C^i(\mathfrak U,\mathscr H^q(\mathscr K^\bullet))$, et les termes
$E_2$ de la seconde suite spectrale de
$C^\bullet(\mathfrak U,\mathscr K^\bullet)$ sont donnés (0, 11.3.2) par
\[
  ''E_2^{pq}
  =\mathrm H^p(C^\bullet(\mathfrak U,\mathscr H^q(\mathscr K^\bullet)))
  =\mathbf H^p(\mathfrak U,\mathscr H^q(\mathscr K^\bullet))
  =\mathrm H^p(X,\mathscr H^q(\mathscr K^\bullet))
\]
en vertu de (1.4.1); puisque $f$ est propre, ce sont des
$\Gamma(Y,\mathscr O_Y)$-modules de type fini (3.2.1). La suite spectrale
$''E(C^\bullet(\mathfrak U,\mathscr K^\bullet))$ étant birégulière, on en
déduit bien que les $\mathbf H^n(X,\mathscr K^\bullet)$ sont des
$\Gamma(Y,\mathscr O_Y)$-modules de type fini (0, 11.1.8).

A fortiori, si $\mathscr K^\bullet$ est un complexe de
$\mathscr O_X$-Modules cohérents, les $\mathscr O_Y$-Modules
$\mathscr H^n(f,\mathscr K^\bullet)$ sont cohérents sous les hypothèses de
(6.2.5) relatives à $Y$ et $f$ (0\textsubscript{I}, 5.3.4).
\end{proof}

\begin{env}[6.2.6]
L'hypercohomologie $\mathscr H^\bullet(f,\mathscr K^\bullet)$ est un foncteur
cohomologique dans la catégorie des complexes de $\mathscr O_X$-Modules
limités inférieurement (0, 12.4.4). C'est un foncteur cohomologique dans la
catégorie de tous les complexes de $\mathscr O_X$-Modules lorsque le
morphisme $f$ est quasi-compact et l'espace sous-jacent à $X$ localement
noethérien : en effet, il résulte alors de (G, II, 3.10.1) que $f_*$ permute
aux limites inductives (la question étant locale sur $Y$), et l'on peut
appliquer (0, 11.5.2).

Enfin, si $f$ est séparé, $\mathscr H^\bullet(f,\mathscr K^\bullet)$ est un
foncteur cohomologique dans la catégorie des complexes de
$\mathscr O_X$-Modules quasi-cohérents. C'est immédiat lorsque $Y$ est
affine, car alors $X$ est un schéma, donc, en vertu de l'isomorphisme
canonique (6.2.2), on est ramené à voir que
$\mathscr K^\bullet\mapsto\mathbf H^\bullet(\mathfrak U,\mathscr K^\bullet)$
est un foncteur cohomologique dans la catégorie des complexes de
$\mathscr O_X$-Modules quasi-cohérents, ce qui est immédiat puisque le
foncteur $\mathscr K^\bullet\mapsto
C^\bullet(\mathfrak U,\mathscr K^\bullet)$ est exact dans cette catégorie
(I, 1.3.7). Dans le cas général, pour tout ouvert affine $V$

\oldpage[III]{141}

de $Y$, $f^{-1}(Y)$ est un schéma, et pour appliquer ce qui précède, il
suffit de vérifier que pour une suite exacte
\[
  0\longrightarrow\mathscr K'^\bullet
  \longrightarrow\mathscr K^\bullet
  \longrightarrow\mathscr K''^\bullet
  \longrightarrow0
\]
de complexes de $\mathscr O_X$-Modules quasi-cohérents, l'homomorphisme
\[
  \partial:
  \mathscr H^n(f,\mathscr K''^\bullet)|V
  \longrightarrow
  \mathscr H^{n+1}(f,\mathscr K'^\bullet)|V
\]
ne dépend pas du recouvrement ouvert affine $\mathfrak U$ de $f^{-1}(V)$
utilisé pour le définir. Mais cela résulte de ce que, si $\mathfrak U'$ est
un recouvrement ouvert affine plus fin que $\mathfrak U$, le diagramme
\[
  \xymatrix@C=34pt@R=11pt{
    \mathbf H^\bullet(\mathfrak U,\mathscr K^\bullet)
      \ar[rrd]^-{\sim}\ar[dd]^{\iota}
    && \\
    &&
    \mathbf H^\bullet(f^{-1}(V),\mathscr K^\bullet)
    \\
    \mathbf H^\bullet(\mathfrak U',\mathscr K^\bullet)
      \ar[rru]_-{\sim}
    &&
  }
\]
d'isomorphismes canoniques est commutatif, ainsi que le diagramme
\[
  \xymatrix@C=46pt@R=25pt{
    \mathbf H^n(\mathfrak U,\mathscr K''^\bullet)
      \ar[r]^-{\partial}\ar[d]^{\iota}
    &
    \mathbf H^{n+1}(\mathfrak U,\mathscr K'^\bullet)
      \ar[d]^{\iota}
    \\
    \mathbf H^n(\mathfrak U',\mathscr K''^\bullet)
      \ar[r]_-{\partial}
    &
    \mathbf H^{n+1}(\mathfrak U',\mathscr K'^\bullet).
  }
\]

Lorsque l'une des conditions précédentes est remplie et que
$\mathscr K^\bullet\to\mathscr K'^\bullet$ est un homotopisme (6.1.4),
l'isomorphisme correspondant
$\mathscr H^\bullet(f,\mathscr K^\bullet)\xrightarrow{\sim}
\mathscr H^\bullet(f,\mathscr K'^\bullet)$ est alors un isomorphisme de
$\partial$-foncteurs (0, 11.4.4).
\end{env}

\begin{env}[6.2.7]
Tout ce qui précède s'applique naturellement sans changement (sinon de
notations) à un complexe $\mathscr K_\bullet$ de $\mathscr O_X$-Modules
quasi-cohérents dont l'opérateur de dérivation est de degré $-1$; il suffit
de considérer le complexe $\mathscr K^\bullet=(\mathscr K^i)$ où
$\mathscr K^i=\mathscr K_{-i}$ pour tout $i\in\mathbf Z$.
\end{env}

\subsection{Hypertor de deux complexes de modules}
\label{subsection:III.6.3.fr}

\begin{env}[6.3.1]
Soient $A$ un anneau commutatif, $P_\bullet$, $Q_\bullet$ deux complexes de
$A$-modules dont les opérateurs de dérivation sont de degré $-1$; soit
$L_{\bullet\bullet}$ (resp. $M_{\bullet\bullet}$) une résolution projective
de Cartan-Eilenberg de $P_\bullet$ (resp. $Q_\bullet$) (0, 11.6.1);
$L_{\bullet\bullet}\otimes_A M_{\bullet\bullet}$ est alors (pour la somme
des premiers degrés et la somme des seconds degrés) un bicomplexe (à
opérateurs de dérivation de degré $-1$), dont l'homologie
$H_\bullet(L_{\bullet\bullet}\otimes_A M_{\bullet\bullet})$ ne dépend pas des
résolutions de Cartan-Eilenberg $L_{\bullet\bullet}$,
$M_{\bullet\bullet}$ choisies, et est par définition l'hyperhomologie du
bifoncteur $P_\bullet\otimes_A Q_\bullet$ en $P_\bullet$ et $Q_\bullet$
(0, 11.6.5). Nous poserons par définition
\[
  \mathbf{Tor}_n^A(P_\bullet,Q_\bullet)
  =H_n(L_{\bullet\bullet}\otimes_A M_{\bullet\bullet})
  \tag{6.3.1.1}\label{III.6.3.1.1.fr}
\]

\oldpage[III]{142}

et nous dirons que cet $A$-module est l'hypertor d'indice $n$ des deux
complexes $P_\bullet$, $Q_\bullet$. On sait que dans la catégorie des
complexes de $A$-modules limités inférieurement, les
$\mathbf{Tor}_n^A(P_\bullet,Q_\bullet)$ forment un bifoncteur homologique en
$P_\bullet$, $Q_\bullet$ (0, 11.6.5). En outre :
\end{env}

\begin{proposition}[6.3.2]
Le bifoncteur $\mathbf{Tor}_\bullet^A(P_\bullet,Q_\bullet)$ est
l'aboutissement commun de deux bifoncteurs spectraux
$'E(P_\bullet,Q_\bullet)$, $''E(P_\bullet,Q_\bullet)$, dont les termes $E_2$
sont
\begin{align}
  'E^2_{pq}
    &=H_p\bigl(\operatorname{Tor}_q^A(P_\bullet,Q_\bullet)\bigr),
    \tag{6.3.2.1}\label{III.6.3.2.1.fr}\\
  ''E^2_{pq}
    &=\bigoplus_{q'+q''=q}
      \operatorname{Tor}_p^A\bigl(H_{q'}(P_\bullet),H_{q''}(Q_\bullet)\bigr).
    \tag{6.3.2.2}\label{III.6.3.2.2.fr}
\end{align}
où, dans (6.3.2.1), $\operatorname{Tor}_q^A(P_\bullet,Q_\bullet)$ désigne le
bicomplexe formé des $A$-modules $\operatorname{Tor}_q^A(P_i,Q_j)$. La suite
spectrale (6.3.2.2) est toujours régulière; si $P_\bullet$ et $Q_\bullet$
sont limités inférieurement, ou si $A$ est de dimension cohomologique finie,
les deux suites spectrales (6.3.2.1) et (6.3.2.2) sont birégulières.
\end{proposition}

\begin{proof}
Cela résulte de (0, 11.6.5), car lorsque $A$ est de dimension cohomologique
finie $n$, tout $A$-module admet une résolution projective de longueur $n$
(M, VI, 2.1).
\end{proof}

\begin{corollary}[6.3.3]
Soient $P'_\bullet$, $Q'_\bullet$ deux complexes de $A$-modules,
$u:P_\bullet\to P'_\bullet$, $v:Q_\bullet\to Q'_\bullet$ deux homomorphismes
de complexes. Si les homomorphismes
$H_\bullet(u):H_\bullet(P_\bullet)\to H_\bullet(P'_\bullet)$,
$H_\bullet(v):H_\bullet(Q_\bullet)\to H_\bullet(Q'_\bullet)$ déduits
respectivement de $u$ et $v$ sont bijectifs, alors l'homomorphisme
$\mathbf{Tor}_\bullet^A(P_\bullet,Q_\bullet)\to
\mathbf{Tor}_\bullet^A(P'_\bullet,Q'_\bullet)$ déduit de $u$ et $v$ est
bijectif.
\end{corollary}

\begin{proof}
En effet, l'homomorphisme de suites spectrales
$''E(P_\bullet,Q_\bullet)\to''E(P'_\bullet,Q'_\bullet)$ déduit de $u$ et $v$
est alors un isomorphisme pour les termes $E_2$ et la conclusion résulte de
ce que ces suites sont régulières en vertu de (6.3.2) (0, 11.1.5).
\end{proof}

\begin{proposition}[6.3.4]
Soient $P_\bullet$, $Q_\bullet$ deux complexes de $A$-modules, limités
inférieurement. Soit $L_{\bullet\bullet}$ (resp. $M_{\bullet\bullet}$) un
bicomplexe formé de $A$-modules plats, tel que pour tout $i$,
$L_{i,\bullet}$ (resp. $M_{i,\bullet}$) soit une résolution de $P_i$
(resp. $Q_i$). On a alors des isomorphismes canoniques
\[
  \mathbf{Tor}_\bullet^A(P_\bullet,Q_\bullet)
  \xrightarrow{\sim}H_\bullet(L_{\bullet\bullet}\otimes_AQ_\bullet)
  \xrightarrow{\sim}H_\bullet(P_\bullet\otimes_AM_{\bullet\bullet})
  \xrightarrow{\sim}H_\bullet(L_{\bullet\bullet}\otimes_AM_{\bullet\bullet})
  \tag{6.3.4.1}\label{III.6.3.4.1.fr}
\]
\end{proposition}

\begin{proof}
Cela résulte de (0, 11.6.5, (ii) et (iii)) et de la définition des
$A$-modules plats.
\end{proof}

\begin{namedtheorem}[6.3.5]{Remarques}
(i) Avec les notations de (6.3.1), les bicomplexes
$L_{\bullet\bullet}\otimes_A M_{\bullet\bullet}$ et
$M_{\bullet\bullet}\otimes_A L_{\bullet\bullet}$ sont canoniquement
isomorphes, d'où un isomorphisme canonique
$\mathbf{Tor}_\bullet^A(P_\bullet,Q_\bullet)\xrightarrow{\sim}
\mathbf{Tor}_\bullet^A(Q_\bullet,P_\bullet)$.

(ii) Si $F$ et $G$ sont deux $A$-modules, $P_\bullet$ et $Q_\bullet$ les
complexes de $A$-modules réduits à $F$ et $G$ respectivement en degré 0 et
nuls dans les autres degrés, alors deux résolutions projectives $L_\bullet$,
$M_\bullet$ de $F$ et $G$ respectivement peuvent être considérées comme des
résolutions de Cartan-Eilenberg de $P_\bullet$ et $Q_\bullet$ en les
complétant par des zéros. On a par suite dans ce cas
$\mathbf{Tor}_\bullet^A(P_\bullet,Q_\bullet)
=\operatorname{Tor}_\bullet^A(F,G)$.
\end{namedtheorem}

\begin{proposition}[6.3.6]
Soient $(P_\bullet^\lambda)$, $(Q_\bullet^\mu)$ deux systèmes inductifs
filtrants de complexes de $A$-modules; on a un isomorphisme canonique
\[
  \varinjlim_{\lambda,\mu}
  \mathbf{Tor}_\bullet^A(P_\bullet^\lambda,Q_\bullet^\mu)
  \xrightarrow{\sim}
  \mathbf{Tor}_\bullet^A
  \left(\varinjlim_\lambda P_\bullet^\lambda,
        \varinjlim_\mu Q_\bullet^\mu\right).
  \tag{6.3.6.1}\label{III.6.3.6.1.fr}
\]
\end{proposition}

\begin{proof}
Posons $P_\bullet=\varinjlim_\lambda P_\bullet^\lambda$,
$Q_\bullet=\varinjlim_\mu Q_\bullet^\mu$; par fonctorialité, il est clair que
les $\mathbf{Tor}_\bullet^A(P_\bullet^\lambda,Q_\bullet^\mu)$ forment un
système inductif et que les applications
$\mathbf{Tor}_\bullet^A(P_\bullet^\lambda,Q_\bullet^\mu)\to
\mathbf{Tor}_\bullet^A(P_\bullet,Q_\bullet)$ déduites

\oldpage[III]{143}

des applications canoniques $P_\bullet^\lambda\to P_\bullet$,
$Q_\bullet^\mu\to Q_\bullet$, forment un système inductif d'homomorphismes,
d'où un homomorphisme canonique (6.3.6.1), et plus généralement un
homomorphisme canonique
\[
  \varinjlim_{\lambda,\mu}''E(P_\bullet^\lambda,Q_\bullet^\mu)
  \longrightarrow ''E(P_\bullet,Q_\bullet)
\]
dont (6.3.6.1) est l'homomorphisme des aboutissements. En outre, la suite
spectrale $''E(P_\bullet,Q_\bullet)$ est régulière (6.3.2), et il en est de
même de la suite spectrale
$\varinjlim_{\lambda,\mu}''E(P_\bullet^\lambda,Q_\bullet^\mu)$, comme il
résulte des définitions (0, 11.1.7) et de la démonstration de (0, 11.3.3);
pour démontrer que (6.3.6.1) est bijectif, il suffit donc (0, 11.1.5) de
prouver que l'homomorphisme
\[
  \varinjlim_{\lambda,\mu}''E(P_\bullet^\lambda,Q_\bullet^\mu)
  \longrightarrow ''E(P_\bullet,Q_\bullet)
  \tag{6.3.6.2}\label{III.6.3.6.2.fr}
\]
est bijectif pour les termes $E^2$. Comme le foncteur $H_\bullet$ commute à
la limite inductive des complexes de modules, on est finalement ramené à
prouver que pour deux systèmes inductifs filtrants $(F^\lambda)$,
$(G^\mu)$ de $A$-modules, l'homomorphisme canonique
\[
  \varinjlim_{\lambda,\mu}\operatorname{Tor}_\bullet^A(F^\lambda,G^\mu)
  \longrightarrow
  \operatorname{Tor}_\bullet^A
  \left(\varinjlim_\lambda F^\lambda,
        \varinjlim_\mu G^\mu\right)
\]
est bijectif. Pour cela, considérons pour chaque $F^\lambda$ la résolution
libre canonique
\[
  L_\bullet^\lambda:\quad
  \cdots\longrightarrow L_{i+1}^\lambda
  \longrightarrow L_i^\lambda\longrightarrow\cdots
  \longrightarrow L_1^\lambda\longrightarrow L_0^\lambda\longrightarrow0
\]
où $L_0^\lambda$ est le $A$-module des combinaisons linéaires formelles
d'éléments de $F^\lambda$ et $L_{i+1}^\lambda$ le $A$-module des combinaisons
linéaires formelles d'éléments de
$\operatorname{Ker}(L_i^\lambda\to L_{i-1}^\lambda)$; on vérifie
immédiatement que les $L_\bullet^\lambda$ forment un système inductif de
complexes, et si l'on pose
$F=\varinjlim_\lambda F^\lambda$, $L_i=\varinjlim_\lambda L_i^\lambda$, les
$L_i$ forment une résolution $L_\bullet$ de $F$, le foncteur
$\varinjlim$ étant exact; en outre, les $L_i$, limites inductives de
$A$-modules libres, sont plats (0\textsubscript{I}, 6.1.2). On considère de
même pour chaque $\mu$ la résolution libre canonique $M_\bullet^\mu$ de
$G^\mu$, et $M_\bullet=\varinjlim_\mu M_\bullet^\mu$ est une résolution plate
de $G=\varinjlim_\mu G^\mu$. On a alors
$\operatorname{Tor}_\bullet^A(\varinjlim_\lambda F^\lambda,
\varinjlim_\mu G^\mu)=H_\bullet(L_\bullet\otimes_A M_\bullet)$ en vertu de
(6.3.5) et (6.3.4); mais
\[
  H_\bullet(L_\bullet\otimes_A M_\bullet)
  =\varinjlim_{\lambda,\mu}
   H_\bullet(L_\bullet^\lambda\otimes_A M_\bullet^\mu)
\]
puisque $H_\bullet$ commute aux limites inductives de complexes de modules;
comme $H_\bullet(L_\bullet^\lambda\otimes_A M_\bullet^\mu)
=\operatorname{Tor}_\bullet^A(F^\lambda,G^\mu)$, cela termine la
démonstration.

Lorsqu'on suppose qu'il existe $i_0$ tel que
$P_i^\lambda=Q_i^\mu=0$ pour $i<i_0$ quels que soient $\lambda$ et $\mu$, on
démontre de la même manière que l'homomorphisme canonique
\[
  \varinjlim_{\lambda,\mu}'E(P_\bullet^\lambda,Q_\bullet^\mu)
  \longrightarrow'E(P_\bullet,Q_\bullet)
  \tag{6.3.6.3}\label{III.6.3.6.3.fr}
\]
est bijectif.
\end{proof}

\begin{proposition}[6.3.7]
Supposons $P_\bullet$ et $Q_\bullet$ limités inférieurement. Si le complexe
$P_\bullet$ est formé de $A$-modules plats, on a un $A$-isomorphisme
canonique de $\partial$-foncteurs en $Q_\bullet$
\[
  \mathbf{Tor}_\bullet^A(P_\bullet,Q_\bullet)
  \xrightarrow{\sim}H_\bullet(P_\bullet\otimes_A Q_\bullet).
  \tag{6.3.7.1}\label{III.6.3.7.1.fr}
\]
\end{proposition}

\begin{proof}
En effet, la suite spectrale (6.3.2.1) est birégulière et dégénérée, et
l'existence de l'isomorphisme (6.3.7.1) résulte de (0, 11.1.6). En outre, en
calculant l'hypertor à partir d'une résolution projective de Cartan-Eilenberg
de $P_\bullet$ (6.3.4), on voit aussitôt que l'isomorphisme ainsi défini est
un isomorphisme de $\partial$-foncteurs en $Q_\bullet$.
\end{proof}
\begin{env}[6.3.8]
\oldpage[III]{144}
Soit $\rho:A\to A'$ un homomorphisme d'anneaux. Nous nous proposons de
définir un $A$-homomorphisme factoriel de degré 0 canoniquement associé à
$\rho$ :
\[
  \rho_{P_\bullet,Q_\bullet}:
  \mathbf{Tor}_\bullet^A(P_\bullet,Q_\bullet)
  \longrightarrow
  \mathbf{Tor}_\bullet^{A'}
  (P_\bullet\otimes_A A',Q_\bullet\otimes_A A').
  \tag{6.3.8.1}\label{III.6.3.8.1.fr}
\]

Pour cela, considérons une résolution de Cartan-Eilenberg projective
$L_{\bullet\bullet}$ de $P_\bullet$; considérons d'autre part une résolution
de Cartan-Eilenberg projective $L'_{\bullet\bullet}$ de
$P_\bullet\otimes_A A'$. Nous allons voir qu'on peut définir un
$A'$-homomorphisme de complexes
$L_{\bullet\bullet}\otimes_A A'\to L'_{\bullet\bullet}$, déterminé à
homotopie près. En effet, la construction de $L_{\bullet\bullet}$ est
entièrement déterminée lorsqu'on se donne (arbitrairement) pour chaque $i$,
une résolution projective $(X_{ij}^{\mathrm B})_{j\geqslant0}$ de
$\mathrm B_i(P_\bullet)$ et une résolution projective
$(X_{ij}^{\mathrm H})_{j\geqslant0}$ de $\mathrm H_i(P_\bullet)$, qui sont
respectivement égales à $\mathrm B_i^{\mathrm I}(L_{\bullet\bullet})$ et
$\mathrm H_i^{\mathrm I}(L_{\bullet\bullet})$; on en déduit successivement
\[
  \mathrm Z_i^{\mathrm I}(L_{\bullet\bullet})
  =\mathrm H_i^{\mathrm I}(L_{\bullet\bullet})
   \oplus\mathrm B_i^{\mathrm I}(L_{\bullet\bullet}),
  \qquad
  L_{i,\bullet}=\mathrm Z_i^{\mathrm I}(L_{\bullet\bullet})
   \oplus\mathrm B_{i-1}^{\mathrm I}(L_{\bullet\bullet}).
\]
Cela étant, $X_{i,\bullet}^{\mathrm B}\otimes_A A'$ n'est plus en général
une résolution de $\mathrm B_i(P_\bullet)\otimes_A A'$, mais est encore un
complexe formé de $A'$-modules projectifs, et il y a donc un
$A'$-homomorphisme
\[
  X_{i,\bullet}^{\mathrm B}\otimes_A A'
  \longrightarrow \mathrm B_i^{\mathrm I}(L'_{\bullet\bullet})
\]
compatible avec les augmentations, et déterminé à homotopie près
(M, V, 1.1). On a de même un $A'$-homomorphisme
\[
  X_{i,\bullet}^{\mathrm H}\otimes_A A'
  \longrightarrow \mathrm H_i^{\mathrm I}(L'_{\bullet\bullet})
\]
déterminé à homotopie près, d'où l'on déduit, par la construction rappelée
plus haut, un $A'$-homomorphisme
$L_{i,\bullet}\otimes_A A'\to L'_{i,\bullet}$ pour tout $i$; ces
homomorphismes (pour $i\in\mathbf Z$) sont compatibles avec les opérateurs de
dérivation $L_{i,\bullet}\to L_{i-1,\bullet}$ et les analogues pour
$L'_{\bullet\bullet}$, en vertu de la même construction, et ils constituent
donc le $A'$-homomorphisme
$L_{\bullet\bullet}\otimes_A A'\to L'_{\bullet\bullet}$ cherché.

Pour définir (6.3.8.1), il suffit alors de considérer de même une résolution
de Cartan-Eilenberg projective $M_{\bullet\bullet}$ (resp.
$M'_{\bullet\bullet}$) de $Q_\bullet$ (resp.
$Q_\bullet\otimes_A A'$), et un $A'$-homomorphisme
$M_{\bullet\bullet}\otimes_A A'\to M'_{\bullet\bullet}$. On déduit de ces
homomorphismes un $A'$-homomorphisme
\[
  (L_{\bullet\bullet}\otimes_A A')\otimes_{A'}
  (M_{\bullet\bullet}\otimes_A A')
  \longrightarrow
  L'_{\bullet\bullet}\otimes_{A'}M'_{\bullet\bullet},
\]
puis par composition un $A$-homomorphisme de bicomplexes
\[
  L_{\bullet\bullet}\otimes_A M_{\bullet\bullet}
  \longrightarrow
  L'_{\bullet\bullet}\otimes_{A'}M'_{\bullet\bullet},
\]
et en passant à l'homologie on obtient (6.3.8.1), qui est bien défini
puisqu'il provient d'un morphisme de complexes défini à homotopie près.

Si $\rho':A'\to A''$ est un second homomorphisme d'anneaux, et
$\rho'':A\to A''$ l'homomorphisme composé $\rho'\circ\rho$, il est clair que
\[
  \rho''_{P_\bullet,Q_\bullet}
  =\rho'_{P'_\bullet,Q'_\bullet}\circ\rho_{P_\bullet,Q_\bullet},
  \qquad
  P'_\bullet=P_\bullet\otimes_A A',
  \quad Q'_\bullet=Q_\bullet\otimes_A A'.
\]
Notons encore que le morphisme de bicomplexes
$L_{\bullet\bullet}\otimes_A M_{\bullet\bullet}
\to L'_{\bullet\bullet}\otimes_{A'}M'_{\bullet\bullet}$ considéré
ci-dessus définit des morphismes factoriels (en $P_\bullet$ et $Q_\bullet$)
de suites spectrales
\[
  'E_{pq}^r(P_\bullet,Q_\bullet)
  \longrightarrow
  'E_{pq}^r(P_\bullet\otimes_A A',Q_\bullet\otimes_A A')
\]
et
\[
  ''E_{pq}^r(P_\bullet,Q_\bullet)
  \longrightarrow
  ''E_{pq}^r(P_\bullet\otimes_A A',Q_\bullet\otimes_A A'),
\]
indépendants des résolutions de Cartan-Eilenberg considérées, et ayant aussi
la propriété de transitivité précédente.
\end{env}

\begin{proposition}[6.3.9]
Soit $\rho:A\to A'$ un homomorphisme d'anneaux tel que $A'$ soit un
$A$-module plat. On a alors des isomorphismes canoniques factoriels
\[
  \mathbf{Tor}_\bullet^{A'}
  (P_\bullet\otimes_A A',Q_\bullet\otimes_A A')
  \xrightarrow{\sim}
  \mathbf{Tor}_\bullet^A(P_\bullet,Q_\bullet)\otimes_A A'.
  \tag{6.3.9.1}\label{III.6.3.9.1.fr}
\]
\[
  \left\{
  \begin{aligned}
  'E(P_\bullet\otimes_A A',Q_\bullet\otimes_A A')
    &\xrightarrow{\sim}'E(P_\bullet,Q_\bullet)\otimes_A A',\\
  ''E(P_\bullet\otimes_A A',Q_\bullet\otimes_A A')
    &\xrightarrow{\sim}''E(P_\bullet,Q_\bullet)\otimes_A A'.
  \end{aligned}
  \right.
  \tag{6.3.9.2}\label{III.6.3.9.2.fr}
\]
\end{proposition}

\begin{proof}
\oldpage[III]{145}
En effet, vu l'exactitude du foncteur $M\otimes_A A'$ en $M$,
$L_{\bullet\bullet}\otimes_A A'$ et $M_{\bullet\bullet}\otimes_A A'$ sont
alors des résolutions projectives de Cartan-Eilenberg de
$P_\bullet\otimes_A A'$ et $Q_\bullet\otimes_A A'$ respectivement, d'où la
conclusion.
\end{proof}

\begin{env}[6.3.10]
Soit $\rho:A\to A'$ un homomorphisme d'anneaux; pour tout complexe
$P'_\bullet$ de $A'$-modules, $P'_{\bullet[\rho]}$ est un complexe de
$A$-modules; en outre, l'application identique
$P'_{\bullet[\rho]}\to P'_\bullet$ peut être considérée comme composée des
applications canoniques
\[
  P'_{\bullet[\rho]}\longrightarrow
  P'_{\bullet[\rho]}\otimes_A A'
  \xrightarrow{\mu}P'_\bullet,
\]
où $\mu$ est le $A'$-homomorphisme $\mu(x\otimes a')=a'x$. Si $Q'_\bullet$
est un second complexe de $A'$-modules, on a donc des homomorphismes
canoniques factoriels de degré 0
\[
  \mathbf{Tor}_\bullet^A(P'_{\bullet[\rho]},Q'_{\bullet[\rho]})
  \longrightarrow
  \mathbf{Tor}_\bullet^{A'}
  (P'_{\bullet[\rho]}\otimes_A A',Q'_{\bullet[\rho]}\otimes_A A')
  \longrightarrow
  \mathbf{Tor}_\bullet^{A'}(P'_\bullet,Q'_\bullet),
  \tag{6.3.10.1}\label{III.6.3.10.1.fr}
\]
où la première flèche est le $A$-homomorphisme défini dans (6.3.8) et la
seconde se déduit des $A'$-homomorphismes
$P'_{\bullet[\rho]}\otimes_A A'\to P'_\bullet$ et
$Q'_{\bullet[\rho]}\otimes_A A'\to Q'_\bullet$ par fonctorialité. On a des
homomorphismes analogues pour les suites spectrales de (6.3.2), et des
propriétés évidentes de transitivité, que nous laissons au lecteur le soin
d'énoncer.
\end{env}

\begin{proposition}[6.3.11]
Soit $\rho:A\to A'$ un homomorphisme d'anneaux faisant de $A'$ un
$A$-module plat. Pour tout complexe $P'_\bullet$ de $A'$-modules et tout
complexe $Q_\bullet$ de $A$-modules limités inférieurement, on a un
isomorphisme canonique factoriel
\[
  \mathbf{Tor}_\bullet^A(P'_{\bullet[\rho]},Q_\bullet)
  \xrightarrow{\sim}
  \mathbf{Tor}_\bullet^{A'}(P'_\bullet,Q_\bullet\otimes_A A').
  \tag{6.3.11.1}\label{III.6.3.11.1.fr}
\]
\end{proposition}

\begin{proof}
En effet, si $M_{\bullet\bullet}$ est une résolution projective de
Cartan-Eilenberg de $Q_\bullet$, $M_{\bullet\bullet}\otimes_A A'$ est une
résolution projective de Cartan-Eilenberg de $Q_\bullet\otimes_A A'$, et
l'on a, à un isomorphisme canonique près,
\[
  P'_{\bullet[\rho]}\otimes_A M_{\bullet\bullet}
  =P'_\bullet\otimes_{A'}(M_{\bullet\bullet}\otimes_A A');
\]
la conclusion résulte de (6.3.4).
\end{proof}

\begin{namedtheorem}[6.3.12]{Remarque}
Soit $(A^\lambda)$ un système inductif filtrant d'anneaux, et soient
$(P_\bullet^\lambda)$, $(Q_\bullet^\lambda)$ deux systèmes inductifs de
complexes de $(A^\lambda)$-modules; on a alors un isomorphisme canonique
généralisant (6.3.6.1)
\[
  \varinjlim_\lambda
  \mathbf{Tor}_\bullet^{A^\lambda}
  (P_\bullet^\lambda,Q_\bullet^\lambda)
  \xrightarrow{\sim}
  \mathbf{Tor}_\bullet^A(P_\bullet,Q_\bullet),
  \tag{6.3.12.1}\label{III.6.3.12.1.fr}
\]
où $A=\varinjlim A^\lambda$, $P_\bullet=\varinjlim P_\bullet^\lambda$,
$Q_\bullet=\varinjlim Q_\bullet^\lambda$. Une fois définis les
homomorphismes
\[
  \mathbf{Tor}_\bullet^{A^\lambda}
  (P_\bullet^\lambda,Q_\bullet^\lambda)
  \longrightarrow
  \mathbf{Tor}_\bullet^{A^\mu}
  (P_\bullet^\mu,Q_\bullet^\mu)
\]
pour $\lambda\leqslant\mu$, à l'aide de (6.3.10), la démonstration est celle
de (6.3.6).
\end{namedtheorem}

\begin{proposition}[6.3.13]
Soient $S$ une partie multiplicative de $A$, $P_\bullet$ et $Q_\bullet$ deux
complexes de $A$-modules, dans lesquels les homothéties définies par les
éléments de $S$ soient bijectives, de sorte que, si $A'=S^{-1}A$,
$P_\bullet$ et $Q_\bullet$ sont formés de $A'$-modules. Alors on a un
isomorphisme canonique
\[
  \mathbf{Tor}_\bullet^A(P_\bullet,Q_\bullet)
  \xrightarrow{\sim}
  \mathbf{Tor}_\bullet^{A'}(P_\bullet,Q_\bullet).
\]
\end{proposition}

\begin{proof}
En effet, l'hypothèse entraîne que les homomorphismes canoniques
$P_\bullet\to P_\bullet\otimes_A A'$,
$Q_\bullet\to Q_\bullet\otimes_A A'$ sont bijectifs. D'autre part, la
fonctorialité de l'hypertor montre que tout $s\in S$ définit une homothétie
bijective dans $\mathbf{Tor}_\bullet^A(P_\bullet,Q_\bullet)$, et par suite
\[
  \mathbf{Tor}_\bullet^A(P_\bullet,Q_\bullet)
  \longrightarrow
  \mathbf{Tor}_\bullet^A(P_\bullet,Q_\bullet)\otimes_A A'
\]

\oldpage[III]{146}

est aussi un homomorphisme bijectif. Comme $A'$ est un $A$-module plat, la
conclusion résulte de (6.3.9), et on a de même des isomorphismes canoniques
pour les suites spectrales.
\end{proof}

\subsection{Foncteurs hypertor locaux de complexes de Modules quasi-cohérents : cas des schémas affines}

\begin{env}[6.4.1]
Soient $S$ un schéma affine d'anneau $A$, $X$, $Y$ deux $S$-schémas affines
d'anneaux $B$, $C$ respectivement, de sorte que $B$ et $C$ sont des algèbres
sur $A$. Tout complexe $\mathscr P_\bullet$ (resp. $\mathscr Q_\bullet$) de
$\mathscr O_X$-Modules (resp. $\mathscr O_Y$-Modules) quasi-cohérents est de
la forme $\widetilde{P_\bullet}$ (resp. $\widetilde{Q_\bullet}$), où
$P_\bullet$ (resp. $Q_\bullet$) est un complexe de $B$-modules (resp.
$C$-modules) (I, 1.3.7 et 1.3.8). On peut évidemment considérer $P_\bullet$
et $Q_\bullet$ comme des complexes de $A$-modules et former les
$\mathbf{Tor}_n^A(P_\bullet,Q_\bullet)$; en outre, en vertu du caractère
bifonctoriel de $\mathbf{Tor}_n^A(P_\bullet,Q_\bullet)$, les $A$-algèbres
$B$ et $C$ opèrent dans ce $A$-module, et ces opérations en font un
$(B,C)$-bimodule, ou, ce qui revient au même, un module sur
$B\otimes_A C=A(X\times_S Y)$. On a donc défini de la sorte un
$\mathscr O_{X\times_S Y}$-Module quasi-cohérent
\[
  \mathscr{T}\!or_n^{\mathscr O_S}(\mathscr P_\bullet,\mathscr Q_\bullet)
  =\bigl(\mathbf{Tor}_n^A(P_\bullet,Q_\bullet)\bigr)^{\sim}
  \tag{6.4.1.1}\label{III.6.4.1.1.fr}
\]
que l'on appelle l'\emph{hypertor local d'indice $n$} des complexes
$\mathscr P_\bullet$ et $\mathscr Q_\bullet$ et que l'on note aussi
$\mathscr{T}\!or_n^S(\mathscr P_\bullet,\mathscr Q_\bullet)$.
\end{env}

\begin{namedtheorem}[6.4.2]{Lemme}
Avec les notations de (6.4.1), supposons que l'anneau $A$ soit de la forme
$R^{-1}A'$, où $A'$ est un anneau et $R$ une partie multiplicative de $A'$.
Soit $S'=\operatorname{Spec}(A')$, de sorte que $X$ et $Y$ peuvent être
considérés comme des $S'$-préschémas et que l'on a
$X\times_{S'}Y=X\times_S Y$ (I, 1.6.2 et 3.2.4). On a alors
\[
  \mathscr{T}\!or_\bullet^{S'}(\mathscr P_\bullet,\mathscr Q_\bullet)
  =\mathscr{T}\!or_\bullet^S(\mathscr P_\bullet,\mathscr Q_\bullet).
\]
\end{namedtheorem}

\begin{proof}
Cela résulte de la formule (6.4.1.1) et de (6.3.13).
\end{proof}

\begin{env}[6.4.3]
Avec les notations et hypothèses de (6.4.1), soient
$\mathscr F=\widetilde F$ un $\mathscr O_X$-Module quasi-cohérent,
$\mathscr G=\widetilde G$ un $\mathscr O_Y$-Module quasi-cohérent;
considérant $\mathscr F$ et $\mathscr G$ comme des complexes de Modules, on
notera $\mathscr{T}\!or_n^{\mathscr O_S}(\mathscr F,\mathscr G)$ ou
$\mathscr{T}\!or_n^S(\mathscr F,\mathscr G)$ leur hypertor d'indice $n$; il
résulte de (6.3.5 (ii)) que l'on a
\[
  \mathscr{T}\!or_n^{\mathscr O_S}(\mathscr F,\mathscr G)
  =\bigl(\operatorname{Tor}_n^A(F,G)\bigr)^{\sim}.
  \tag{6.4.3.1}\label{III.6.4.3.1.fr}
\]

Revenons alors au cas général de deux complexes de Modules quasi-cohérents
$\mathscr P_\bullet$, $\mathscr Q_\bullet$. Les formules (6.4.1.1) et
(6.4.3.1) montrent, compte tenu de la prop. (6.3.2), que
$\mathscr{T}\!or_\bullet^S(\mathscr P_\bullet,\mathscr Q_\bullet)$ est
l'aboutissement de deux suites spectrales
$'\mathscr E(\mathscr P_\bullet,\mathscr Q_\bullet)$,
$''\mathscr E(\mathscr P_\bullet,\mathscr Q_\bullet)$, dont les termes
$E_2$ sont donnés par
\[
  '\mathscr E_{pq}^2
  =\mathscr H_p\bigl(\mathscr{T}\!or_q^S
    (\mathscr P_\bullet,\mathscr Q_\bullet)\bigr)
  \tag{6.4.3.2}\label{III.6.4.3.2.fr}
\]
\[
  ''\mathscr E_{pq}^2
  =\bigoplus_{q'+q''=q}
  \mathscr{T}\!or_p^S
  \bigl(\mathscr H_{q'}(\mathscr P_\bullet),
        \mathscr H_{q''}(\mathscr Q_\bullet)\bigr)
  \tag{6.4.3.3}\label{III.6.4.3.3.fr}
\]
où $\mathscr{T}\!or_q^S(\mathscr P_\bullet,\mathscr Q_\bullet)$ est le
bicomplexe de $\mathscr O_{X\times_S Y}$-Modules quasi-cohérents
$\mathscr{T}\!or_q^S(\mathscr P_i,\mathscr Q_j)$.
\end{env}

\begin{env}[6.4.4]
Considérons maintenant deux autres schémas affines
$X^{(1)}=\operatorname{Spec}(B^{(1)})$,
$Y^{(1)}=\operatorname{Spec}(C^{(1)})$, où $B^{(1)}$ et $C^{(1)}$ sont des
$A$-algèbres, et supposons donnés deux

\oldpage[III]{147}
S-morphismes $u:X^{(1)}\to X$, $v:Y^{(1)}\to Y$, correspondant à des
$A$-homomorphismes $\varphi:B\to B^{(1)}$, $\psi:C\to C^{(1)}$.
Considérons les complexes
$u^*(\mathscr P_\bullet)=(P_\bullet\otimes_B B^{(1)})^{\sim}$ de
$\mathscr O_{X^{(1)}}$-Modules,
$v^*(\mathscr Q_\bullet)=(Q_\bullet\otimes_C C^{(1)})^{\sim}$ de
$\mathscr O_{Y^{(1)}}$-Modules. Les $A$-homomorphismes canoniques
\[
  P_\bullet\longrightarrow P_\bullet\otimes_B B^{(1)},
  \qquad
  Q_\bullet\longrightarrow Q_\bullet\otimes_C C^{(1)}
\]
donnent par fonctorialité un $A$-homomorphisme
\[
  \mathbf{Tor}_\bullet^A(P_\bullet,Q_\bullet)
  \longrightarrow
  \mathbf{Tor}_\bullet^A
  (P_\bullet\otimes_B B^{(1)},Q_\bullet\otimes_C C^{(1)});
\]
en outre, toujours par fonctorialité, cet homomorphisme est en fait un
homomorphisme de $(B\otimes_A C)$-modules. D'où l'on conclut que l'on a défini
ainsi un $(u\times_S v)$-morphisme
\[
  \theta:\mathscr{T}\!or_\bullet^S(\mathscr P_\bullet,\mathscr Q_\bullet)
  \longrightarrow
  \mathscr{T}\!or_\bullet^S
  (u^*(\mathscr P_\bullet),v^*(\mathscr Q_\bullet))
  \tag{6.4.4.1}\label{III.6.4.4.1.fr}
\]
et par suite, un homomorphisme de
$\mathscr O_{X^{(1)}\times_S Y^{(1)}}$-Modules
\[
  \theta^\sharp:(u\times_S v)^*
  \bigl(\mathscr{T}\!or_\bullet^S(\mathscr P_\bullet,\mathscr Q_\bullet)\bigr)
  \longrightarrow
  \mathscr{T}\!or_\bullet^S
  (u^*(\mathscr P_\bullet),v^*(\mathscr Q_\bullet))
  \tag{6.4.4.2}\label{III.6.4.4.2.fr}
\]
qui est évidemment un morphisme de bi-$\partial$-foncteurs dans les catégories
de Modules quasi-cohérents limités inférieurement.

L'homomorphisme (6.4.4.2) n'est pas nécessairement bijectif; toutefois :
\end{env}

\begin{namedtheorem}[6.4.5]{Lemme}
Avec les notations de (6.4.4), supposons que $u$ et $v$ soient des immersions
ouvertes; alors l'homomorphisme (6.4.4.2) est bijectif.
\end{namedtheorem}

\begin{proof}
Identifions $X^{(1)}$ (resp. $Y^{(1)}$) à un ouvert de $X$ (resp. $Y$);
$X^{(1)}$ (resp. $Y^{(1)}$) est alors réunion d'ouverts de la forme
$\mathrm D(f)$ (resp. $\mathrm D(g)$), où $f\in B$ (resp. $g\in C$), et les
préschémas induits $\mathrm D(\varphi(f))$ et $\mathrm D(f)$ (resp.
$\mathrm D(\psi(g))$ et $\mathrm D(g)$) sont isomorphes. Il suffira de prouver
le lemme lorsque $X^{(1)}$ (resp. $Y^{(1)}$) est de la forme $\mathrm D(f)$
(resp. $\mathrm D(g)$); en effet, si ce point est établi, et si l'on revient au
cas général, il suffira de prouver que la restriction de $\theta^\sharp$ à
chaque ouvert $\mathrm D(f)\times_S\mathrm D(g)$ est un isomorphisme; or, si
$u_1:\mathrm D(f)\to X^{(1)}$, $v_1:\mathrm D(g)\to Y^{(1)}$ sont les
injections canoniques, la restriction précédente n'est autre que
$(u_1\times_S v_1)^*(\theta^\sharp)$; mais il est immédiat, en vertu des
définitions (6.4.4) et de ($0_{\mathrm I}$, 4.4.8), qu'en la composant avec
l'homomorphisme canonique
\[
  (u_1\times_S v_1)^*
  \bigl(\mathscr{T}\!or_\bullet^S
    (u^*(\mathscr P_\bullet),v^*(\mathscr Q_\bullet))\bigr)
  \longrightarrow
  \mathscr{T}\!or_\bullet^S
  ((u')^*(\mathscr P_\bullet),(v')^*(\mathscr Q_\bullet))
  \tag{6.4.5.1}\label{III.6.4.5.1.fr}
\]
où $u'=u\circ u_1$ et $v'=v\circ v_1$, on obtient l'homomorphisme canonique
\[
  (u'\times_S v')^*
  \bigl(\mathscr{T}\!or_\bullet^S(\mathscr P_\bullet,\mathscr Q_\bullet)\bigr)
  \longrightarrow
  \mathscr{T}\!or_\bullet^S
  ((u')^*(\mathscr P_\bullet),(v')^*(\mathscr Q_\bullet))
  \tag{6.4.5.2}\label{III.6.4.5.2.fr}
\]
et si l'on sait que (6.4.5.1) et (6.4.5.2) sont des isomorphismes, il en
résultera qu'il en est de même de $(u_1\times_S v_1)^*(\theta^\sharp)$.

Supposons donc que $X^{(1)}=\mathrm D(f)$ et $Y^{(1)}=\mathrm D(g)$, de sorte
que $B^{(1)}=B_f$ et $C^{(1)}=C_g$; $u^*(\mathscr P_\bullet)$ (resp.
$v^*(\mathscr Q_\bullet)$) s'identifie alors à
$((P_\bullet)_f)^{\sim}$ (resp. $((Q_\bullet)_g)^{\sim}$); d'autre part,
$X^{(1)}\times_S Y^{(1)}$ s'identifie au sous-schéma ouvert
$\mathrm D(f\otimes g)$ de
$X\times_S Y=\operatorname{Spec}(B\otimes_A C)$ (II, 4.3.2.4); il s'agit de
prouver que l'homomorphisme
\[
  \bigl(\mathbf{Tor}_\bullet^A(P_\bullet,Q_\bullet)\bigr)_{f\otimes g}
  \longrightarrow
  \mathbf{Tor}_\bullet^A((P_\bullet)_f,(Q_\bullet)_g)
  \tag{6.4.5.3}\label{III.6.4.5.3.fr}
\]

\oldpage[III]{148}
déduit par fonctorialité des homomorphismes canoniques
$P_\bullet\to(P_\bullet)_f$, $Q_\bullet\to(Q_\bullet)_g$, est bijectif. Or
($0_{\mathrm I}$, 1.6.1), on peut écrire
$(P_\bullet)_f=\varinjlim P_\bullet^{(n)}$, où les $P_\bullet^{(n)}$ sont tous
des complexes de $B$-modules identiques à $P_\bullet$, l'application
$P_\bullet^{(m)}\to P_\bullet^{(n)}$ pour $m\leqslant n$ étant la
multiplication par $f^{n-m}$; on a un résultat analogue pour $Q_\bullet$ en
remplaçant $f$ par $g$; d'autre part, il est clair que l'homomorphisme
\[
  \mathbf{Tor}_\bullet^A(P_\bullet^{(m)},Q_\bullet^{(m)})
  \longrightarrow
  \mathbf{Tor}_\bullet^A(P_\bullet^{(n)},Q_\bullet^{(n)})
\]
correspondant aux homomorphismes
$P_\bullet^{(m)}\to P_\bullet^{(n)}$ et
$Q_\bullet^{(m)}\to Q_\bullet^{(n)}$ est par définition la multiplication
par $(f\otimes g)^{n-m}$. La conclusion résulte alors de ($0_{\mathrm I}$,
1.6.1) appliqué au premier membre de (6.4.5.3) et de (6.3.6).
\end{proof}

\begin{env}[6.4.6]
Avec les notations de (6.4.4), on définit de même des homomorphismes canoniques
de foncteurs spectraux
\[
  \left\{
  \begin{aligned}
  (u\times_S v)^*\bigl('\mathscr E(\mathscr P_\bullet,\mathscr Q_\bullet)\bigr)
    &\longrightarrow
      '\mathscr E(u^*(\mathscr P_\bullet),v^*(\mathscr Q_\bullet)),\\
  (u\times_S v)^*\bigl(''\mathscr E(\mathscr P_\bullet,\mathscr Q_\bullet)\bigr)
    &\longrightarrow
      ''\mathscr E(u^*(\mathscr P_\bullet),v^*(\mathscr Q_\bullet)).
  \end{aligned}
  \right.
  \tag{6.4.6.1}\label{III.6.4.6.1.fr}
\]
et le raisonnement de (6.4.5) montre que lorsque $u$ et $v$ sont des
immersions ouvertes, les homomorphismes (6.4.6.1) sont bijectifs : en effet,
compte tenu de (6.3.6.2) et (6.3.6.3), il prouve que c'est un isomorphisme pour
les termes $E^2$, et (6.4.5) montre que c'est un isomorphisme pour les
aboutissements; on conclut donc à l'aide de (0, 11.1.2) et (0, 11.2.4).
\end{env}

\subsection{Foncteurs hypertor locaux de complexes de Modules quasi-cohérents : cas général}

\begin{env}[6.5.1]
Considérons maintenant un préschéma $S$ quelconque et deux $S$-préschémas
quelconques $X$, $Y$; soit $\mathscr P_\bullet$ (resp. $\mathscr Q_\bullet$)
un complexe de $\mathscr O_X$-Modules (resp. $\mathscr O_Y$-Modules)
quasi-cohérents. Posons $Z=X\times_S Y$; nous allons définir des
$\mathscr O_Z$-Modules quasi-cohérents
$\mathscr{T}\!or_n^S(\mathscr P_\bullet,\mathscr Q_\bullet)$ dits
\emph{hypertor locaux} de $\mathscr P_\bullet$ et $\mathscr Q_\bullet$ qui se
réduiront à ceux déjà définis dans (6.4) lorsque $S$, $X$ et $Y$ sont affines.

Lorsque $\mathscr P_\bullet$ et $\mathscr Q_\bullet$ se réduisent
respectivement à leurs termes de degré 0, $\mathscr F$ et $\mathscr G$ (les
autres étant nuls), on écrira
$\mathscr{T}\!or_n^S(\mathscr F,\mathscr G)$ au lieu de
$\mathscr{T}\!or_n^S(\mathscr P_\bullet,\mathscr Q_\bullet)$.
\end{env}

\begin{env}[6.5.2]
Supposons d'abord $S$ affine, et soient $(X_\lambda)$, $(Y_\mu)$ des
recouvrements de $X$ et $Y$ respectivement par des ouverts affines; alors les
$Z_{\lambda\mu}=X_\lambda\times_S Y_\mu$ forment un recouvrement ouvert affine
de $Z$. Posons
$\mathscr P_{\lambda\bullet}=\mathscr P_\bullet|X_\lambda$,
$\mathscr Q_{\mu\bullet}=\mathscr Q_\bullet|Y_\mu$; nous avons donc pour tout
couple $(\lambda,\mu)$ un $\mathscr O_{Z_{\lambda\mu}}$-Module quasi-cohérent
\[
  \mathscr F_{\lambda\mu}
  =\mathscr{T}\!or_n^S
   (\mathscr P_{\lambda\bullet},\mathscr Q_{\mu\bullet}),
\]
et il faut montrer que les $\mathscr F_{\lambda\mu}$ vérifient la condition de
recollement ($0_{\mathrm I}$, 3.3.1). Pour cela, il suffit de vérifier que pour
tout ouvert affine $U\subset X_\lambda\cap X_{\lambda'}$ (resp.
$V\subset Y_\mu\cap Y_{\mu'}$), les restrictions de
$\mathscr F_{\lambda\mu}$ et de $\mathscr F_{\lambda'\mu'}$ à $U\times_S V$
sont canoniquement isomorphes; mais cela découle aussitôt de l'existence
d'isomorphismes canoniques de ces restrictions sur
$\mathscr{T}\!or_n^S(\mathscr P_\bullet|U,\mathscr Q_\bullet|V)$ (6.4.5). En
outre, il résulte aussitôt de cette définition et de (6.4.5) que le
$\mathscr O_Z$-Module ainsi défini ne dépend pas (à un isomorphisme près) des
recouvrements ouverts consi\oldpage[III]{149}dérés;
nous le noterons donc
$\mathscr{T}\!or_n^S(\mathscr P_\bullet,\mathscr Q_\bullet)$; il résulte enfin
de (6.4.5) que pour toute partie ouverte $U$ (resp. $V$) de $X$ (resp. $Y$), la
restriction de $\mathscr{T}\!or_n^S(\mathscr P_\bullet,\mathscr Q_\bullet)$ à
$U\times_S V$ est canoniquement isomorphe à
$\mathscr{T}\!or_n^S(\mathscr P_\bullet|U,\mathscr Q_\bullet|V)$.
\end{env}

\begin{env}[6.5.3]
Passons maintenant au cas général où $S$ est quelconque, et soit $(S_\alpha)$
un recouvrement de $S$ formé d'ouverts affines; désignons par $X_\alpha$
(resp. $Y_\alpha$) l'image réciproque de $S_\alpha$ dans $X$ (resp. $Y$); il
faut encore prouver que les faisceaux
\[
  \mathscr{T}\!or_n^{S_\alpha}
  (\mathscr P_\bullet|X_\alpha,\mathscr Q_\bullet|Y_\alpha)
  =\mathscr G_\alpha
\]
vérifient la condition de recollement. Il suffit de définir, pour tout ouvert
affine $T$ contenu dans $S_\alpha\cap S_\beta$, des isomorphismes canoniques des
restrictions de $\mathscr G_\alpha$ et $\mathscr G_\beta$ à $U\times_S V$ (en
désignant par $U$ et $V$ les images réciproques de $T$ dans $X$ et $Y$
respectivement) sur
$\mathscr{T}\!or_n^T(\mathscr P_\bullet|U,\mathscr Q_\bullet|V)$; on peut, en
outre, se borner au cas où $T$ s'écrit à la fois $\mathrm D(f_\alpha)$ et
$\mathrm D(f_\beta)$, $f_\alpha$ (resp. $f_\beta$) étant une section de
$\mathscr O_S$ au-dessus de $S_\alpha$ (resp. $S_\beta$); mais alors
$\mathscr{T}\!or_n^T(\mathscr P_\bullet|U,\mathscr Q_\bullet|V)$ est
canoniquement isomorphe à
$\mathscr{T}\!or_n^{S_\alpha}(\mathscr P_\bullet|U,\mathscr Q_\bullet|V)$
d'une part, à
$\mathscr{T}\!or_n^{S_\beta}(\mathscr P_\bullet|U,\mathscr Q_\bullet|V)$
d'autre part, en vertu de (6.4.2); comme on vient de définir des isomorphismes
canoniques de $\mathscr G_\alpha$ sur
$\mathscr{T}\!or_n^{S_\alpha}(\mathscr P_\bullet|U,\mathscr Q_\bullet|V)$ et
de $\mathscr G_\beta$ sur
$\mathscr{T}\!or_n^{S_\beta}(\mathscr P_\bullet|U,\mathscr Q_\bullet|V)$
(6.5.2), cela achève de définir le $\mathscr O_Z$-Module
$\mathscr{T}\!or_n^S(\mathscr P_\bullet,\mathscr Q_\bullet)$. En outre, pour
toute partie ouverte $U$ (resp. $V$) de $X$ (resp. $Y$),
$\mathscr{T}\!or_n^S(\mathscr P_\bullet|U,\mathscr Q_\bullet|V)$ est
canoniquement isomorphe à la restriction de
$\mathscr{T}\!or_n^S(\mathscr P_\bullet,\mathscr Q_\bullet)$ à $U\times_S V$.

Il est immédiat que l'on a ainsi défini (dans les catégories de complexes de
Modules quasi-cohérents limités inférieurement) un bi-$\partial$-foncteur
$\mathscr{T}\!or_\bullet^S(\mathscr P_\bullet,\mathscr Q_\bullet)$ à valeurs
dans la catégorie des $\mathscr O_Z$-Modules, car il est clair que la question
est locale sur $X$, $Y$ et $S$, en vertu de (6.4.5) et de la remarque que
(6.4.4.2) est un morphisme de bi-$\partial$-foncteurs. On notera que si
$\mathscr P_\bullet$ et $\mathscr Q_\bullet$ sont réduits respectivement à
leurs termes de degré 0, $\mathscr F$ et $\mathscr G$,
$\mathscr{T}\!or_0^S(\mathscr F,\mathscr G)$ n'est autre, en vertu de
(6.4.1.1), que le produit tensoriel externe $\mathscr F\otimes_S\mathscr G$
défini dans (I, 9.1.2); cela résulte en effet de (I, 9.1.3).
\end{env}

\begin{env}[6.5.4]
Il résulte de la construction précédente et des remarques faites dans (6.4.6)
que $\mathscr{T}\!or_\bullet^S(\mathscr P_\bullet,\mathscr Q_\bullet)$ est
l'aboutissement de deux foncteurs spectraux,
$'\mathscr E(\mathscr P_\bullet,\mathscr Q_\bullet)$,
$''\mathscr E(\mathscr P_\bullet,\mathscr Q_\bullet)$, de termes $E_2$ égaux à
\[
  '\mathscr E_{pq}^2
  =\mathscr H_p\bigl(\mathscr{T}\!or_q^S
    (\mathscr P_\bullet,\mathscr Q_\bullet)\bigr)
  \tag{6.5.4.1}\label{III.6.5.4.1.fr}
\]
\[
  ''\mathscr E_{pq}^2
  =\bigoplus_{q'+q''=q}
   \mathscr{T}\!or_p^S
   \bigl(\mathscr H_{q'}(\mathscr P_\bullet),
         \mathscr H_{q''}(\mathscr Q_\bullet)\bigr).
  \tag{6.5.4.2}\label{III.6.5.4.2.fr}
\]
La suite spectrale (6.5.4.2) est toujours régulière; les deux suites spectrales
sont birégulières si $\mathscr P_\bullet$ et $\mathscr Q_\bullet$ sont limités
inférieurement. Un autre cas où les deux suites précédentes sont birégulières
est le suivant :
\end{env}

\begin{env}[6.5.5]
Nous dirons que sur un espace topologique $T$ un faisceau d'anneaux
$\mathscr A$ est \emph{de dimension cohomologique $\leqslant n$} si, pour tout
$t\in T$ l'anneau $\mathscr A_t$ est de dimension cohomologique $\leqslant n$;
on dira alors aussi que l'\emph{espace annelé} $(T,\mathscr A)$ est
\emph{de dimension cohomologique $\leqslant n$}. On dira qu'un faisceau
d'anneaux (resp. un espace annelé) est de dimension cohomologique \emph{finie}
s'il existe un entier $n$ tel qu'il soit de dimension cohomologique
$\leqslant n$. On notera que si les $\mathscr A_t$ sont des anneaux locaux
(commutatifs) noethériens,

\oldpage[III]{150}
dire qu'ils sont de dimension cohomologique $\leqslant n$ signifie qu'ils sont
réguliers et de dimension (de Krull) $\leqslant n$ ($0_{\mathrm{IV}}$, 17.3.1).
Avec la terminologie de la théorie de la dimension que nous introduirons au
chap. IV, il revient au même de dire qu'un préschéma localement noethérien $T$
est de dimension cohomologique $\leqslant n$, ou de dire qu'il est régulier
($0_{\mathrm I}$, 4.1.4) et de dimension $\leqslant n$; cela signifie que pour
tout ouvert affine $U$ de $T$, l'anneau $\Gamma(U,\mathscr O_T)$ est de
dimension cohomologique $\leqslant n$ ($0_{\mathrm{IV}}$, 17.2.6). Cela étant,
cette dernière remarque, jointe à (6.3.2), prouve que si $S$ est localement
noethérien et de dimension cohomologique finie, les suites spectrales
$'\mathscr E(\mathscr P_\bullet,\mathscr Q_\bullet)$ et
$''\mathscr E(\mathscr P_\bullet,\mathscr Q_\bullet)$ sont birégulières.

Il est clair que
$\mathscr{T}\!or_\bullet^S(\mathscr P_\bullet,\mathscr Q_\bullet)$ se
transforme en
$\mathscr{T}\!or_\bullet^S(\mathscr Q_\bullet,\mathscr P_\bullet)$ (à un
isomorphisme près) par l'isomorphisme canonique de $X\times_S Y$ sur
$Y\times_S X$.
\end{env}

\begin{proposition}[6.5.6]
Soit $(\mathscr P_{\alpha\bullet})$ un système inductif filtrant de complexes
de $\mathscr O_X$-Modules quasi-cohérents; il existe alors un isomorphisme
canonique
\[
  \varinjlim_\alpha
  \mathscr{T}\!or_\bullet^S
  (\mathscr P_{\alpha\bullet},\mathscr Q_\bullet)
  \xrightarrow{\sim}
  \mathscr{T}\!or_\bullet^S
  \left(\varinjlim_\alpha\mathscr P_{\alpha\bullet},\mathscr Q_\bullet\right).
  \tag{6.5.6.1}\label{III.6.5.6.1.fr}
\]
\end{proposition}

\begin{proof}
La question étant locale sur $S$, $X$ et $Y$, on peut supposer $S$, $X$, $Y$
affines et la proposition se réduit alors à (6.3.6).
\end{proof}

\begin{namedtheorem}[6.5.7]{Remarques}
\textnormal{(i)} Considérons en particulier le cas où $S=X=Y$,
$\mathscr P_\bullet$ et $\mathscr Q_\bullet$ étant donc deux complexes de
$\mathscr O_S$-Modules quasi-cohérents; alors les
$\mathscr{T}\!or_n^S(\mathscr P_\bullet,\mathscr Q_\bullet)$ sont des
$\mathscr O_S$-Modules quasi-cohérents; en outre, pour tout point $z\in S$, il
résulte de (6.5.6) que l'on a un isomorphisme canonique
\[
  \bigl(\mathscr{T}\!or_n^S
    (\mathscr P_\bullet,\mathscr Q_\bullet)\bigr)_z
  \xrightarrow{\sim}
  \mathbf{Tor}_n^{\mathscr O_z}
  \bigl((\mathscr P_\bullet)_z,(\mathscr Q_\bullet)_z\bigr)
  \tag{6.5.7.1}\label{III.6.5.7.1.fr}
\]
car la question est locale et on est ramené au cas des modules, en vertu de
(6.4.1.1).

\textnormal{(ii)} On peut généraliser la définition des hypertor au cas de
deux complexes de $\mathscr O_X$-Modules $\mathscr P_\bullet$,
$\mathscr Q_\bullet$ sur un même espace annelé $(X,\mathscr O_X)$; pour tout
ouvert $U$ de $X$, posons en effet
$A(U)=\Gamma(U,\mathscr O_X)$,
$P_\bullet(U)=\Gamma(U,\mathscr P_\bullet)$,
$Q_\bullet(U)=\Gamma(U,\mathscr Q_\bullet)$; les $A(U)$-modules
\[
  \mathbf{Tor}_n^{A(U)}(P_\bullet(U),Q_\bullet(U))
\]
forment alors un préfaisceau sur $X$, et l'on désigne par
$\mathscr{T}\!or_n^X(\mathscr P_\bullet,\mathscr Q_\bullet)$ le
$\mathscr O_X$-Module associé à ce préfaisceau. Lorsque $X$ est un préschéma,
il résulte de (6.3.12) que ce $\mathscr O_X$-Module est canoniquement isomorphe
à l'hypertor défini ci-dessus. Nous ne développerons pas davantage cette
généralisation.
\end{namedtheorem}

\begin{proposition}[6.5.8]
Soient $X$, $Y$ deux $S$-préschémas, $\mathscr F$ (resp. $\mathscr G$) un
$\mathscr O_X$-Module (resp. un $\mathscr O_Y$-Module) quasi-cohérent. Si
$\mathscr F$ ou $\mathscr G$ est $S$-plat, on a
$\mathscr{T}\!or_n^S(\mathscr F,\mathscr G)=0$ pour $n\neq0$.
\end{proposition}

\begin{proof}
La question étant locale sur $X$ et $Y$, on peut supposer $X$, $Y$ et $S$
affines, d'anneaux respectifs $B$, $C$, $A$, et
$\mathscr F=\widetilde M$, $\mathscr G=\widetilde N$, $M$ (resp. $N$) étant
un $B$-module (resp. un $C$-module). Supposons par exemple que $\mathscr F$
soit $S$-plat, ce qui signifie que pour tout $s\in S$, $M_s$ est un
$A_s$-module plat ($0_{\mathrm I}$, 6.7.1); par suite $M$ est un $A$-module plat
($0_{\mathrm I}$, 6.3.3), et l'on sait que
$\operatorname{Tor}_n^A(M,N)=0$ pour $n>0$ et pour tout $C$-module $N$
($0_{\mathrm I}$, 6.1.1), d'où la conclusion par (6.4.1.1).
\end{proof}

\begin{namedtheorem}[6.5.9]{Corollaire}
Soient $X$, $Y$ deux $S$-préschémas, $\mathscr P_\bullet$ (resp.
$\mathscr Q_\bullet$) un complexe de

\oldpage[III]{151}
$\mathscr O_X$-Modules (resp. de $\mathscr O_Y$-Modules) quasi-cohérents limité
inférieurement. Supposons que tous les $\mathscr P_i$ soient $S$-plats. Alors
il existe un isomorphisme canonique de $\partial$-foncteurs en
$\mathscr Q_\bullet$
\[
  \mathscr{T}\!or_\bullet^S(\mathscr P_\bullet,\mathscr Q_\bullet)
  \xrightarrow{\sim}
  \mathscr H_\bullet(\mathscr P_\bullet\otimes_S\mathscr Q_\bullet).
  \tag{6.5.9.1}\label{III.6.5.9.1.fr}
\]
\end{namedtheorem}

\begin{proof}
Ce n'est autre que (6.3.7) lorsque $S$, $X$, $Y$ sont affines; on passe de là
au cas général par les raisonnements de (6.5.2) et (6.5.3).
\end{proof}

\begin{namedtheorem}[6.5.10]{Corollaire}
Supposons que $X$ soit plat sur $S$ ($0_{\mathrm I}$, 6.7.1), que
$\mathscr P_\bullet$ et $\mathscr Q_\bullet$ soient limités inférieurement et
que tous les $\mathscr P_i$ soient des $\mathscr O_X$-Modules localement
libres (non nécessairement de type fini). Alors l'homomorphisme (6.5.9.1) est
bijectif.
\end{namedtheorem}

\begin{proof}
En effet, l'hypothèse de (6.5.9) est remplie, la platitude étant une propriété
ponctuelle sur $X$ par définition et toute somme directe de modules plats étant
un module plat ($0_{\mathrm I}$, 6.1.2).
\end{proof}

\begin{proposition}[6.5.11]
Soient $X'$, $Y'$ deux $S$-préschémas, $f:X\to X'$, $g:Y\to Y'$ deux
$S$-morphismes affines. Soit $\mathscr P_\bullet$ (resp.
$\mathscr Q_\bullet$) un complexe de $\mathscr O_X$-Modules (resp. de
$\mathscr O_Y$-Modules) quasi-cohérents; on a alors un isomorphisme canonique
fonctoriel
\[
  (f\times_S g)_*
  \bigl(\mathscr{T}\!or_\bullet^S
    (\mathscr P_\bullet,\mathscr Q_\bullet)\bigr)
  \xrightarrow{\sim}
  \mathscr{T}\!or_\bullet^S
  (f_*(\mathscr P_\bullet),g_*(\mathscr Q_\bullet)).
  \tag{6.5.11.1}\label{III.6.5.11.1.fr}
\]
\end{proposition}

\begin{proof}
Comme $f$ et $g$ sont affines, $f_*(\mathscr P_\bullet)$ et
$g_*(\mathscr Q_\bullet)$ sont des complexes de Modules quasi-cohérents
(II, 1.2.6), et si l'on pose $Z'=X'\times_S Y'$, les deux membres de (6.5.11.1)
sont des $\mathscr O_{Z'}$-Modules quasi-cohérents (6.5.1); on se ramène
aisément au cas où $S$, $X'$ et $Y'$ sont affines; mais alors il en est de même
par hypothèse de $X$ et $Y$ et la vérification résulte aussitôt de (6.4.1.1) et
(I, 1.6.3).
\end{proof}

\begin{namedtheorem}[6.5.12]{Remarque}
Soient $X'$, $Y'$ deux $S$-préschémas et supposons que $X'$ soit $S$-plat;
soient $X$ un sous-préschéma fermé de $X'$, $i:X\to X'$ l'injection canonique,
$\mathscr P_\bullet$ (resp. $\mathscr Q_\bullet$) un complexe de
$\mathscr O_X$-Modules (resp. de $\mathscr O_{Y'}$-Modules) quasi-cohérents,
limité inférieurement. Soit enfin $\mathscr L'_{\bullet\bullet}$ une résolution
de $i_*(\mathscr P_\bullet)$ formé de $\mathscr O_{X'}$-Modules localement
libres, telle que tout point de $X'$ ait un voisinage ouvert affine $U$ pour
lequel $\mathscr L'_{j,\bullet}|U$ soit une résolution libre de
$i_*(\mathscr P_j)|U$ pour tout $j$. On a alors un isomorphisme canonique
\[
  (i\times_S 1)_*
  \bigl(\mathscr{T}\!or_\bullet^S
    (\mathscr P_\bullet,\mathscr Q_\bullet)\bigr)
  \xrightarrow{\sim}
  \mathscr H_\bullet
  (\mathscr L'_{\bullet\bullet}\otimes_S\mathscr Q_\bullet).
  \tag{6.5.12.1}\label{III.6.5.12.1.fr}
\]

Si $S$, $X'$, $Y'$ sont affines et si $\mathscr L'_{j,\bullet}$ est une
résolution libre de $i_*(\mathscr P_j)$ pour tout $j$, on est ramené, en vertu
de (6.5.11), au cas où $X'=X$ est $S$-plat, et il suffit d'appliquer (6.3.4).
Dans le cas général, on définit localement l'isomorphisme (6.5.12.1), et il
s'agit de vérifier que cette définition donne bien un isomorphisme global.
Pour cela, il faut se reporter à la définition du premier isomorphisme
(6.3.4.1) qui provient d'un isomorphisme de suites spectrales (0, 11.6.5 et
11.5.3), obtenu lui-même à partir d'un morphisme de bicomplexes
$L_{\bullet\bullet}\to L''_{\bullet\bullet}$, où $L_{\bullet\bullet}$ est la
résolution donnée de $P_\bullet$, $L''_{\bullet\bullet}$ une résolution
projective de $P_\bullet$ dans la catégorie des complexes de $A$-modules
limités

\oldpage[III]{152}
inférieurement (cf. (0, 11.5.2.2)); notre assertion résulte de ce que
l'isomorphisme (6.3.4.1) ne dépend pas de la résolution projective
$L''_{\bullet\bullet}$ choisie, en vertu de l'existence d'un homotopisme entre
deux telles résolutions (M, V, 1.2).
\end{namedtheorem}

\begin{proposition}[6.5.13]
Soient $X$, $Y$ deux $S$-préschémas, et supposons vérifiée l'une des conditions
suivantes :
\begin{enumerate}
  \item[(i)] $X$ et $Z=X\times_S Y$ sont localement noethériens et $X$ est plat
    sur $S$.
  \item[(ii)] $S$ et $X$ sont localement noethériens et $Y$ est de type fini
    sur $S$.
\end{enumerate}
Soit $\mathscr P_\bullet$ (resp. $\mathscr Q_\bullet$) un complexe de
$\mathscr O_X$-Modules (resp. de $\mathscr O_Y$-Modules) quasi-cohérents,
limité inférieurement. On suppose en outre que, pour tout $n$,
$\mathscr H_n(\mathscr P_\bullet)$ (resp.
$\mathscr H_n(\mathscr Q_\bullet)$) est un $\mathscr O_X$-Module (resp. un
$\mathscr O_Y$-Module) de type fini. Alors les
$\mathscr{T}\!or_n^S(\mathscr P_\bullet,\mathscr Q_\bullet)$ sont des
$\mathscr O_Z$-Modules cohérents.
\end{proposition}

\begin{proof}
Comme $\mathscr P_\bullet$ et $\mathscr Q_\bullet$ sont limités
inférieurement, la suite spectrale
$''\mathscr E(\mathscr P_\bullet,\mathscr Q_\bullet)$ est birégulière (6.5.4),
et en vertu de (0, 11.1.8), il suffit (puisque dans les deux cas (i), (ii), $Z$
est localement noethérien) de prouver que les termes
$''\mathscr E_{pq}^2$ sont cohérents. L'hypothèse sur les
$\mathscr H_n(\mathscr P_\bullet)$ et
$\mathscr H_n(\mathscr Q_\bullet)$ et l'expression (6.5.4.2) des
$''\mathscr E_{pq}^2$ montrent donc que la proposition est équivalente à son
cas particulier correspondant à $\mathscr P_\bullet$ et
$\mathscr Q_\bullet$ réduits à leurs termes de degré 0, autrement dit à son
\end{proof}

\begin{corollary}[6.5.14]
Supposons vérifiée l'une des conditions (i), (ii) de (6.5.13), et soit
$\mathscr F$ (resp. $\mathscr G$) un $\mathscr O_X$-Module (resp. un
$\mathscr O_Y$-Module) quasi-cohérent de type fini; alors les
$\mathscr{T}\!or_n^S(\mathscr F,\mathscr G)$ sont des $\mathscr O_Z$-Modules
cohérents.
\end{corollary}

\begin{proof}
La question étant locale sur $X$ et $Y$, on peut supposer $S$, $X$, $Y$
affines.

\textnormal{(i)} Sous les hypothèses de (i), $S$, $X$ et $Z$ sont noethériens.
Il existe donc une résolution localement libre $\mathscr L_\bullet$ de
$\mathscr F$ formée de $\mathscr O_X$-Modules de type fini (I, 1.3.7); comme
$X$ est plat sur $S$, il résulte de (6.3.4) que l'on a
\[
  \mathscr{T}\!or_n^S(\mathscr F,\mathscr G)
  =\mathscr H_n(\mathscr L_\bullet\otimes_S\mathscr G);
\]
or, les $\mathscr L_i\otimes_S\mathscr G$ sont des $\mathscr O_Z$-Modules
quasi-cohérents de type fini (I, 9.1.1), donc cohérents. On en conclut que
$\mathscr H_n(\mathscr L_\bullet\otimes_S\mathscr G)$ est cohérent
($0_{\mathrm I}$, 5.3.4).

\textnormal{(ii)} Supposons maintenant vérifiées les conditions (ii). Comme
l'anneau $A(Y)$ est quotient d'une $A(S)$-algèbre de polynômes à un nombre fini
d'indéterminées (I, 6.3.3), $Y$ est un sous-$S$-préschéma fermé d'un
$S$-préschéma affine $Y'$, plat et de type fini sur $S$; $Y'$ étant noethérien
(I, 6.3.7), il existe une résolution localement libre $\mathscr M_\bullet$ de
$j_*(\mathscr G)$ par des $\mathscr O_{Y'}$-Modules de type fini
($j:Y\to Y'$ étant l'injection canonique); en vertu de (6.5.12),
$\mathscr{T}\!or_n^S(\mathscr F,\mathscr G)$ est l'image réciproque par
$1\times j$ du $\mathscr O_{Z'}$-Module
$\mathscr H_n(\mathscr F\otimes_S\mathscr M_\bullet)$ (où
$Z'=X\times_S Y'$); on voit comme dans (i) que les
$\mathscr F\otimes_S\mathscr M_i$ sont des $\mathscr O_{Z'}$-Modules
cohérents, et l'on en tire encore la conclusion par ($0_{\mathrm I}$, 5.3.4).
\end{proof}

\begin{env}[6.5.15]
La théorie développée ci-dessus pour deux complexes $\mathscr P_\bullet$,
$\mathscr Q_\bullet$ de faisceaux quasi-cohérents sur deux $S$-préschémas
$X$, $Y$ se généralise sans peine au cas où l'on considère un nombre fini
quelconque de $S$-préschémas $X^{(i)}$ ($1\leqslant i\leqslant m$) et sur
chaque $X^{(i)}$ un complexe $\mathscr P_\bullet^{(i)}$ de
$\mathscr O_{X^{(i)}}$-Modules quasi-cohérents; si
$Z=X^{(1)}\times_S X^{(2)}\times_S\cdots\times_S X^{(m)}$, on définit ainsi un
$\mathscr O_Z$-Module quasi-cohérent
$\mathscr{T}\!or_q^S(\mathscr P_\bullet^{(1)},\mathscr P_\bullet^{(2)},
\ldots,\mathscr P_\bullet^{(m)})$. Nous laissons au lecteur le soin de
développer la théorie dans ce cas général et nous nous bornerons à écrire,
pour référence ultérieure, le terme $E_2$ de la seconde suite spectrale
(régulière) dont l'aboutissement est
$\mathscr{T}\!or_\bullet^S(\mathscr P_\bullet^{(1)},
\mathscr P_\bullet^{(2)},\ldots,\mathscr P_\bullet^{(m)})$ :

\oldpage[III]{153}

\[
  ''\mathscr E_{pq}^2
  =\bigoplus_{q_1+q_2+\cdots+q_m=q}
   \mathscr{T}\!or_p^S
   \bigl(\mathscr H_{q_1}(\mathscr P_\bullet^{(1)}),\ldots,
         \mathscr H_{q_m}(\mathscr P_\bullet^{(m)})\bigr).
  \tag{6.5.15.1}\label{III.6.5.15.1.fr}
\]
Nous étudierons dans (6.8) les suites spectrales d'associativité auxquelles
donnent lieu ces foncteurs hypertor d'un nombre quelconque de complexes.
\end{env}

\subsection{Foncteurs hypertor globaux de complexes de Modules quasi-cohérents
et suites spectrales de Künneth : cas de la base affine}
\label{subsection:III.6.6.fr}

\begin{env}[6.6.1]
Considérons un schéma affine $S=\operatorname{Spec}(A)$ et deux $S$-schémas
quasi-compacts $X^{(i)}$ ($i=1,2$); soit $\mathscr P_\bullet^{(i)}$ un complexe
de $\mathscr O_{X^{(i)}}$-Modules quasi-cohérents, limité inférieurement, dont
l'opérateur de dérivation est de degré $-1$ ($i=1,2$). Considérons d'autre
part un recouvrement fini
$\mathfrak U^{(i)}=(U_\alpha^{(i)})$ de $X^{(i)}$ par des ouverts affines;
soit
\[
  X=X^{(1)}\times_S X^{(2)},
\]
qui est un $S$-schéma quasi-compact (I, 5.5.1 et 6.6.4), et soit
$\mathfrak U=\mathfrak U^{(1)}\times_S\mathfrak U^{(2)}$ le recouvrement de
$X$ formé des ouverts affines
$U_\alpha^{(1)}\times_S U_\beta^{(2)}$. Pour tout couple d'entiers
$p\leqslant0$, $q\in\mathbf Z$, le groupe des $(-p)$-cochaînes alternées
\[
  C^{-p}(\mathfrak U^{(i)},\mathscr P_q^{(i)})
\]
du recouvrement $\mathfrak U^{(i)}$ à coefficients dans le faisceau
$\mathscr P_q^{(i)}$ (G, II, 5.1) est un $A$-module; pour $p>0$, on posera
$C^p(\mathfrak U^{(i)},\mathscr P_q^{(i)})=0$; on a ainsi défini un
bicomplexe
\[
  C^\bullet(\mathfrak U^{(i)},\mathscr P_\bullet^{(i)})
\]
de $A$-modules, dont les deux opérateurs de dérivation sont de degré $-1$. Il
résulte des définitions ($0_{\mathrm I}$, 12.1.2) que le $A$-module
d'homologie
\[
  \mathrm H_n\bigl(C^\bullet
  (\mathfrak U^{(i)},\mathscr P_\bullet^{(i)})\bigr)
\]
de ce bicomplexe (considéré à l'ordinaire comme un complexe simple pour le
degré total) n'est autre que le $A$-module d'hypercohomologie
\[
  \mathbf H^{-n}
  (\mathfrak U^{(i)},\mathscr P^{(i)\bullet}),
\]
où $\mathscr P^{(i)\bullet}$ est le complexe à opérateur de dérivation de
degré $+1$ obtenu en prenant $\mathscr P_{-q}^{(i)}$ pour composante de degré
$q$; par abus de notation, nous l'écrirons
$\mathbf H^{-n}(\mathfrak U^{(i)},\mathscr P_\bullet^{(i)})$. Il résulte alors
de (6.2.2) que ce $A$-module est canoniquement isomorphe au $A$-module
d'hypercohomologie
$\mathbf H^{-n}(X^{(i)},\mathscr P^{(i)\bullet})$, que nous écrirons de même
$\mathbf H^{-n}(X^{(i)},\mathscr P_\bullet^{(i)})$; il ne dépend donc pas du
recouvrement fini $\mathfrak U^{(i)}$ choisi.
\end{env}

\begin{env}[6.6.2]
Nous allons appliquer aux deux bicomplexes de $A$-modules
\[
  L_{\bullet\bullet}^{(i)}
  =C^\bullet(\mathfrak U^{(i)},\mathscr P_\bullet^{(i)})
  \qquad(i=1,2)
\]
et au bifoncteur covariant
$L_{\bullet\bullet}^{(1)}\otimes_A L_{\bullet\bullet}^{(2)}$ en ces deux
bicomplexes, la théorie générale de l'hyperhomologie des foncteurs par rapport
aux bicomplexes ($0_{\mathrm I}$, 11.7.4). Comme les cochaînes considérées sont
alternées et les recouvrements $\mathfrak U^{(i)}$ finis, on notera que les
modules $C^{-p}(\mathfrak U^{(i)},\mathscr P_q^{(i)})$ ne sont $\ne0$ que pour
un nombre fini (indépendant de $q$) de valeurs de $p$, et en particulier les
deux degrés de chacun des $L_{\bullet\bullet}^{(i)}$ sont limités
inférieurement. Nous désignerons par
\[
  \mathbf{Tor}_n^A(L_{\bullet\bullet}^{(1)},L_{\bullet\bullet}^{(2)})
\]
ou
\[
  \mathbf{Tor}_n^S
  (\mathfrak U^{(1)},\mathfrak U^{(2)};
   \mathscr P_\bullet^{(1)},\mathscr P_\bullet^{(2)})
\]
le $n$-ième module d'hyperhomologie de
$L_{\bullet\bullet}^{(1)}\otimes_A L_{\bullet\bullet}^{(2)}$, que nous
appellerons l'\emph{hypertor} d'indice $n$ de
$\mathscr P_\bullet^{(1)}$ et $\mathscr P_\bullet^{(2)}$, relatif aux
recouvrements $\mathfrak U^{(1)}$ et $\mathfrak U^{(2)}$. Lorsque
$\mathscr P_\bullet^{(1)}$ et $\mathscr P_\bullet^{(2)}$ sont réduits à leurs
termes de degré 0, $\mathscr F^{(1)}$ et $\mathscr F^{(2)}$, on écrit
\[
  \mathbf{Tor}_n^S
  (\mathfrak U^{(1)},\mathfrak U^{(2)};
   \mathscr F^{(1)},\mathscr F^{(2)})
\]
leur hypertor. On désignera par
\[
  \mathbf{Tor}_n^S
  (\mathfrak U^{(1)},\mathfrak U^{(2)};
   \mathscr P_\bullet^{(1)},\mathscr P_\bullet^{(2)}),
\]
suivant les conventions générales, le bicomplexe dont le composant d'indices
$(j,k)$ est
\[
  \mathbf{Tor}_n^S
  (\mathfrak U^{(1)},\mathfrak U^{(2)};
   \mathscr P_j^{(1)},\mathscr P_k^{(2)}).
\]

Comme $L_{\bullet\bullet}^{(i)}$ est un foncteur exact en
$\mathscr P_\bullet^{(i)}$, puisque les intersections des ensembles
\oldpage[III]{154}
de $\mathfrak U^{(i)}$ sont affines (I, 5.5.6 et 1.3.11),
$\mathbf{Tor}_\bullet^S
(\mathfrak U^{(1)},\mathfrak U^{(2)};
 \mathscr P_\bullet^{(1)},\mathscr P_\bullet^{(2)})$ est un
bi-$\partial$-foncteur covariant en $\mathscr P_\bullet^{(1)}$,
$\mathscr P_\bullet^{(2)}$, à valeurs dans la catégorie des $A$-modules
($0_{\mathrm I}$, 11.7.3). En outre, on sait ($0_{\mathrm I}$, 11.7.2) que ce
bifoncteur est l'aboutissement commun de six foncteurs spectraux biréguliers,
que nous désignerons par la notation
\[
  {}^{(t)}\mathscr E^S
  (\mathfrak U^{(1)},\mathfrak U^{(2)};
   \mathscr P_\bullet^{(1)},\mathscr P_\bullet^{(2)})
\]
ou
\[
  {}^{(t)}\mathscr E
  (\mathfrak U^{(1)},\mathfrak U^{(2)};
   \mathscr P_\bullet^{(1)},\mathscr P_\bullet^{(2)}),
\]
où $t$ doit être remplacé par une des lettres $a$, $b$, $a'$, $b'$, $c$,
$d$, et dont les termes $E_2$ sont les suivants :
\[
\begin{aligned}
  {}^{(a)}\mathscr E_{pq}^2
  &=\bigoplus_{q_1+q_2=q}
    \mathbf{Tor}_p^A
    \bigl(\mathrm H_{q_1}(L_{\bullet\bullet}^{(1)}),
          \mathrm H_{q_2}(L_{\bullet\bullet}^{(2)})\bigr),\\
  {}^{(b)}\mathscr E_{pq}^2
  &=\mathrm H_p\bigl(
    \mathbf{Tor}_q^{A,\mathrm{II}}
    (L_{\bullet\bullet}^{(1)},L_{\bullet\bullet}^{(2)})\bigr),\\
  {}^{(a')}\mathscr E_{pq}^2
  &=\bigoplus_{q_1+q_2=q}
    \mathbf{Tor}_p^A
    \bigl(\mathrm H_{q_1}^{\mathrm I}(L_{\bullet\bullet}^{(1)}),
          \mathrm H_{q_2}^{\mathrm I}(L_{\bullet\bullet}^{(2)})\bigr),\\
  {}^{(b')}\mathscr E_{pq}^2
  &=\mathrm H_p\bigl(
    \mathbf{Tor}_q^A
    (L_{\bullet\bullet}^{(1)},L_{\bullet\bullet}^{(2)})\bigr),\\
  {}^{(c)}\mathscr E_{pq}^2
  &=\bigoplus_{q_1+q_2=q}
    \mathbf{Tor}_p^A
    \bigl(\mathrm H_{q_1}^{\mathrm{II}}(L_{\bullet\bullet}^{(1)}),
          \mathrm H_{q_2}^{\mathrm{II}}(L_{\bullet\bullet}^{(2)})\bigr),\\
  {}^{(d)}\mathscr E_{pq}^2
  &=\mathrm H_p\bigl(
    \mathbf{Tor}_q^{A,\mathrm I}
    (L_{\bullet\bullet}^{(1)},L_{\bullet\bullet}^{(2)})\bigr),
\end{aligned}
\]
où les notations sont conformes à celles de la théorie générale de
l'hyperhomologie. Nous allons dans ce qui suit expliciter davantage ces
termes initiaux.
\end{env}

\begin{env}[6.6.3]
\emph{Suites spectrales $(a)$ et $(a')$.} Nous avons vu en (6.6.1) que le
module d'homologie $\mathrm H_n(L_{\bullet\bullet}^{(i)})$ du bicomplexe
$L_{\bullet\bullet}^{(i)}$ était égal à
$\mathbf H^{-n}(X^{(i)},\mathscr P_\bullet^{(i)})$; donc
\[
  {}^{(a)}\mathscr E_{pq}^2
  =\bigoplus_{q_1+q_2=q}
   \mathbf{Tor}_p^A
   \bigl(\mathbf H^{-q_1}(X^{(1)},\mathscr P_\bullet^{(1)}),
         \mathbf H^{-q_2}(X^{(2)},\mathscr P_\bullet^{(2)})\bigr).
\]
Par définition, le complexe
$\mathrm H_n^{\mathrm I}(L_{\bullet\bullet}^{(i)})$ a pour terme de degré $k$
le module d'homologie
$\mathrm H_n(C^\bullet(\mathfrak U^{(i)},\mathscr P_k^{(i)}))$, c'est-à-dire,
par définition, le module de cohomologie
$\mathbf H^{-n}(\mathfrak U^{(i)},\mathscr P_k^{(i)})$; on sait (1.4.1) que
ce module est canoniquement isomorphe à
$\mathbf H^{-n}(X^{(i)},\mathscr P_k^{(i)})$; donc
\[
  {}^{(a')}\mathscr E_{pq}^2
  =\bigoplus_{q_1+q_2=q}
   \mathbf{Tor}_p^A
   \bigl(\mathbf H^{-q_1}(X^{(1)},\mathscr P_\bullet^{(1)}),
         \mathbf H^{-q_2}(X^{(2)},\mathscr P_\bullet^{(2)})\bigr).
\]
\end{env}

\begin{env}[6.6.4]
\emph{Suites spectrales $(b)$ et $(b')$.} Par définition,
$\mathbf{Tor}_q^{A,\mathrm{II}}
(L_{\bullet\bullet}^{(1)},L_{\bullet\bullet}^{(2)})$ est un bicomplexe dont
le terme de degré $(h,k)$ est le $A$-module
\[
  \mathbf{Tor}_q^A
  \bigl(C^{-h}(\mathfrak U^{(1)},\mathscr P_\bullet^{(1)}),
        C^{-k}(\mathfrak U^{(2)},\mathscr P_\bullet^{(2)})\bigr).
\]
Soit $\Phi^{(i)}$ l'ensemble d'indices de $\mathfrak U^{(i)}$; par
définition, le complexe de modules
$C^r(\mathfrak U^{(i)},\mathscr P_\bullet^{(i)})$ ($r\geqslant0$) est somme
directe des complexes
$\Gamma(U_\rho^{(i)},\mathscr P_\bullet^{(i)})$, où $U_\rho^{(i)}$ est
l'intersection des $U_\xi^{(i)}$ pour $\xi\in\rho$, et $\rho$ parcourt
$\mathfrak P(\Phi^{(i)})$; donc le $A$-module
\[
  \mathbf{Tor}_q^A
  \bigl(C^{-h}(\mathfrak U^{(1)},\mathscr P_\bullet^{(1)}),
        C^{-k}(\mathfrak U^{(2)},\mathscr P_\bullet^{(2)})\bigr)
\]
est somme directe des $A$-modules
\[
  \mathbf{Tor}_q^A
  \bigl(\Gamma(U_\sigma^{(1)},\mathscr P_\bullet^{(1)}),
        \Gamma(U_\tau^{(2)},\mathscr P_\bullet^{(2)})\bigr),
\]
où $\sigma$ (resp. $\tau$) parcourt les éléments de
$\mathfrak P(\Phi^{(1)})$ (resp. $\mathfrak P(\Phi^{(2)})$) tels que
$\operatorname{Card}(\sigma)=-(h+1)$ (resp.
$\operatorname{Card}(\tau)=-(k+1)$). Comme $X^{(1)}$ et $X^{(2)}$ sont des
schémas, les $U_\rho^{(i)}$ sont affines, donc on a (6.4.1.1)
\[
  \mathbf{Tor}_q^A
  \bigl(\Gamma(U_\sigma^{(1)},\mathscr P_\bullet^{(1)}),
        \Gamma(U_\tau^{(2)},\mathscr P_\bullet^{(2)})\bigr)
  =\Gamma\bigl(U_\sigma^{(1)}\times_S U_\tau^{(2)},
    \mathscr{T}\!or_q^S
    (\mathscr P_\bullet^{(1)},\mathscr P_\bullet^{(2)})\bigr).
\]
\oldpage[III]{155}
On voit donc que ${}^{(b)}\mathscr E_{pq}^2$ est le $(-p)$-ème module de
cohomologie du complexe
\[
  L^\bullet(\Phi^{(1)},\Phi^{(2)};\mathscr S)
\]
des cochaînes bi-alternées sur $\Phi^{(1)}$ et $\Phi^{(2)}$ à valeurs dans le
système de coefficients
\[
  \mathscr S:(\sigma,\tau)\leadsto
  \Gamma\bigl(U_\sigma^{(1)}\times_S U_\tau^{(2)},
    \mathscr{T}\!or_q^S
    (\mathscr P_\bullet^{(1)},\mathscr P_\bullet^{(2)})\bigr)
\]
($0_{\mathrm I}$, 11.8.4). On sait alors ($0_{\mathrm I}$, 11.8.5 et 11.8.6)
que la cohomologie de ce complexe est la même que celle du complexe
$C^\bullet(\Phi^{(1)},\Phi^{(2)};\mathscr S)$ de toutes les cochaînes sur
$\Phi^{(1)}$ et $\Phi^{(2)}$ à valeurs dans $\mathscr S$, et aussi que celle du
complexe $P^\bullet(\Phi^{(1)},\Phi^{(2)};\mathscr S)$, dont les éléments sont
les combinaisons linéaires des
\[
  \lambda(\sigma,\tau)\in
  \Gamma\bigl(U_\sigma^{(1)}\times_S U_\tau^{(2)},
    \mathscr{T}\!or_q^S
    (\mathscr P_\bullet^{(1)},\mathscr P_\bullet^{(2)})\bigr),
\]
où $\sigma=(\alpha_0,\ldots,\alpha_h)$ et
$\tau=(\beta_0,\ldots,\beta_h)$ sont des suites ayant même nombre d'éléments.
Mais on a alors
\[
  U_\sigma^{(1)}\times_S U_\tau^{(2)}
  =(U_{\alpha_0}^{(1)}\times_S U_{\beta_0}^{(2)})\cap\cdots\cap
    (U_{\alpha_h}^{(1)}\times_S U_{\beta_h}^{(2)})
  \qquad\text{(I, 3.2.7)}.
\]
Si l'on désigne par $\mathfrak U$ le recouvrement de
$Z=X^{(1)}\times_S X^{(2)}$ par les ouverts affines
$U_\alpha^{(1)}\times_S U_\beta^{(2)}$, on voit finalement, compte tenu de ce
que $X^{(1)}\times_S X^{(2)}$ est un schéma, que l'on a, en vertu de (1.3.1),
\[
  {}^{(b)}\mathscr E_{pq}^2
  =\mathbf H^{-p}\bigl(X^{(1)}\times_S X^{(2)},
    \mathscr{T}\!or_q^S
    (\mathscr P_\bullet^{(1)},\mathscr P_\bullet^{(2)})\bigr).
\]

En second lieu,
$\mathbf{Tor}_q^A(L_{\bullet\bullet}^{(1)},L_{\bullet\bullet}^{(2)})$ est un
bicomplexe dont le terme de degré $(h,k)$ est la somme directe des $A$-modules
\[
  \mathbf{Tor}_q^A
  \bigl(C^{-h_1}(\mathfrak U^{(1)},\mathscr P_{k_1}^{(1)}),
        C^{-h_2}(\mathfrak U^{(2)},\mathscr P_{k_2}^{(2)})\bigr)
\]
tels que $h_1+h_2=h$ et $k_1+k_2=k$; explicitant les modules
$C^r(\mathfrak U^{(i)},\mathscr P_s^{(i)})$ comme ci-dessus, on voit encore
que ce terme est somme directe des $A$-modules
\[
  \Gamma\bigl(U_\sigma^{(1)}\times_S U_\tau^{(2)},
    \mathscr{T}\!or_q^S
    (\mathscr P_{k_1}^{(1)},\mathscr P_{k_2}^{(2)})\bigr),
\]
où $k_1+k_2=k$, et $\sigma$ (resp. $\tau$) parcourt les éléments de
$\mathfrak P(\Phi^{(1)})$ (resp. $\mathfrak P(\Phi^{(2)})$) tels que
$\operatorname{Card}(\sigma)+\operatorname{Card}(\tau)=-h-2$. Le terme
${}^{(b')}\mathscr E_{pq}^2$ que nous calculons est le $(-p)$-ème module de
cohomologie d'un bicomplexe $N^{\bullet\bullet}=(N^{hk})$, où le complexe
simple $N^{\bullet k}$ est le complexe de cochaînes bi-alternées sur
$\Phi^{(1)}$ et $\Phi^{(2)}$, à valeurs dans le système de coefficients
\[
  \mathscr S_k:(\sigma,\tau)\leadsto
  \Gamma\left(U_\sigma^{(1)}\times_S U_\tau^{(2)},
    \bigoplus_{k_1+k_2=k}
    \mathscr{T}\!or_q^S
    (\mathscr P_{-k_1}^{(1)},\mathscr P_{-k_2}^{(2)})\right),
\]
ces systèmes de coefficients formant un complexe $\mathscr S^\bullet$ où la
différentielle provient de celle du complexe simple associé au bicomplexe
$\mathscr{T}\!or_q^S
(\mathscr P_\bullet^{(1)},\mathscr P_\bullet^{(2)})$. On sait que la
cohomologie de $N^{\bullet\bullet}$ est la même que celle du bicomplexe
$C^\bullet(\Phi^{(1)},\Phi^{(2)};\mathscr S^\bullet)$
($0_{\mathrm I}$, 11.8.9), et aussi la même que celle du bicomplexe
$P^\bullet(\Phi^{(1)},\Phi^{(2)};\mathscr S^\bullet)$, où les éléments de
degrés $(h,k)$ sont les combinaisons linéaires des
\[
  \lambda(\sigma,\tau)\in
  \Gamma\left(U_\sigma^{(1)}\times_S U_\tau^{(2)},
    \bigoplus_{k_1+k_2=k}
    \mathscr{T}\!or_q^S
    (\mathscr P_{-k_1}^{(1)},\mathscr P_{-k_2}^{(2)})\right),
\]
$\sigma=(\alpha_0,\ldots,\alpha_h)$,
$\tau=(\beta_0,\ldots,\beta_h)$ étant des suites ayant même nombre
d'éléments ($0_{\mathrm I}$, 11.8.10). On voit alors comme ci-dessus que
${}^{(b')}\mathscr E_{pq}^2$ est le $(-p)$-ème module de cohomologie du
bicomplexe $C^\bullet(\mathfrak U,\mathscr Q^\bullet)$, où
$\mathscr Q^\bullet$ est le complexe simple associé au bicomplexe
$\mathscr{T}\!or_q^S
(\mathscr P_\bullet^{(1)},\mathscr P_\bullet^{(2)})$ de
$\mathscr O_Z$-Modules. Avec les conventions faites dans (6.6.1), on a donc
\[
  {}^{(b')}\mathscr E_{pq}^2
  =\mathbf H^{-p}\bigl(X^{(1)}\times_S X^{(2)},
    \mathscr{T}\!or_q^S
    (\mathscr P_\bullet^{(1)},\mathscr P_\bullet^{(2)})\bigr).
\]
\end{env}

\oldpage[III]{156}
\begin{env}[6.6.5]
\emph{Suites spectrales $(c)$ et $(d)$.} Par définition, le complexe
$\mathrm H_n^{\mathrm{II}}(L_{\bullet\bullet}^{(i)})$ a pour terme de degré
$h$ le $A$-module
$C^{-h}(\mathfrak U^{(i)},\mathscr H_n(\mathscr P_\bullet^{(i)}))$, en vertu
de l'exactitude du foncteur $C^{-h}$. On a donc, par définition de l'hypertor
de deux Modules relatif à deux recouvrements (6.6.2)
\[
  {}^{(c)}\mathscr E_{pq}^2
  =\bigoplus_{q_1+q_2=q}
   \mathbf{Tor}_p^S
   \bigl(\mathfrak U^{(1)},\mathfrak U^{(2)};
         \mathscr H_{q_1}(\mathscr P_\bullet^{(1)}),
         \mathscr H_{q_2}(\mathscr P_\bullet^{(2)})\bigr).
\]
Enfin, par définition,
$\mathbf{Tor}_q^{A,\mathrm I}
(L_{\bullet\bullet}^{(1)},L_{\bullet\bullet}^{(2)})$ est un bicomplexe dont
le terme de degré $(h,k)$ est le $A$-module
$\mathbf{Tor}_q^S
(\mathfrak U^{(1)},\mathfrak U^{(2)};
 \mathscr P_h^{(1)},\mathscr P_k^{(2)})$. On a donc
\[
  {}^{(d)}\mathscr E_{pq}^2
  =\mathrm H_p\bigl(
    \mathbf{Tor}_q^S
    (\mathfrak U^{(1)},\mathfrak U^{(2)};
     \mathscr P_\bullet^{(1)},\mathscr P_\bullet^{(2)})\bigr).
\]
\end{env}

\begin{env}[6.6.6]
La théorie de l'hyperhomologie des foncteurs de bicomplexes
($0_{\mathrm I}$, 11.7.3) montre, comme dans (6.3.4), que, pour toute
résolution \emph{plate} de Cartan-Eilenberg $M_{\bullet\bullet}^{(i)}$ de
$L_{\bullet\bullet}^{(i)}$ (dans la catégorie des complexes de modules
limités inférieurement) ($i=1,2$), on a des isomorphismes canoniques de
bi-$\partial$-foncteurs
\[
\begin{aligned}
  \mathbf{Tor}_\bullet^S
  (\mathfrak U^{(1)},\mathfrak U^{(2)};
   \mathscr P_\bullet^{(1)},\mathscr P_\bullet^{(2)})
  &\xrightarrow{\sim}
   \mathrm H_\bullet
   (M_{\bullet\bullet}^{(1)}\otimes_A M_{\bullet\bullet}^{(2)})\\
  &\xrightarrow{\sim}
   \mathrm H_\bullet
   (M_{\bullet\bullet}^{(1)}\otimes_A L_{\bullet\bullet}^{(2)})
   \xrightarrow{\sim}
   \mathrm H_\bullet
   (L_{\bullet\bullet}^{(1)}\otimes_A M_{\bullet\bullet}^{(2)}).
\end{aligned}
\tag{6.6.6.1}\label{III.6.6.6.1.fr}
\]
\end{env}

\begin{env}[6.6.7]
Nous allons maintenant montrer que l'hypertor global défini dans (6.6.2), et
les six suites spectrales correspondantes, ne dépendent pas des recouvrements
ouverts affines finis $\mathfrak U^{(i)}$ qui ont servi à les définir (à des
isomorphismes canoniques près). Il suffira pour cela de montrer que si
$\mathfrak V^{(i)}$ sont deux autres recouvrements de même nature, tels que
$\mathfrak V^{(i)}$ soit \emph{plus fin} que $\mathfrak U^{(i)}$ pour
$i=1,2$, alors on a des \emph{isomorphismes canoniques} de foncteurs spectraux
\[
  {}^{(t)}\mathscr E
  (\mathfrak U^{(1)},\mathfrak U^{(2)};
   \mathscr P_\bullet^{(1)},\mathscr P_\bullet^{(2)})
  \xrightarrow{\sim}
  {}^{(t)}\mathscr E
  (\mathfrak V^{(1)},\mathfrak V^{(2)};
   \mathscr P_\bullet^{(1)},\mathscr P_\bullet^{(2)})
  \tag{6.6.7, $t$}\label{III.6.6.7.t.fr}
\]
où $t$ est remplacé par $a$, $b$, $a'$, $b'$, $c$ ou $d$.

Or, on a pour $i=1,2$ des homomorphismes de bicomplexes
\[
  C^\bullet(\mathfrak U^{(i)},\mathscr P_\bullet^{(i)})
  \longrightarrow
  C^\bullet(\mathfrak V^{(i)},\mathscr P_\bullet^{(i)})
\]
bien définis à homotopies près (G, II, 5.7.1); il en résulte déjà des
homomorphismes (6.6.7, $t$) canoniquement définis et compatibles avec les
opérateurs bords dans les aboutissements ($0_{\mathrm I}$, 11.3.2). En outre,
le calcul des termes $E_2$ des suites spectrales $(a)$, $(b)$, $(a')$, $(b')$,
montre que pour ces suites spectrales l'homomorphisme (6.6.7, $t$) est un
isomorphisme sur les termes $E_2$; comme ces suites spectrales sont
birégulières, on voit que (6.6.7, $t$) est un isomorphisme pour ces trois
foncteurs spectraux, donc un \emph{isomorphisme de bi-$\partial$-foncteurs}
pour leur aboutissement commun ($0_{\mathrm I}$, 11.1.5).

En particulier, pour des $\mathscr O_{X^{(i)}}$-Modules quasi-cohérents
$\mathscr F^{(i)}$ ($i=1,2$), l'homomorphisme canonique
\[
  \mathbf{Tor}_\bullet^S
  (\mathfrak U^{(1)},\mathfrak U^{(2)};
   \mathscr F^{(1)},\mathscr F^{(2)})
  \longrightarrow
  \mathbf{Tor}_\bullet^S
  (\mathfrak V^{(1)},\mathfrak V^{(2)};
   \mathscr F^{(1)},\mathscr F^{(2)})
\]
est bijectif; vu le calcul de (6.6.5), on voit que (6.6.7, $t$) est aussi un
isomorphisme des termes $E_2$ pour $t=c$ et $t=d$. On conclut comme ci-dessus
que (6.6.7, $t$) est aussi un isomorphisme de suites spectrales pour $t=c$ et
$t=d$.
\oldpage[III]{157}
On peut considérer que les isomorphismes (6.6.7, $t$) définissent des systèmes
inductifs de foncteurs spectraux sur l'ensemble filtrant des couples
$(\mathfrak U^{(1)},\mathfrak U^{(2)})$ de recouvrements ouverts affines finis
de $X^{(1)}$ et $X^{(2)}$. Nous désignerons par
\[
  {}^{(t)}\mathscr E
  (X^{(1)},X^{(2)};
   \mathscr P_\bullet^{(1)},\mathscr P_\bullet^{(2)})
  \quad\text{ou}\quad
  {}^{(t)}\mathscr E^S
  (X^{(1)},X^{(2)};
   \mathscr P_\bullet^{(1)},\mathscr P_\bullet^{(2)})
\]
la limite inductive de ce système, par
$\mathbf{Tor}_\bullet^S
(X^{(1)},X^{(2)};\mathscr P_\bullet^{(1)},\mathscr P_\bullet^{(2)})$
l'aboutissement de ce foncteur spectral, que nous appellerons l'\emph{hypertor
global} des deux complexes $\mathscr P_\bullet^{(1)}$ et
$\mathscr P_\bullet^{(2)}$; si $\mathscr P_\bullet^{(1)}$ et
$\mathscr P_\bullet^{(2)}$ sont réduits à leurs termes de degré 0,
$\mathscr F^{(1)}$ et $\mathscr F^{(2)}$, nous écrirons
\[
  \mathbf{Tor}_\bullet^S
  (X^{(1)},X^{(2)};\mathscr F^{(1)},\mathscr F^{(2)}),
\]
et conformément aux conventions générales,
\[
  \mathbf{Tor}_q^S
  (X^{(1)},X^{(2)};\mathscr P_\bullet^{(1)},\mathscr P_\bullet^{(2)})
\]
sera donc le bicomplexe des
\[
  \mathbf{Tor}_q^S
  (X^{(1)},X^{(2)};\mathscr P_h^{(1)},\mathscr P_k^{(2)}).
\]
\end{env}

\begin{env}[6.6.8]
Les hypothèses étant celles de (6.6.1), considérons maintenant deux
$S$-morphismes $f_i:X^{(i)}\to Y^{(i)}$, où
$Y^{(i)}=\operatorname{Spec}(B_i)$ est un $S$-schéma \emph{affine}, $B_i$
étant donc une $A$-algèbre ($i=1,2$); cela définit donc un $A$-homomorphisme
$B_i\to\Gamma(X^{(i)},\mathscr O_{X^{(i)}})$ (I, 2.2.4), et par suite chacun
des $L_{\bullet\bullet}^{(i)}$ définis dans (6.6.2) est un bicomplexe de
$B_i$-modules; on en conclut que
$L_{\bullet\bullet}^{(1)}\otimes_A L_{\bullet\bullet}^{(2)}$ est un
quadricomplexe de $(B_1\otimes_A B_2)$-modules, et ses six foncteurs spectraux
d'hyperhomologie peuvent donc être considérés comme prenant leurs valeurs
dans la catégorie des suites spectrales de $(B_1\otimes_A B_2)$-modules. Si
l'on pose
$Y=Y^{(1)}\times_S Y^{(2)}=\operatorname{Spec}(B_1\otimes_A B_2)$, on peut
considérer les $\mathscr O_Y$-Modules quasi-cohérents associés à ces modules
(I, 1.3.4); nous noterons
${}^{(t)}\mathscr E
(f_1,f_2;\mathscr P_\bullet^{(1)},\mathscr P_\bullet^{(2)})$
(pour $t=a$, $b$, $a'$, $b'$, $c$ ou $d$) les six suites spectrales
\[
  \bigl({}^{(t)}\mathscr E
  (X^{(1)},X^{(2)};
   \mathscr P_\bullet^{(1)},\mathscr P_\bullet^{(2)})\bigr)^{\sim}
\]
de $\mathscr O_Y$-Modules, et
$\mathfrak{Tor}_\bullet^S
(f_1,f_2;\mathscr P_\bullet^{(1)},\mathscr P_\bullet^{(2)})$ leur
aboutissement commun
\[
  \bigl(\mathbf{Tor}_\bullet^S
  (X^{(1)},X^{(2)};
   \mathscr P_\bullet^{(1)},\mathscr P_\bullet^{(2)})\bigr)^{\sim}.
\]
On le notera
$\mathfrak{Tor}_\bullet^S(f_1,f_2;\mathscr F^{(1)},\mathscr F^{(2)})$
lorsque $\mathscr P_\bullet^{(i)}$ est réduit à son terme de degré 0,
$\mathscr F^{(i)}$ ($i=1,2$).
\end{env}

\subsection{Foncteurs hypertor globaux de complexes de Modules quasi-cohérents
et suites spectrales de Künneth : cas général.}

\begin{env}[6.7.1]
Nous allons maintenant généraliser les définitions de (6.6.8) au cas où $S$
est un préschéma quelconque, $X^{(i)}$, $Y^{(i)}$ des $S$-préschémas et
$f_i:X^{(i)}\to Y^{(i)}$ des morphismes \emph{séparés et quasi-compacts}. Il
s'agit alors, pour tout couple de complexes $\mathscr P_\bullet^{(i)}$ de
$\mathscr O_{X^{(i)}}$-Modules quasi-cohérents, limités inférieurement
($i=1,2$), de définir pour tout $n$ un $\mathscr O_Y$-Module quasi-cohérent
$\mathfrak{Tor}_n^S
(f_1,f_2;\mathscr P_\bullet^{(1)},\mathscr P_\bullet^{(2)})$ ainsi que 6
foncteurs spectraux, se réduisant aux définitions de (6.6.8) lorsque $S$,
$Y^{(1)}$ et $Y^{(2)}$ sont affines (on pose
$Y=Y^{(1)}\times_S Y^{(2)}$).

Supposons d'abord $S=\operatorname{Spec}(A)$ \emph{affine}, mais $Y^{(1)}$
et $Y^{(2)}$ quelconques; soit $W^{(i)}$ un ouvert affine de $Y^{(i)}$;
$f_i^{-1}(W^{(i)})$ est alors un $S$-schéma quasi-compact,
$W=W^{(1)}\times_S W^{(2)}$ un ouvert affine de $Y$; soit
$f'_i:f_i^{-1}(W^{(i)})\to W^{(i)}$ la restriction de $f_i$, et
$\mathscr P'_\bullet{}^{(i)}$ la restriction
$\mathscr P_\bullet^{(i)}|f_i^{-1}(W^{(i)})$ ($i=1,2$). On a alors d'après
(6.6.8) les suites spectrales
${}^{(t)}\mathscr E
(f'_1,f'_2;\mathscr P'_\bullet{}^{(1)},\mathscr P'_\bullet{}^{(2)})$
de $(\mathscr O_Y|W)$-Modules quasi-cohérents, et il s'agit de vérifier
qu'elles satisfont aux conditions de recollement ($0_{\mathrm I}$, 3.3.1).
On est aussitôt ramené au cas où $Y^{(i)}=\operatorname{Spec}(B_i)$ est affine
et où $W^{(i)}=D(g_i)$, où $g_i\in B_i$, de sorte que
$W=D(g_1\otimes g_2)$
\oldpage[III]{158}
dans $Y=\operatorname{Spec}(B_1\otimes_A B_2)$ (II, 4.3.2.1); si
$X'^{(i)}=f_i^{-1}(W^{(i)})$, il s'agit d'établir un isomorphisme canonique
de foncteurs spectraux
\[
  {}^{(t)}\mathscr E
  (X'^{(1)},X'^{(2)};
   \mathscr P'_\bullet{}^{(1)},\mathscr P'_\bullet{}^{(2)})
  \xrightarrow{\sim}
  {}^{(t)}\mathscr E
  (X^{(1)},X^{(2)};
   \mathscr P_\bullet^{(1)},\mathscr P_\bullet^{(2)})
  \otimes_B B_g
  \tag{6.7.1.1}\label{III.6.7.1.1.fr}
\]
où l'on a posé $B=B_1\otimes_A B_2$ et $g=g_1\otimes g_2$. Pour cela,
partons de recouvrements ouverts affines finis $\mathfrak U^{(i)}$ de
$X^{(i)}$ ($i=1,2$), et soit $\mathfrak U'^{(i)}$ la trace de
$\mathfrak U^{(i)}$ sur $X'^{(i)}$, qui est encore formé d'ouverts affines
(I, 5.5.10); de façon précise, on a
\[
  C^\bullet(\mathfrak U'^{(i)},\mathscr P'_\bullet{}^{(i)})
  =C^\bullet(\mathfrak U^{(i)},\mathscr P_\bullet^{(i)})
   \otimes_{B_i}(B_i)_{g_i}.
\]
Si on pose
$L'_{\bullet\bullet}{}^{(i)}
=C^\bullet(\mathfrak U'^{(i)},\mathscr P'_\bullet{}^{(i)})$, on a donc
\[
  L'_{\bullet\bullet}{}^{(1)}\otimes_A L'_{\bullet\bullet}{}^{(2)}
  =\bigl(L_{\bullet\bullet}^{(1)}\otimes_{B_1}(B_1)_{g_1}\bigr)
   \otimes_A
   \bigl(L_{\bullet\bullet}^{(2)}\otimes_{B_2}(B_2)_{g_2}\bigr);
\]
comme on a $B_g=(B_1)_{g_1}\otimes_A(B_2)_{g_2}$ à un isomorphisme
canonique près, on a, à un isomorphisme canonique près,
\[
  L'_{\bullet\bullet}{}^{(1)}\otimes_A L'_{\bullet\bullet}{}^{(2)}
  =(L_{\bullet\bullet}^{(1)}\otimes_A L_{\bullet\bullet}^{(2)})
    \otimes_B B_g.
\]
Si $M_{\bullet\bullet\bullet}^{(i)}$ est une résolution projective de
Cartan-Eilenberg de $L_{\bullet\bullet}^{(i)}$, qu'on peut supposer formée de
$B_i$-modules, il résulte du fait que $(B_i)_{g_i}$ est \emph{plat} sur $B_i$
que
$M'_{\bullet\bullet\bullet}{}^{(i)}
=M_{\bullet\bullet\bullet}^{(i)}\otimes_{B_i}(B_i)_{g_i}$ est une résolution
projective de Cartan-Eilenberg du bicomplexe
$L'_{\bullet\bullet}{}^{(i)}$; en outre, on a
\[
  M'_{\bullet\bullet\bullet}{}^{(1)}\otimes_A M'_{\bullet\bullet\bullet}{}^{(2)}
  =(M_{\bullet\bullet\bullet}^{(1)}\otimes_A M_{\bullet\bullet\bullet}^{(2)})
    \otimes_B B_g.
\]
L'isomorphisme cherché (6.7.1.1) résulte alors aussitôt des définitions de
l'hyperhomologie d'un bicomplexe, et de l'exactitude du foncteur
$G\otimes_B B_g$ en le $B$-module $G$.
\end{env}

\begin{env}[6.7.2]
Supposons maintenant $S$ quelconque, et soient
$u_i:Y^{(i)}\to S$ les morphismes structuraux ($i=1,2$). Soit $(S_\alpha)$ un
recouvrement ouvert affine de $S$, posons
$Y_\alpha^{(i)}=u_i^{-1}(S_\alpha)$,
$X_\alpha^{(i)}=f_i^{-1}(Y_\alpha^{(i)})$, et soit
$f_{i\alpha}:X_\alpha^{(i)}\to Y_\alpha^{(i)}$ la restriction de $f_i$, qui
est un morphisme \emph{séparé et quasi-compact}. Les
$Y_\alpha=Y_\alpha^{(1)}\times_{S_\alpha}Y_\alpha^{(2)}$ forment un
recouvrement ouvert de $Y$, et sur chaque $Y_\alpha$ sont définis par (6.7.1)
des foncteurs spectraux
\[
  {}^{(t)}\mathscr E^{S_\alpha}
  \bigl(f_{1\alpha},f_{2\alpha};
    \mathscr P_\bullet^{(1)}|X_\alpha^{(1)},
    \mathscr P_\bullet^{(2)}|X_\alpha^{(2)}\bigr);
\]
il s'agit encore de montrer que ces foncteurs vérifient les conditions de
recollement.

On se ramène aussitôt à la situation suivante :
$S=\operatorname{Spec}(A)$ est affine, $S'=D(h)$, où $h\in A$, et
$u_i(Y^{(i)})\subset S'$; on peut en outre supposer
$Y^{(i)}=\operatorname{Spec}(B_i)$ affine; il s'agit de définir des
isomorphismes canoniques
\[
  {}^{(t)}\mathscr E^S
  (X^{(1)},X^{(2)};
   \mathscr P_\bullet^{(1)},\mathscr P_\bullet^{(2)})
  \xrightarrow{\sim}
  {}^{(t)}\mathscr E^{S'}
  (X^{(1)},X^{(2)};
   \mathscr P_\bullet^{(1)},\mathscr P_\bullet^{(2)}).
  \tag{6.7.2.1}\label{III.6.7.2.1.fr}
\]
Or, avec les notations de (6.6.2), les $L_{\bullet\bullet}^{(i)}$ sont formés
de $A_h$-modules, et l'on a donc
$L_{\bullet\bullet}^{(1)}\otimes_{A_h}L_{\bullet\bullet}^{(2)}
=L_{\bullet\bullet}^{(1)}\otimes_A L_{\bullet\bullet}^{(2)}$ à un
isomorphisme canonique près; comme on peut prendre une résolution projective
de Cartan-Eilenberg $M_{\bullet\bullet\bullet}^{(i)}$ de
$L_{\bullet\bullet}^{(i)}$ formée de $A_h$-modules, cela donne aussitôt
l'isomorphisme canonique cherché.
\end{env}

Nous avons en résumé démontré le

\begin{theorem}[6.7.3]
Soient $S$ un préschéma, $f_i:X_i\to Y_i$ un $S$-morphisme séparé et
quasi-compact de $S$-préschémas, $\mathscr P_\bullet^{(i)}$ un complexe de
$\mathscr O_{X^{(i)}}$-Modules quasi-cohérents limité inférieurement
($i=1,2$); on pose $Y=Y^{(1)}\times_S Y^{(2)}$. Il existe un
bi-$\partial$-foncteur
  $\mathfrak{Tor}_\bullet^S
(f_1,f_2;\mathscr P_\bullet^{(1)},\mathscr P_\bullet^{(2)})$
à valeurs dans la catégorie des $\mathscr O_Y$-Modules quasi-cohérents, tel
que si $V^{(i)}$ est un ouvert affine de $Y^{(i)}$ ($i=1,2$) et
$V=V^{(1)}\times_S V^{(2)}$, on ait
\[
  \mathfrak{Tor}_\bullet^S
  (f_1,f_2;\mathscr P_\bullet^{(1)},\mathscr P_\bullet^{(2)})|V
  =\Bigl(\mathbf{Tor}_\bullet^S
    \bigl(f_1^{-1}(V^{(1)}),f_2^{-1}(V^{(2)});
      \mathscr P_\bullet^{(1)}|f_1^{-1}(V^{(1)}),
      \mathscr P_\bullet^{(2)}|f_2^{-1}(V^{(2)})\bigr)\Bigr)^{\sim}.
\]
\oldpage[III]{159}
Ce bifoncteur est l'aboutissement de six foncteurs spectraux biréguliers
\[
  {}^{(t)}\mathscr E
  (f_1,f_2;\mathscr P_\bullet^{(1)},\mathscr P_\bullet^{(2)})
  \qquad(t=a,b,a',b',c,d)
\]
dont les termes $E_2$ sont donnés par
\[
\begin{aligned}
  {}^{(a)}\mathscr E_{pq}^2
  &=\bigoplus_{q_1+q_2=q}
    \mathscr{T}\!or_p^S
    \bigl(\mathscr H^{-q_1}(f_1,\mathscr P_\bullet^{(1)}),
          \mathscr H^{-q_2}(f_2,\mathscr P_\bullet^{(2)})\bigr),\\
  {}^{(b)}\mathscr E_{pq}^2
  &=\mathscr H^{-p}\bigl(f_1\times_S f_2,
    \mathfrak{Tor}_q^S
    (\mathscr P_\bullet^{(1)},\mathscr P_\bullet^{(2)})\bigr),\\
  {}^{(a')}\mathscr E_{pq}^2
  &=\bigoplus_{q_1+q_2=q}
    \mathfrak{Tor}_p^S
    \bigl(\mathscr H^{-q_1}(f_1,\mathscr P_\bullet^{(1)}),
          \mathscr H^{-q_2}(f_2,\mathscr P_\bullet^{(2)})\bigr),\\
  {}^{(b')}\mathscr E_{pq}^2
  &=\mathscr H^{-p}\bigl(f_1\times_S f_2,
    \mathscr{T}\!or_q^S
    (\mathscr P_\bullet^{(1)},\mathscr P_\bullet^{(2)})\bigr),\\
  {}^{(c)}\mathscr E_{pq}^2
  &=\bigoplus_{q_1+q_2=q}
    \mathfrak{Tor}_p^S
    \bigl(f_1,f_2;
      \mathscr H_{q_1}(\mathscr P_\bullet^{(1)}),
      \mathscr H_{q_2}(\mathscr P_\bullet^{(2)})\bigr),\\
  {}^{(d)}\mathscr E_{pq}^2
  &=\mathscr H_p\bigl(\mathfrak{Tor}_q^S
    (f_1,f_2;\mathscr P_\bullet^{(1)},\mathscr P_\bullet^{(2)})\bigr).
\end{aligned}
\]
On dit que les suites spectrales $(a)$ et $(b)$ sont les \emph{suites
spectrales de Künneth}.

On notera que les suites spectrales $(a)$ et $(a')$ (resp. $(b)$ et $(b')$)
sont identiques lorsque $\mathscr P_\bullet^{(1)}$ et
$\mathscr P_\bullet^{(2)}$ se réduisent à leurs termes de degré 0; dans ce cas,
les suites $(c)$ et $(d)$ sont dégénérées et sont donc sans intérêt.
\end{theorem}

\begin{remark}[6.7.4]
Les hypertor globaux que nous avons définis ci-dessus comprennent comme cas
particuliers, à la fois les \emph{Modules d'hypercohomologie} définis dans
(6.2.1) et les \emph{hypertor locaux} définis dans (6.5.3). Montrons que l'on
a, pour tout morphisme $f:X\to Y$ \emph{quasi-compact et séparé} et tout
complexe $\mathscr P_\bullet$ de $\mathscr O_X$-Modules quasi-cohérents,
limité inférieurement, un isomorphisme canonique de $\partial$-foncteurs en
$\mathscr P_\bullet$
\[
  \mathfrak{Tor}_n^Y(f,1_Y;\mathscr P_\bullet,\mathscr O_Y)
  \xrightarrow{\sim}
  \mathscr H^{-n}(f,\mathscr P_\bullet)
  \qquad\text{(pour tout $n\in\mathbf Z$)}.
  \tag{6.7.4.1}\label{III.6.7.4.1.fr}
\]
En effet, les méthodes de recollement de (6.7.2) ramènent aussitôt au cas où
$Y$ est affine; on peut alors, en vertu de (6.2.2), calculer les deux membres
de (6.7.4.1) à l'aide d'un même recouvrement fini $\mathfrak U$ de $X$ par des
ouverts affines, et (pour le premier membre) du recouvrement de $Y$ formé de
$Y$ lui-même; avec les notations de (6.6.2), le bicomplexe
$L_{\bullet\bullet}^{(2)}$ est alors réduit à son terme de degrés $(0,0)$,
égal à $A$, et la conclusion résulte de ($0_{\mathrm I}$, 11.7.5). Pour une
généralisation de ce résultat, voir (6.7.7); mais on notera que lorsque dans
le premier membre de (6.7.4.1), on remplace $\mathscr O_Y$ par un
$\mathscr O_Y$-Module quasi-cohérent \emph{quelconque} $\mathscr F$, on n'a
plus en général un isomorphisme avec
$\mathscr H^{-n}(f,\mathscr P_\bullet\otimes_{\mathscr O_Y}\mathscr F)$, bien
que, dans le calcul précédent, le bicomplexe
$L_{\bullet\bullet}^{(1)}\otimes_A L_{\bullet\bullet}^{(2)}$ s'identifie
encore au bicomplexe
$C^\bullet(\mathfrak U,\mathscr P_\bullet\otimes_Y\mathscr F)$.

D'autre part, on a un isomorphisme canonique de bi-$\partial$-foncteurs
\[
  \mathfrak{Tor}_\bullet^S
  (1_{X^{(1)}},1_{X^{(2)}};
   \mathscr P_\bullet^{(1)},\mathscr P_\bullet^{(2)})
  \xrightarrow{\sim}
  \mathscr{T}\!or_\bullet^S
  (\mathscr P_\bullet^{(1)},\mathscr P_\bullet^{(2)}).
  \tag{6.7.4.2}\label{III.6.7.4.2.fr}
\]
En effet, on se ramène encore, par (6.7.1) et (6.7.2), au cas où $S$ et les
$X^{(i)}$ sont affines; dans le calcul du premier membre de (6.7.4.2), on peut
alors prendre pour
\oldpage[III]{160}
recouvrement $\mathfrak U^{(i)}$ la famille réduite au seul élément $X^{(i)}$,
de sorte qu'avec les notations de (6.6.2), $L_{\bullet\bullet}^{(i)}$ se réduit
à $\Gamma(X^{(i)},\mathscr P_\bullet^{(i)})$ (considéré comme bicomplexe dont
les termes de premier degré $\ne0$ sont nuls), et l'égalité des deux membres
de (6.7.4.2) résulte de (6.4.1.1) et (6.3.1).
\end{remark}

\begin{proposition}[6.7.5]
Soit $u:\mathscr P_\bullet^{(1)}\to\mathscr Q_\bullet^{(1)}$ un
homomorphisme de complexes de $\mathscr O_{X^{(1)}}$-Modules quasi-cohérents,
limités inférieurement, tel que l'homomorphisme
\[
  \mathscr H_\bullet(u):
  \mathscr H_\bullet(\mathscr P_\bullet^{(1)})
  \longrightarrow
  \mathscr H_\bullet(\mathscr Q_\bullet^{(1)})
\]
déduit de $u$ soit un isomorphisme. Alors les homomorphismes
\[
  {}^{(t)}\mathscr E
  (f_1,f_2;\mathscr P_\bullet^{(1)},\mathscr P_\bullet^{(2)})
  \longrightarrow
  {}^{(t)}\mathscr E
  (f_1,f_2;\mathscr Q_\bullet^{(1)},\mathscr P_\bullet^{(2)})
\]
déduits de $u$ sont des isomorphismes pour $t=a$, $t=b$ et $t=c$.
\end{proposition}

\begin{proof}
L'assertion relative à la suite spectrale $(c)$ résulte de ce que cette suite
est birégulière et de ce que l'homomorphisme considéré est un isomorphisme pour
les termes $E_2$ par hypothèse ($0_{\mathrm I}$, 11.1.5). Ceci montre déjà que
\[
  \mathfrak{Tor}_\bullet^S
  (f_1,f_2;\mathscr P_\bullet^{(1)},\mathscr P_\bullet^{(2)})
  \longrightarrow
  \mathfrak{Tor}_\bullet^S
  (f_1,f_2;\mathscr Q_\bullet^{(1)},\mathscr P_\bullet^{(2)})
\]
est un isomorphisme. Appliquant les relations (6.7.4.1) et (6.7.4.2) on voit
d'abord que les homomorphismes
\[
  \mathscr H^{-n}(f_1,\mathscr P_\bullet^{(1)})
  \longrightarrow
  \mathscr H^{-n}(f_1,\mathscr Q_\bullet^{(1)})
  \quad\text{et}\quad
  \mathscr{T}\!or_n^S
  (\mathscr P_\bullet^{(1)},\mathscr P_\bullet^{(2)})
  \longrightarrow
  \mathscr{T}\!or_n^S
  (\mathscr Q_\bullet^{(1)},\mathscr P_\bullet^{(2)})
\]
déduits de $u$ sont des isomorphismes. L'assertion relative aux suites $(a)$
et $(b)$ résulte alors de ce que ces suites sont birégulières (6.7.3) et que
les homomorphismes considérés sont bijectifs pour les termes $E_2$
($0_{\mathrm I}$, 11.1.5).

Notons d'autre part que, si
$u:\mathscr P_\bullet^{(1)}\to\mathscr Q_\bullet^{(1)}$ est un
\emph{homotopisme}, on en déduit des isomorphismes canoniques
\[
  {}^{(t)}\mathscr E
  (f_1,f_2;\mathscr P_\bullet^{(1)},\mathscr P_\bullet^{(2)})
  \longrightarrow
  {}^{(t)}\mathscr E
  (f_1,f_2;\mathscr Q_\bullet^{(1)},\mathscr P_\bullet^{(2)})
\]
pour les six suites spectrales. En effet, si $S$ et les $Y^{(i)}$ sont affines,
on déduit de $u$ un homotopisme de bicomplexes
\[
  C^\bullet(\mathfrak U^{(1)},\mathscr P_\bullet^{(1)})
  \longrightarrow
  C^\bullet(\mathfrak U^{(1)},\mathscr Q_\bullet^{(1)}),
\]
et la proposition résulte de la théorie générale de l'hyperhomologie
($0_{\mathrm I}$, 11.3.2); le passage au cas général se fait par recollement,
en utilisant le fait que, d'un homotopisme de complexes, on déduit un
homotopisme de résolutions projectives de Cartan-Eilenberg de ces complexes
(M, XVII, 1.2).
\end{proof}

\begin{proposition}[6.7.6]
Supposons que le complexe $\mathscr P_\bullet^{(1)}$ ou le complexe
$\mathscr P_\bullet^{(2)}$ soit formé de Modules $S$-plats (les deux complexes
étant limités inférieurement). On a alors un isomorphisme canonique de
bi-$\partial$-foncteurs
\[
  \mathfrak{Tor}_n^S
  (f_1,f_2;\mathscr P_\bullet^{(1)},\mathscr P_\bullet^{(2)})
  \xrightarrow{\sim}
  \mathscr H^{-n}\bigl(f_1\times_S f_2,
    \mathscr P_\bullet^{(1)}\otimes_S\mathscr P_\bullet^{(2)}\bigr).
  \tag{6.7.6.1}\label{III.6.7.6.1.fr}
\]
\end{proposition}

\begin{proof}
Supposons d'abord $S$, $Y^{(1)}$ et $Y^{(2)}$ affines, de sorte qu'on est
dans la situation de (6.6.2), dont nous conservons les notations. Supposons par
exemple que $\mathscr P_\bullet^{(1)}$ soit formé de Modules $S$-plats, et
calculons l'hypertor en utilisant la remarque (6.6.6) : c'est donc l'homologie
de $L_{\bullet\bullet}^{(1)}\otimes_A M_{\bullet\bullet\bullet}^{(2)}$, où
$M_{\bullet\bullet\bullet}^{(2)}$ est une résolution projective de
Cartan-Eilenberg de $L_{\bullet\bullet}^{(2)}$, au sens de
($0_{\mathrm I}$, 11.7.1). D'autre part, les modules
$L_{ij}^{(1)}$ sont \emph{plats} sur $A$ en vertu de l'hypothèse (I, 4.15.1);
on déduit alors de ($0_{\mathrm I}$, 11.7.5) un isomorphisme canonique
\[
  \mathbf{Tor}_\bullet^S
  (\mathfrak U^{(1)},\mathfrak U^{(2)};
   \mathscr P_\bullet^{(1)},\mathscr P_\bullet^{(2)})
  \xrightarrow{\sim}
  \mathscr H_\bullet
  (L_{\bullet\bullet}^{(1)}\otimes_A L_{\bullet\bullet}^{(2)}).
  \tag{6.7.6.2}\label{III.6.7.6.2.fr}
\]
On a d'autre part un homomorphisme naturel de bicomplexes de
$L_{\bullet\bullet}^{(1)}\otimes_A L_{\bullet\bullet}^{(2)}$ dans
$C^\bullet(\mathfrak U,\mathscr Q_\bullet)$, où $\mathfrak U$ est le
recouvrement de $Z=X^{(1)}\times_S X^{(2)}$ par les ouverts affines
\oldpage[III]{161}
$U_\alpha^{(1)}\times_S U_\beta^{(2)}$ et
$\mathscr Q_\bullet
=\mathscr P_\bullet^{(1)}\otimes_S\mathscr P_\bullet^{(2)}$ (considéré comme
complexe simple pour le degré total); en effet, la définition de cet
homomorphisme a en substance été donnée au cours du calcul de la suite $(b')$
dans (6.6.4), pour $q=0$; il suffit simplement (en gardant les notations de
(6.6.4)) de tenir compte de ce qu'il y a d'une part un homomorphisme naturel
du complexe $N^{\bullet k}$ dans le complexe
$C^\bullet(\Phi^{(1)},\Phi^{(2)};\mathscr S_k)$
($0_{\mathrm I}$, 11.8.5), d'autre part un homomorphisme naturel de ce dernier
complexe dans le complexe
$P^\bullet(\Phi^{(1)},\Phi^{(2)};\mathscr S_k)$
($0_{\mathrm I}$, 11.8.6), et enfin un homomorphisme naturel de ce dernier
complexe de cochaînes dans le sous-complexe des cochaînes alternées
($0_{\mathrm I}$, 11.8.7). Par ailleurs, l'homomorphisme de bicomplexes ainsi
défini donne un isomorphisme en homologie, comme on l'a vu en (6.6.4); on a
donc, en composant avec (6.7.6.2), obtenu un isomorphisme
\[
  \mathbf{Tor}_n^S
  (\mathfrak U^{(1)},\mathfrak U^{(2)};
   \mathscr P_\bullet^{(1)},\mathscr P_\bullet^{(2)})
  \xrightarrow{\sim}
  \mathbf H^{-n}
  (\mathfrak U,\mathscr P_\bullet^{(1)}\otimes_S\mathscr P_\bullet^{(2)}).
  \tag{6.7.6.3}\label{III.6.7.6.3.fr}
\]
Il faut ensuite prouver que l'isomorphisme ainsi défini ne dépend pas des
recouvrements ouverts choisis (le second membre de (6.7.6.3) étant
canoniquement isomorphe à
$\mathbf H^{-n}
(X^{(1)}\times_S X^{(2)},
 \mathscr P_\bullet^{(1)}\otimes_S\mathscr P_\bullet^{(2)})$ par (6.2.2));
cela se fait à l'aide de (6.6.7) en remarquant (avec les notations de (6.6.7))
que l'on a un diagramme commutatif à homotopismes près
\[
\xymatrix@C=18pt@R=28pt{
  C^\bullet(\mathfrak U^{(1)},\mathscr P_\bullet^{(1)})
  \otimes_A
  C^\bullet(\mathfrak U^{(2)},\mathscr P_\bullet^{(2)})
  \ar[r]\ar[d]
  & C^\bullet(\mathfrak U,\mathscr Q_\bullet)\ar[d]
  \\
  C^\bullet(\mathfrak V^{(1)},\mathscr P_\bullet^{(1)})
  \otimes_A
  C^\bullet(\mathfrak V^{(2)},\mathscr P_\bullet^{(2)})
  \ar[r]
  & C^\bullet(\mathfrak V,\mathscr Q_\bullet)
}
\]
où les flèches horizontales sont les homomorphismes définis ci-dessus. Enfin,
il faut passer au cas général par recollement, ce qui se fait sans difficulté
comme dans (6.7.1) et (6.7.2); nous laissons les détails au lecteur.
\end{proof}

\begin{proposition}[6.7.7]
Supposons que $\mathscr P_\bullet^{(1)}$ et
$\mathscr P_\bullet^{(2)}$ soient limités inférieurement, et que tous les
Modules $\mathscr H^{-n}(f_1,\mathscr P_\bullet^{(1)})$ ou tous les Modules
$\mathscr H^{-n}(f_2,\mathscr P_\bullet^{(2)})$ soient $S$-plats. On a alors un
isomorphisme canonique de bi-$\partial$-foncteurs ($n$ parcourant $\mathbf Z$)
\[
  \mathfrak{Tor}_n^S
  (f_1,f_2;\mathscr P_\bullet^{(1)},\mathscr P_\bullet^{(2)})
  \xrightarrow{\sim}
  \bigoplus_{q_1+q_2=n}
  \mathscr H^{-q_1}(f_1,\mathscr P_\bullet^{(1)})
  \otimes_S
  \mathscr H^{-q_2}(f_2,\mathscr P_\bullet^{(2)}).
  \tag{6.7.7.1}\label{III.6.7.7.1.fr}
\]
\end{proposition}

\begin{proof}
En effet, vu (6.5.8), la suite spectrale $(a)$ de (6.7.3) est dégénérée, et
la proposition résulte aussitôt de ($0_{\mathrm I}$, 11.1.6), cette suite
étant birégulière (6.7.3).
\end{proof}

\begin{theorem}[6.7.8]
Supposons que : $1^\circ$ les complexes $\mathscr P_\bullet^{(1)}$ et
$\mathscr P_\bullet^{(2)}$ soient limités inférieurement; $2^\circ$ le
complexe $\mathscr P_\bullet^{(1)}$ ou le complexe
$\mathscr P_\bullet^{(2)}$ soit formé de Modules $S$-plats; $3^\circ$ tous les
\oldpage[III]{162}
Modules $\mathscr H^{-n}(f_1,\mathscr P_\bullet^{(1)})$ ou tous les Modules
$\mathscr H^{-n}(f_2,\mathscr P_\bullet^{(2)})$ soient $S$-plats. On a alors un
isomorphisme canonique de bi-$\partial$-foncteurs ($n$ parcourant $\mathbf Z$)
\[
  \mathscr H^n\bigl(f_1\times_S f_2,
    \mathscr P_\bullet^{(1)}\otimes_S\mathscr P_\bullet^{(2)}\bigr)
  \xrightarrow{\sim}
  \bigoplus_{n_1+n_2=n}
  \mathscr H^{n_1}(f_1,\mathscr P_\bullet^{(1)})
  \otimes_S
  \mathscr H^{n_2}(f_2,\mathscr P_\bullet^{(2)}).
  \tag{6.7.8.1}\label{III.6.7.8.1.fr}
\]
(« formule de Künneth »).
\end{theorem}

\begin{proof}
Cela résulte de (6.7.6) et (6.7.7).

Lorsque $S$, $Y^{(1)}$ et $Y^{(2)}$ sont affines, l'isomorphisme réciproque de
(6.7.8.1) se déduit (avec les notations de (6.7.6)) de l'homomorphisme de
bicomplexes
\[
  C^\bullet(\mathfrak U^{(1)},\mathscr P_\bullet^{(1)})
  \otimes_A
  C^\bullet(\mathfrak U^{(2)},\mathscr P_\bullet^{(2)})
  \longrightarrow
  C^\bullet(\mathfrak U,\mathscr Q_\bullet)
\]
par le procédé défini dans (G, I, 2.7), comme il résulte de (G, I, 5.5).
\end{proof}

\begin{proposition}[6.7.9]
Supposons vérifiées les trois conditions suivantes : $1^\circ$ $S$,
$Y^{(1)}$ et $Y^{(2)}$ sont localement noethériens, $f_1$ et $f_2$ sont
propres, $Y^{(1)}$ ou $Y^{(2)}$ de type fini sur $S$; $2^\circ$
$\mathscr P_\bullet^{(1)}$ et $\mathscr P_\bullet^{(2)}$ sont limités
inférieurement; $3^\circ$ pour tout $n\in\mathbf Z$,
$\mathscr H_n(\mathscr P_\bullet^{(i)})$ est un Module cohérent ($i=1,2$).
Dans ces conditions,
$\mathfrak{Tor}_n^S
(f_1,f_2;\mathscr P_\bullet^{(1)},\mathscr P_\bullet^{(2)})$ est un
$\mathscr O_Y$-Module cohérent (avec
$Y=Y^{(1)}\times_S Y^{(2)}$).
\end{proposition}

\begin{proof}
Il résulte de (6.5.13) que les hypertor locaux
$\mathscr{T}\!or_n^S(\mathscr P_\bullet^{(1)},\mathscr P_\bullet^{(2)})$ sont
des $\mathscr O_X$-Modules cohérents
($X=X^{(1)}\times_S X^{(2)}$ étant localement noethérien, car un des
$X^{(i)}$ est par hypothèse de type fini sur $S$ (I, 6.3.4 et 6.3.8)). Comme
$Y$ est localement noethérien et $f_1\times_S f_2$ propre (II, 5.4.2), il
résulte de (6.2.5) que les termes ${}^{(b)}\mathscr E_{pq}^2$ de (6.7.3) sont
des $\mathscr O_Y$-Modules cohérents. Comme toutes les suites spectrales de
(6.7.3) sont birégulières en vertu de l'hypothèse $2^\circ$, on conclut par
($0_{\mathrm I}$, 11.1.8).
\end{proof}

\begin{env}[6.7.10]
Soient maintenant $Y'^{(i)}$ deux $S$-préschémas,
$v_i:Y'^{(i)}\to Y^{(i)}$ deux $S$-morphismes ($i=1,2$),
$v=v_1\times_S v_2$ leur produit, qui est un $S$-morphisme $Y'\to Y$, où
l'on pose $Y'=Y'^{(1)}\times_S Y'^{(2)}$. Considérons d'autre part, pour
$i=1,2$, un $S$-préschéma $X'^{(i)}$, et deux $S$-morphismes
$u_i:X'^{(i)}\to X^{(i)}$, $f'_i:X'^{(i)}\to Y'^{(i)}$ de sorte que les
diagrammes
\[
\xymatrix@C=36pt@R=28pt{
  X'^{(i)}\ar[r]^{u_i}\ar[d]_{f'_i}
  & X^{(i)}\ar[d]^{f_i}
  \\
  Y'^{(i)}\ar[r]_{v_i}
  & Y^{(i)}
}
\tag{6.7.10.1}\label{III.6.7.10.1.fr}
\]
soient commutatifs, les morphismes $f'_i$ étant \emph{séparés et
quasi-compacts}. On a alors des $\mathscr O_{Y'}$-homomorphismes canoniques de
foncteurs spectraux
\[
  v^*\bigl({}^{(t)}\mathscr E
  (f_1,f_2;\mathscr P_\bullet^{(1)},\mathscr P_\bullet^{(2)})\bigr)
  \longrightarrow
  {}^{(t)}\mathscr E
  (f'_1,f'_2;u_1^*(\mathscr P_\bullet^{(1)}),
                 u_2^*(\mathscr P_\bullet^{(2)}))
  \tag{6.7.10.2}\label{III.6.7.10.2.fr}
\]
pour $t=a$, $a'$, $b$, $b'$, $c$, $d$. Pour les définir, supposons d'abord
$S=\operatorname{Spec}(A)$,
$Y^{(i)}=\operatorname{Spec}(B_i)$,
$Y'^{(i)}=\operatorname{Spec}(B'_i)$ affines; les $X^{(i)}$ et $X'^{(i)}$
sont alors des schémas quasi-compacts. Pour calculer les suites spectrales
${}^{(t)}\mathscr E
(f_1,f_2;\mathscr P_\bullet^{(1)},\mathscr P_\bullet^{(2)})$, nous
considérons comme dans (6.6.1) des recouvrements finis
$\mathfrak U^{(i)}$ par des ouverts affines de $X^{(i)}$ ($i=1,2$); pour
calculer
${}^{(t)}\mathscr E
(f'_1,f'_2;u_1^*(\mathscr P_\bullet^{(1)}),
u_2^*(\mathscr P_\bullet^{(2)}))$, nous considérons des recouvrements finis
$\mathfrak U'^{(i)}$

\oldpage[III]{163}
de $X'^{(i)}$ par des ouverts affines, \emph{plus fins} respectivement que les
recouvrements $u_i^{-1}(\mathfrak U^{(i)})$ ($i=1,2$). Il est clair que le
bicomplexe
$C^\bullet(\mathfrak U^{(i)},\mathscr P_\bullet^{(i)})
=L_{\bullet\bullet}^{(i)}$ peut être considéré canoniquement comme un
sous-bicomplexe de
$C^\bullet(u_i^{-1}(\mathfrak U^{(i)}),u_i^*(\mathscr P_\bullet^{(i)}))$
($0_{\mathrm I}$, 4.4.3.2); en outre, en choisissant une application
simpliciale (G, II, 5.7) de $\mathfrak U'^{(i)}$ dans
$u_i^{-1}(\mathfrak U^{(i)})$, on définit un homomorphisme de bicomplexes
\[
  C^\bullet(u_i^{-1}(\mathfrak U^{(i)}),u_i^*(\mathscr P_\bullet^{(i)}))
  \longrightarrow
  C^\bullet(\mathfrak U'^{(i)},u_i^*(\mathscr P_\bullet^{(i)}))
\]
d'où, par composition, un homomorphisme de bicomplexes
\[
  L_{\bullet\bullet}^{(i)}
  \longrightarrow
  L_{\bullet\bullet}^{\prime(i)}
  =C^\bullet(\mathfrak U'^{(i)},u_i^*(\mathscr P_\bullet^{(i)})).
\]
En outre, cet homomorphisme est remplacé par un homomorphisme homotope quand
on change d'application simpliciale (G, II, 5.7.1); on a ainsi un
homomorphisme bien défini de foncteurs spectraux :
\[
  {}^{(t)}\mathscr E
  (\mathfrak U^{(1)},\mathfrak U^{(2)};
   \mathscr P_\bullet^{(1)},\mathscr P_\bullet^{(2)})
  \longrightarrow
  {}^{(t)}\mathscr E
  (\mathfrak U'^{(1)},\mathfrak U'^{(2)};
   u_1^*(\mathscr P_\bullet^{(1)}),u_2^*(\mathscr P_\bullet^{(2)})).
  \tag{6.7.10.3}\label{III.6.7.10.3.fr}
\]
On vérifie aussitôt que si $\mathfrak V^{(i)}$ est un recouvrement affine
fini de $X^{(i)}$ plus fin que $\mathfrak U^{(i)}$,
$\mathfrak V'^{(i)}$ un recouvrement affine fini de $X'^{(i)}$, plus fin que
$u_i^{-1}(\mathfrak V^{(i)})$ et que $\mathfrak U'^{(i)}$, le diagramme
\[
\xymatrix@C=34pt@R=30pt{
  C^\bullet(\mathfrak U^{(i)},\mathscr P_\bullet^{(i)})
  \ar[r]\ar[d]
  & C^\bullet(\mathfrak V^{(i)},\mathscr P_\bullet^{(i)})\ar[d]
  \\
  C^\bullet(\mathfrak U'^{(i)},u_i^*(\mathscr P_\bullet^{(i)}))
  \ar[r]
  & C^\bullet(\mathfrak V'^{(i)},u_i^*(\mathscr P_\bullet^{(i)}))
}
\]
est commutatif, ce qui implique que l'homomorphisme (6.7.10.3) ne dépend pas
essentiellement des recouvrements $\mathfrak U^{(i)}$ et
$\mathfrak U'^{(i)}$ considérés. On a donc en fait défini un homomorphisme de
$A$-modules
\[
  {}^{(t)}E
  (X^{(1)},X^{(2)};\mathscr P_\bullet^{(1)},\mathscr P_\bullet^{(2)})
  \longrightarrow
  {}^{(t)}E
  (X'^{(1)},X'^{(2)};
   u_1^*(\mathscr P_\bullet^{(1)}),u_2^*(\mathscr P_\bullet^{(2)})).
  \tag{6.7.10.4}\label{III.6.7.10.4.fr}
\]
mais il est clair par définition des $u_i^*(\mathscr P_\bullet^{(i)})$ et en
vertu de la commutativité de (6.7.10.1) que cet homomorphisme est aussi un
homomorphisme de $(B_1\otimes_A B_2)$-modules; comme le second membre de
(6.7.10.4) est formé de $(B'_1\otimes_A B'_2)$-modules, on déduit
canoniquement de (6.7.10.4) un homomorphisme de
$(B'_1\otimes_A B'_2)$-modules
\[
  {}^{(t)}E
  (X^{(1)},X^{(2)};\mathscr P_\bullet^{(1)},\mathscr P_\bullet^{(2)})
  \otimes_{B_1\otimes_A B_2}(B'_1\otimes_A B'_2)
  \longrightarrow
  {}^{(t)}E
  (X'^{(1)},X'^{(2)};
   u_1^*(\mathscr P_\bullet^{(1)}),u_2^*(\mathscr P_\bullet^{(2)}))
  \tag{6.7.10.5}\label{III.6.7.10.5.fr}
\]
ce qui, compte tenu de (I, 1.6.5) n'est autre que l'homomorphisme cherché
(6.7.10.2) dans le cas particulier considéré.

Il reste à passer au cas général en suivant les recollements de (6.7.1) et
(6.7.2); le second passage est immédiat; en ce qui concerne le premier, on
considère comme dans (6.7.1) des éléments $g_i\in B_i$, et leurs images
$g'_i\in B'_i$, le produit tensoriel $g=g_1\otimes g_2$ dans
$B=B_1\otimes_A B_2$ et son image $g'=g'_1\otimes g'_2$ dans
$B'=B'_1\otimes_A B'_2$, et tout revient à utiliser

\oldpage[III]{164}
l'isomorphisme canonique
\[
  (M\otimes_B B')_{g'}
  \xrightarrow{\sim}
  M_g\otimes_{B_g}B'_{g'}
\]
($0_{\mathrm I}$, 1.5.4); nous laissons les détails au lecteur.
\end{env}

\begin{env}[6.7.11]
La théorie des hypertor globaux, développée ci-dessus pour deux
$S$-morphismes $X^{(i)}\to Y^{(i)}$ et deux complexes
$\mathscr P_\bullet^{(i)}$ de Modules quasi-cohérents limités inférieurement,
s'étend aussitôt au cas général suivant : on a un préschéma $S$, une famille
finie de $S$-préschémas $Y^{(i)}$ ($i\in I$), une famille finie de
$S$-morphismes séparés et quasi-compacts
$f_i:X^{(i)}\to Y^{(i)}$, et pour chaque $i$ un complexe de
$\mathscr O_{X^{(i)}}$-Modules quasi-cohérents
$\mathscr P_\bullet^{(i)}$ limités inférieurement. Si $Y$ est le produit des
$S$-préschémas $Y^{(i)}$, on définit alors pour chaque entier
$n\in\mathbf Z$, un $\mathscr O_Y$-Module quasi-cohérent
$\mathfrak{Tor}_n^S((f_i)_{i\in I};
(\mathscr P_\bullet^{(i)})_{i\in I})$,
ces Modules formant un $\partial$-foncteur covariant en chacun des complexes
$\mathscr P_\bullet^{(i)}$; en outre, ce foncteur est l'aboutissement commun
de six foncteurs spectraux
${}^{(t)}\mathscr E((f_i)_{i\in I};(\mathscr P_\bullet^{(i)})_{i\in I})$.
Nous laissons au lecteur le soin de répéter pour ce cas général les
définitions et les raisonnements faits ci-dessus pour $I=\{1,2\}$. Notons
simplement que lorsque $I$ se réduit à un seul élément, on retrouve
l'hypercohomologie $\mathscr H^\bullet(f,\mathscr P_\bullet)$ définie dans
(6.2.7) (comme on l'a déjà observé dans (6.7.4)). Lorsque $I$ est l'intervalle
$1\leq i\leq m$ de $\mathbf N$, nous écrirons
\[
  \mathfrak{Tor}_n^S
  (f_1,\ldots,f_m;
   \mathscr P_\bullet^{(1)},\ldots,\mathscr P_\bullet^{(m)})
  \quad\text{pour}\quad
  \mathfrak{Tor}_n^S
  ((f_i)_{i\in I};(\mathscr P_\bullet^{(i)})_{i\in I}).
\]
\end{env}

\begin{proposition}[6.7.12]
Les notations étant celles de (6.7.11), soit $J$ une partie de $I$ telle que,
pour $i\in I-J$, on ait $X^{(i)}=Y^{(i)}=S$, $f_i$ étant réduit à l'identité,
et $\mathscr P_\bullet^{(i)}$ égal au complexe réduit au terme de degré $0$
égal à $\mathscr O_S$. Il y a alors un isomorphisme canonique de
$\partial$-foncteurs
\[
  \mathfrak{Tor}_\bullet^S
  ((f_i)_{i\in I};(\mathscr P_\bullet^{(i)})_{i\in I})
  \xrightarrow{\sim}
  \mathfrak{Tor}_\bullet^S
  ((f_i)_{i\in J};(\mathscr P_\bullet^{(i)})_{i\in J}).
  \tag{6.7.12.1}\label{III.6.7.12.1.fr}
\]
\end{proposition}

\begin{proof}
On peut se borner à définir cet isomorphisme lorsque $S$ et les $Y^{(i)}$
sont affines, le recollement se faisant comme d'ordinaire. Pour $i\in I-J$,
on peut prendre le recouvrement $\mathfrak U^{(i)}$ formé du seul ensemble
$S$, et alors
$L_{\bullet\bullet}^{(i)}=C^\bullet(\mathfrak U^{(i)},
\mathscr P_\bullet^{(i)})$ est réduit à son seul terme de degrés $(0,0)$,
égal à $\Gamma(S,\mathscr O_S)=A(S)$; l'isomorphisme (6.7.12.1) est alors
évident.
\end{proof}

\begin{remark}[6.7.13]
Les notations étant celles de (6.7.3), considérons le $S$-isomorphisme
canonique
$Y^{(1)}\times_S Y^{(2)}\to Y^{(2)}\times_S Y^{(1)}$ (I, 3.3.5); alors
l'image par cet isomorphisme de
$\mathfrak{Tor}_\bullet^S(f_1,f_2;
\mathscr P_\bullet^{(1)},\mathscr P_\bullet^{(2)})$ est
$\mathfrak{Tor}_\bullet^S(f_2,f_1;
\mathscr P_\bullet^{(2)},\mathscr P_\bullet^{(1)})$; la question étant
locale, on est ramené au cas envisagé dans (6.6.2), et si on désigne par
$M_{\bullet\bullet\bullet}^{(i)}$ une résolution projective de
Cartan-Eilenberg de $L_{\bullet\bullet}^{(i)}$ ($i=1,2$), l'isomorphisme
considéré transforme
$M_{\bullet\bullet\bullet}^{(1)}\otimes_A
M_{\bullet\bullet\bullet}^{(2)}$ en
$M_{\bullet\bullet\bullet}^{(2)}\otimes_A
M_{\bullet\bullet\bullet}^{(1)}$, d'où notre assertion en considérant
l'homologie des complexes simples associés à ces tricomplexes.
\end{remark}

\subsection{Les suites spectrales d'associativité des hypertor globaux.}

\begin{env}[6.8.1]
Les hypothèses et notations étant celles de (6.7.11) (et en particulier les
$\mathscr P_\bullet^{(i)}$ étant supposés limités inférieurement, supposons
donnée une partition $(I_j)_{j\in J}$ de l'ensemble d'indices $I$; nous nous
proposons de donner une relation d'« associativité » entre les hypertor
$\mathfrak{Tor}_n^S((f_i)_{i\in I};(\mathscr P_\bullet^{(i)})_{i\in I})$ et
chacun des hypertor « partiels »
\[
  \mathcal T_{\bullet j}
  =\mathfrak{Tor}_\bullet^S
  ((f_i)_{i\in I_j};(\mathscr P_\bullet^{(i)})_{i\in I_j}).
\]

\oldpage[III]{165}
Pour simplifier l'écriture, nous nous bornerons au cas où $I$ est
l'intervalle $1\leq i\leq m$, et où la partition $(I_j)$ se compose des deux
intervalles $\{1,2,\ldots,r\}$ et $\{r+1,\ldots,m\}$.
\end{env}

\begin{proposition}[6.8.2]
Il existe un foncteur spectral canonique birégulier (dit « foncteur spectral
d'associativité ») noté
\[
  {}^{(e)}\mathscr E^S
  (f_1,\ldots,f_m;
   \mathscr P_\bullet^{(1)},\ldots,\mathscr P_\bullet^{(m)})
\]
(ou simplement ${}^{(e)}\mathscr E(f_1,\ldots,f_m;
\mathscr P_\bullet^{(1)},\ldots,\mathscr P_\bullet^{(m)})$) dont
l'aboutissement est
$\mathfrak{Tor}_\bullet^S(f_1,\ldots,f_m;
\mathscr P_\bullet^{(1)},\ldots,\mathscr P_\bullet^{(m)})$, et dont le terme
$E_2$ est donné par
\[
\begin{split}
  {}^{(e)}\mathscr E_{pq}^2
  =\bigoplus_{q_1+q_2=q}
  \mathscr{T}\!or_p^S\Bigl(
  &\mathfrak{Tor}_{q_1}^S
  (f_1,\ldots,f_r;
   \mathscr P_\bullet^{(1)},\ldots,\mathscr P_\bullet^{(r)}),\\
  &\mathfrak{Tor}_{q_2}^S
  (f_{r+1},\ldots,f_m;
   \mathscr P_\bullet^{(r+1)},\ldots,\mathscr P_\bullet^{(m)})
  \Bigr).
\end{split}
\]
\end{proposition}

\begin{proof}
Dans cet énoncé, on a identifié canoniquement $Y$ au produit
$Z^{(1)}\times_S Z^{(2)}$, où
$Z^{(1)}=Y^{(1)}\times_S Y^{(2)}\times_S\cdots\times_S Y^{(r)}$, et
$Z^{(2)}=Y^{(r+1)}\times_S\cdots\times_S Y^{(m)}$. Nous nous bornerons au
cas où $S$ et les $Y^{(i)}$ sont affines; on passe de ce cas particulier au
cas général par les méthodes développées dans (6.7.1) et (6.7.2), et nous
laissons les détails du raisonnement (sans difficulté) au lecteur. Nous
démontrerons donc le corollaire suivant.
\end{proof}

\begin{corollary}[6.8.3]
Soient $A$ un anneau, $S=\operatorname{Spec}(A)$, $X^{(i)}$
($1\leq i\leq m$) des $S$-schémas quasi-compacts et, pour chaque $i$, soit
$\mathscr P_\bullet^{(i)}$ un complexe de
$\mathscr O_{X^{(i)}}$-Modules quasi-cohérents limité inférieurement. Il
existe un foncteur spectral canonique birégulier ayant pour aboutissement
\[
  \mathbf{Tor}_\bullet^S
  (X^{(1)},\ldots,X^{(m)};
   \mathscr P_\bullet^{(1)},\ldots,\mathscr P_\bullet^{(m)})
\]
et dont le terme $E_2$ est donné par
\[
\begin{split}
  {}^{(e)}E_{pq}^2
  =\bigoplus_{q_1+q_2=q}
  \operatorname{Tor}_p^A\Bigl(
  &\mathbf{Tor}_{q_1}^S
  (X^{(1)},\ldots,X^{(r)};
   \mathscr P_\bullet^{(1)},\ldots,\mathscr P_\bullet^{(r)}),\\
  &\mathbf{Tor}_{q_2}^S
  (X^{(r+1)},\ldots,X^{(m)};
   \mathscr P_\bullet^{(r+1)},\ldots,\mathscr P_\bullet^{(m)})
  \Bigr).
\end{split}
\]
\end{corollary}

\begin{proof}
Suivant la définition donnée en (6.6.2), le calcul de l'hypertor considéré se
fait en prenant pour chaque $i$ un recouvrement ouvert affine fini
$\mathfrak U^{(i)}$ de $X^{(i)}$, en considérant les bicomplexes
$L_{\bullet\bullet}^{(i)}=C^\bullet(\mathfrak U^{(i)},
\mathscr P_\bullet^{(i)})$, une résolution projective
$M_{\bullet\bullet\bullet}^{(i)}$ de Cartan-Eilenberg de chacun de ces
bicomplexes (au sens de (0, 11.7.1)), le produit tensoriel
\[
  M_{\bullet\bullet\bullet}
  =\bigotimes_{i=1}^m M_{\bullet\bullet\bullet}^{(i)}
\]
de ces tricomplexes, et en prenant l'homologie de
$M_{\bullet\bullet\bullet}$. Considérons $M_{\bullet\bullet\bullet}$ comme
un complexe simple $N_\bullet$, produit tensoriel des deux complexes simples
\[
  N'_\bullet=\bigotimes_{i=1}^r M_{\bullet\bullet\bullet}^{(i)},
  \qquad
  N''_\bullet=\bigotimes_{i=r+1}^m M_{\bullet\bullet\bullet}^{(i)},
\]
où $N'_\bullet$ et $N''_\bullet$ sont gradués par la somme des degrés totaux
des $M_{\bullet\bullet\bullet}^{(i)}$. En outre, les $A$-modules des
complexes $N'_\bullet$ et $N''_\bullet$ sont projectifs, donc il résulte de
(6.5.9) que l'on a
\[
  H_\bullet(M_{\bullet\bullet\bullet})
  =\mathbf{Tor}_\bullet^A(N'_\bullet,N''_\bullet);
\]
la suite spectrale cherchée n'est autre alors que la suite (6.3.2.2)
appliquée aux complexes $N'_\bullet$ et $N''_\bullet$, compte tenu de
l'interprétation des modules d'homologie de ces complexes qui résulte de ce
qui précède (quand on applique les remarques du début à chacun des produits
partiels
$X^{(1)}\times_S\cdots\times_S X^{(r)}$ et
$X^{(r+1)}\times_S\cdots\times_S X^{(m)}$). Enfin, les propriétés de
régularité résultent de (6.3.2) et du fait que, les
$\mathscr P_\bullet^{(i)}$ étant limités inférieurement, il en est de même
des $M_{\bullet\bullet\bullet}^{(i)}$; par suite, $N'_\bullet$ et
$N''_\bullet$ sont limités inférieurement.
\end{proof}

\oldpage[III]{166}
\subsection{Les suites spectrales de changement de base dans les hypertor
globaux.}

\begin{env}[6.9.1]
Les hypothèses et notations étant toujours celles de (6.7.11) (et en
particulier les $\mathscr P_\bullet^{(i)}$ étant supposés limités
inférieurement), considérons un morphisme $g:S'\to S$ de préschémas, et
posons
$Y'^{(i)}=Y_{(S')}^{(i)}$, $X'^{(i)}=X_{(S')}^{(i)}$ et
$\mathscr P_\bullet'^{(i)}=
\mathscr P_\bullet^{(i)}\otimes_{\mathscr O_S}\mathscr O_{S'}$,
$\mathscr P_\bullet'^{(i)}$ étant donc un complexe de
$\mathscr O_{X'^{(i)}}$-Modules quasi-cohérents; soit
$f'_i=(f_i)_{(S')}:X'^{(i)}\to Y'^{(i)}$, qui est un morphisme séparé et
quasi-compact (I, 5.5.1 et 6.6.4). Nous nous proposons d'étudier les relations
entre les $\mathscr O_{Y'}$-Modules quasi-cohérents
$\mathfrak{Tor}_n^{S'}((f'_i)_{i\in I};
(\mathscr P_\bullet'^{(i)})_{i\in I})$ et
$\mathfrak{Tor}_n^S((f_i)_{i\in I};
(\mathscr P_\bullet^{(i)})_{i\in I})
\otimes_{\mathscr O_S}\mathscr O_{S'}$, où
$Y'=Y\times_S S'=Y_{(S')}$. Un cas particulièrement simple est le suivant,
qui se réduit à (1.4.15) lorsque $I$ se réduit à un seul élément et
$\mathscr P_\bullet$ à un seul module :
\end{env}

\begin{proposition}[6.9.2]
Si le morphisme $g:S'\to S$ est plat, on a un isomorphisme canonique de
$\partial$-foncteurs (en les $\mathscr P_\bullet^{(i)}$) :
\[
  \mathfrak{Tor}_\bullet^S
  ((f_i)_{i\in I};(\mathscr P_\bullet^{(i)})_{i\in I})
  \otimes_{\mathscr O_S}\mathscr O_{S'}
  \xrightarrow{\sim}
  \mathfrak{Tor}_\bullet^{S'}
  ((f'_i)_{i\in I};(\mathscr P_\bullet'^{(i)})_{i\in I}).
  \tag{6.9.2.1}\label{III.6.9.2.1.fr}
\]
\end{proposition}

\begin{proof}
On peut encore se borner au cas où $S$, $S'$ et les $Y^{(i)}$ sont affines,
le recollement se faisant suivant les méthodes de (6.7.1) et (6.7.2). Soient
$S=\operatorname{Spec}(A)$, $S'=\operatorname{Spec}(A')$, et prenons pour
chaque $i$ un recouvrement ouvert affine $\mathfrak U^{(i)}$ de $X^{(i)}$;
si $u_i:X'^{(i)}\to X^{(i)}$ est la projection canonique,
$u_i^{-1}(\mathfrak U^{(i)})$ est un recouvrement ouvert affine de
$X'^{(i)}$ (II, 1.5.5), que nous noterons $\mathfrak U'^{(i)}$; il est clair
alors que
\[
  C^\bullet(\mathfrak U'^{(i)},\mathscr P_\bullet'^{(i)})
  =C^\bullet(\mathfrak U^{(i)},\mathscr P_\bullet^{(i)})\otimes_A A',
\]
et l'existence de l'isomorphisme (6.9.2.1) est immédiate, car si
$M_{\bullet\bullet\bullet}^{(i)}$ est une résolution projective de
Cartan-Eilenberg de
$L_{\bullet\bullet}^{(i)}=C^\bullet(\mathfrak U^{(i)},
\mathscr P_\bullet^{(i)})$ au sens de (0, 11.7.1), formée de $A$-modules,
$M_{\bullet\bullet\bullet}^{(i)}\otimes_A A'$ est une résolution projective
de Cartan-Eilenberg (au même sens) de
$L_{\bullet\bullet}^{(i)}\otimes_A A'$ formée de $A'$-modules, en vertu de
l'hypothèse que $A'$ est un $A$-module plat; cette même hypothèse montre en
outre que
\[
  H_\bullet\left(
    \bigotimes_{i=1}^m
    (M_{\bullet\bullet\bullet}^{(i)}\otimes_A A')
  \right)
  =H_\bullet\left(
    \bigotimes_{i=1}^m M_{\bullet\bullet\bullet}^{(i)}
  \right)\otimes_A A'.
\]
\end{proof}

\begin{env}[6.9.2.2]
On notera que lorsque $I$ est réduit au seul élément $1$, la formule
(6.9.2.1) se déduit directement de (6.7.7), appliqué en prenant
$Y_2=X_2=S'$, $f_2=1_{S'}$, et le complexe $\mathscr P_\bullet^{(2)}$
réduit à son terme de degré $0$, égal à $\mathscr O_{S'}$; on sait alors que
l'hypercohomologie $\mathscr H^n(1_{S'},\mathscr O_{S'})$ est nulle pour tout
$n\ne0$ et se réduit à $\mathscr O_{S'}$ pour $n=0$ (6.2.1).
\end{env}

Dans le cas général, nous allons introduire à la place de $\mathscr O_{S'}$
un complexe $\mathscr Q'_\bullet$ de $\mathscr O_{S'}$-Modules
quasi-cohérents limités inférieurement, de sorte que si, pour simplifier, on
prend $I=\{1,2,\ldots,m\}$, on peut considérer le $\partial$-foncteur
\[
  \mathfrak{Tor}_\bullet^S
  (f_1,\ldots,f_m,1_{S'};
   \mathscr P_\bullet^{(1)},\ldots,\mathscr P_\bullet^{(m)},
   \mathscr Q'_\bullet).
\]

\begin{proposition}[6.9.3]
Il existe trois foncteurs spectraux canoniques biréguliers notés
${}^{(t)}\mathscr E(f_1,\ldots,f_m;
\mathscr P_\bullet^{(1)},\ldots,\mathscr P_\bullet^{(m)},
\mathscr Q'_\bullet)$ (avec $t=e$, $f$ ou $f'$) ayant pour aboutissement
commun
$\mathfrak{Tor}_\bullet^S(f_1,\ldots,f_m,1_{S'};
\mathscr P_\bullet^{(1)},\ldots,\mathscr P_\bullet^{(m)},
\mathscr Q'_\bullet)$ et dont les termes $E_2$ sont respectivement

\oldpage[III]{167}
\[
\begin{split}
  {}^{(e)}\mathscr E_{pq}^2
  =\bigoplus_{q'+q''=q}
  \mathscr{T}\!or_p^S\Bigl(
  &\mathfrak{Tor}_{q'}^S
  (f_1,\ldots,f_m;
   \mathscr P_\bullet^{(1)},\ldots,\mathscr P_\bullet^{(m)}),\\
  &\mathscr H_{q''}(\mathscr Q'_\bullet)
  \Bigr),
\end{split}
\]
\[
\begin{split}
  {}^{(f)}\mathscr E_{pq}^2
  =\bigoplus_{q_1+q_2+\cdots+q_{m+1}=q}
  \mathscr{T}\!or_p^{S'}\Bigl(
  &\mathfrak{Tor}_{q_1}^S
  (f_1,1_{S'};\mathscr P_\bullet^{(1)},\mathscr O_{S'}),\ldots,\\
  &\mathfrak{Tor}_{q_m}^S
  (f_m,1_{S'};\mathscr P_\bullet^{(m)},\mathscr O_{S'}),
  \mathscr H_{q_{m+1}}(\mathscr Q'_\bullet)
  \Bigr),
\end{split}
\]
\[
\begin{split}
  {}^{(f')}\mathscr E_{pq}^2
  =\bigoplus_{q_1+\cdots+q_m=q}
  \mathfrak{Tor}_p^{S'}\Bigl(
  &f'_1,\ldots,f'_m,1_{S'};\\
  &\mathscr{T}\!or_{q_1}^S
  (\mathscr P_\bullet^{(1)},\mathscr O_{S'}),\ldots,
  \mathscr{T}\!or_{q_m}^S
  (\mathscr P_\bullet^{(m)},\mathscr O_{S'}),
  \mathscr Q'_\bullet
  \Bigr).
\end{split}
\]
\end{proposition}

\begin{proof}
La suite $(e)$ n'est autre que la suite d'associativité de (6.8.2) pour
$r=m$. Pour définir les deux autres suites spectrales, on va encore se limiter
au cas où $S$, $S'$ et les $Y^{(i)}$ sont affines, le passage au cas général
se faisant par les méthodes de (6.7.1) et (6.7.2) et étant laissé au lecteur.
Nous démontrerons donc le
\end{proof}

\begin{corollary}[6.9.4]
Soient $A$ un anneau, $A'$ une $A$-algèbre,
$S=\operatorname{Spec}(A)$, $S'=\operatorname{Spec}(A')$,
$X^{(i)}$ ($1\leq i\leq m$) des $S$-schémas quasi-compacts et, pour chaque
$i$, soit $X'^{(i)}=X_{(S')}^{(i)}$, qui est un $S'$-schéma quasi-compact.
Pour chaque $i$, soit $\mathscr P_\bullet^{(i)}$ un complexe de
$\mathscr O_{X^{(i)}}$-Modules quasi-cohérents; soit enfin $Q'_\bullet$ un
complexe de $A'$-modules, ces complexes étant limités inférieurement. Il
existe trois foncteurs spectraux biréguliers en les
$\mathscr P_\bullet^{(i)}$ et en $Q'_\bullet$, ayant pour aboutissement commun
\[
  \mathbf{Tor}_\bullet^S
  (X^{(1)},\ldots,X^{(m)},S';
   \mathscr P_\bullet^{(1)},\ldots,\mathscr P_\bullet^{(m)},
   \widetilde{Q'_\bullet})
\]
et dont les termes $E_2$ sont respectivement
\[
\begin{split}
  {}^{(e)}E_{pq}^2
  =\bigoplus_{q'+q''=q}
  \operatorname{Tor}_p^A\Bigl(
  &\mathbf{Tor}_{q'}^S
  (X^{(1)},\ldots,X^{(m)};
   \mathscr P_\bullet^{(1)},\ldots,\mathscr P_\bullet^{(m)}),\\
  &H_{q''}(Q'_\bullet)
  \Bigr),
\end{split}
\]
\[
\begin{split}
  {}^{(f)}E_{pq}^2
  =\bigoplus_{q_1+q_2+\cdots+q_{m+1}=q}
  \operatorname{Tor}_p^{A'}\Bigl(
  &\mathbf{Tor}_{q_1}^S
  (X^{(1)},S';\mathscr P_\bullet^{(1)},\mathscr O_{S'}),\ldots,\\
  &\mathbf{Tor}_{q_m}^S
  (X^{(m)},S';\mathscr P_\bullet^{(m)},\mathscr O_{S'}),
  H_{q_{m+1}}(Q'_\bullet)
  \Bigr),
\end{split}
\]
\[
\begin{split}
  {}^{(f')}E_{pq}^2
  =\bigoplus_{q_1+\cdots+q_m=q}
  \mathbf{Tor}_p^{S'}\Bigl(
  &X'^{(1)},\ldots,X'^{(m)},S';\\
  &\mathscr{T}\!or_{q_1}^S
  (\mathscr P_\bullet^{(1)},\mathscr O_{S'}),\ldots,
  \mathscr{T}\!or_{q_m}^S
  (\mathscr P_\bullet^{(m)},\mathscr O_{S'}),
  \widetilde{Q'_\bullet}
  \Bigr).
\end{split}
\]
\end{corollary}

\begin{proof}
Nous ne reviendrons pas sur le premier de ces foncteurs spectraux, qui a été
traité dans (6.8.3) et n'est inclus ici que pour mémoire. Pour définir les
autres, considérons pour chaque $i$ un recouvrement ouvert affine fini
$\mathfrak U^{(i)}$ de $X^{(i)}$ et, si
$u_i:X'^{(i)}\to X^{(i)}$ est la projection canonique, le recouvrement ouvert
affine fini correspondant
$\mathfrak U'^{(i)}=u_i^{-1}(\mathfrak U^{(i)})$.

En vertu de (6.6.6),
$\mathbf{Tor}_\bullet^S(X^{(1)},\ldots,X^{(m)},S';
\mathscr P_\bullet^{(1)},\ldots,\mathscr P_\bullet^{(m)},
\widetilde{Q'_\bullet})$ s'obtient en prenant pour $1\leq i\leq m$ une
résolution projective de Cartan-Eilenberg
$M_{\bullet\bullet\bullet}^{(i)}$ de
$L_{\bullet\bullet}^{(i)}=C^\bullet(\mathfrak U^{(i)},
\mathscr P_\bullet^{(i)})$ (au sens de (0, 11.7.1)), considérant le
tricomplexe
\[
  M_{\bullet\bullet\bullet}
  =M_{\bullet\bullet\bullet}^{(1)}\otimes_A
   M_{\bullet\bullet\bullet}^{(2)}\otimes_A\cdots\otimes_A
   M_{\bullet\bullet\bullet}^{(m)}\otimes_A Q'_\bullet
\]
(où $Q'_\bullet$ est considéré comme un tricomplexe dont les deux derniers
degrés se réduisent à $0$), et en en prenant l'homologie. Si l'on pose
$M_{\bullet\bullet\bullet}'^{(i)}=
M_{\bullet\bullet\bullet}^{(i)}\otimes_A A'$, on a (en se rappelant que
$Q'_\bullet$ est un complexe de $A'$-modules)
\[
  M_{\bullet\bullet\bullet}
  =M_{\bullet\bullet\bullet}'^{(1)}\otimes_{A'}
   M_{\bullet\bullet\bullet}'^{(2)}\otimes_{A'}\cdots\otimes_{A'}
   M_{\bullet\bullet\bullet}'^{(m)}\otimes_{A'} Q'_\bullet.
\]
Or, considérons chacun des complexes
$M_{\bullet\bullet\bullet}'^{(i)}$ comme un complexe simple (pour son degré
total) et notons que ce complexe est formé de $A'$-modules projectifs; il
résulte de (6.3.7) (étendu à un nombre quelconque de complexes) que
$H_\bullet(M_{\bullet\bullet\bullet})$ est aussi égal à
\[
  \operatorname{Tor}_\bullet^{A'}
  (M_{\bullet\bullet\bullet}'^{(1)},\ldots,
   M_{\bullet\bullet\bullet}'^{(m)},Q'_\bullet);
\]
c'est donc (6.5.15) l'aboutissement d'une suite spectrale ayant les propriétés
de régularité voulues (les trois degrés de
$M_{\bullet\bullet\bullet}^{(i)}$ étant limités inférieurement lorsque
$\mathscr P_\bullet^{(i)}$ est limité inférieurement) et dont le terme $E_2$
est donné par
\[
  E_{pq}^2
  =\bigoplus_{q_1+\cdots+q_{m+1}=q}
  \operatorname{Tor}_p^{A'}
  (H_{q_1}(M_{\bullet\bullet\bullet}'^{(1)}),\ldots,
   H_{q_m}(M_{\bullet\bullet\bullet}'^{(m)}),H_{q_{m+1}}(Q'_\bullet)).
\]

\oldpage[III]{168}
On a d'ailleurs
\[
  H_{q_i}(M_{\bullet\bullet\bullet}'^{(i)})
  =H_{q_i}(M_{\bullet\bullet\bullet}^{(i)}\otimes_A A')
  =\mathbf{Tor}_{q_i}^S
  (X^{(i)},S';\mathscr P_\bullet^{(i)},\mathscr O_{S'})
\]
en vertu de la définition des hypertor globaux, ce qui donne la suite $(f)$
cherchée. On peut d'autre part considérer
$M_{\bullet\bullet\bullet}'^{(i)}$ comme un bicomplexe dans lequel le premier
degré est la somme du premier et du second degré du tricomplexe
$M_{\bullet\bullet\bullet}^{(i)}$, le second degré étant le troisième degré
de ce tricomplexe; comme les modules formant les
$M_{\bullet\bullet\bullet}'^{(i)}$ sont des $A'$-modules projectifs, la théorie
générale de l'hyperhomologie montre que l'homologie du bicomplexe
\[
  M_{\bullet\bullet\bullet}'^{(1)}\otimes_{A'}
  M_{\bullet\bullet\bullet}'^{(2)}\otimes_{A'}\cdots\otimes_{A'}
  M_{\bullet\bullet\bullet}'^{(m)}\otimes_{A'} Q'_\bullet
\]
est canoniquement isomorphe à son hyperhomologie (0, 11.6.5); c'est donc
l'aboutissement d'une suite spectrale de terme $E_2$ égal à
\[
  E_{pq}^2
  =\bigoplus_{q_1+\cdots+q_{m+1}=q}
  \operatorname{Tor}_p^{A'}
  (H_{q_1}^{\mathrm{II}}(M_{\bullet\bullet\bullet}'^{(1)}),\ldots,
   H_{q_m}^{\mathrm{II}}(M_{\bullet\bullet\bullet}'^{(m)}),
   H_{q_{m+1}}^{\mathrm{II}}(Q'_\bullet)).
\]
Or, comme le second degré de $Q'_\bullet$ se réduit à $0$, on a
$H_n^{\mathrm{II}}(Q'_\bullet)=0$ pour $n\ne0$ et
$H_0^{\mathrm{II}}(Q'_\bullet)=Q'_\bullet$; la formule précédente s'écrit
aussi
\[
  E_{pq}^2
  =\bigoplus_{q_1+\cdots+q_m=q}
  \operatorname{Tor}_p^{A'}
  (H_{q_1}^{\mathrm{II}}(M_{\bullet\bullet\bullet}'^{(1)}),\ldots,
   H_{q_m}^{\mathrm{II}}(M_{\bullet\bullet\bullet}'^{(m)}),Q'_\bullet).
\]
En outre, on a
\[
  H_{q_i}^{\mathrm{II}}(M_{\bullet\bullet\bullet}'^{(i)})
  =H_{q_i}^{\mathrm{II}}(M_{\bullet\bullet\bullet}^{(i)}\otimes_A A')
  =\operatorname{Tor}_{q_i}^A(L_{\bullet\bullet}^{(i)},A')
\]
en vertu de (6.3.4); mais
$L_{-j,k}^{(i)}=C^j(\mathfrak U^{(i)},\mathscr P_k^{(i)})$, somme directe des
$\Gamma(V,\mathscr P_k^{(i)})$, où $V$ parcourt les intersections (affines) de
$j+1$ ensembles du recouvrement $\mathfrak U^{(i)}$; si
$V'=u_i^{-1}(V)$, $V'$ est affine dans $X'^{(i)}$ et il résulte de (6.4.1.1)
que l'on a
\[
  \Gamma\bigl(V',\mathscr{T}\!or_{q_i}^S
  (\mathscr P_k^{(i)},\mathscr O_{S'})\bigr)
  =\operatorname{Tor}_{q_i}^A(\Gamma(V,\mathscr P_k^{(i)}),A')
\]
d'où pour le bicomplexe
$H_{q_i}^{\mathrm{II}}(M_{\bullet\bullet\bullet}'^{(i)})$ l'expression
\[
  C^\bullet\bigl(\mathfrak U'^{(i)},
  \mathscr{T}\!or_{q_i}^S(\mathscr P_\bullet^{(i)},\mathscr O_{S'})\bigr),
\]
ce qui donne finalement l'expression cherchée pour le terme $E_2$ de la
suite $(f')$. Le fait que cette suite soit birégulière sous les conditions
indiquées se vérifie comme d'ordinaire, tenant compte de ce que, si
$\mathscr P_\bullet^{(i)}$ est limité inférieurement, tous les degrés de
$M_{\bullet\bullet\bullet}'^{(i)}$ sont limités inférieurement.
\end{proof}

\begin{remark}[6.9.5]
On voit comme dans (6.7.6) que le remplacement des
$\mathscr P_\bullet^{(i)}$ et de $\mathscr Q'_\bullet$ par des complexes qui
leur sont respectivement homotopes ne change pas les suites $(e)$, $(f)$ et
$(f')$ à un isomorphisme canonique près. En outre, pour la suite $(f)$, des
homomorphismes
$\mathscr P_\bullet^{(i)}\to\mathscr R_\bullet^{(i)}$,
$\mathscr Q'_\bullet\to\mathscr T'_\bullet$ de complexes qui donnent des
isomorphismes en homologie
$\mathscr H_\bullet(\mathscr P_\bullet^{(i)})\xrightarrow{\sim}
\mathscr H_\bullet(\mathscr R_\bullet^{(i)})$,
$\mathscr H_\bullet(\mathscr Q'_\bullet)\xrightarrow{\sim}
\mathscr H_\bullet(\mathscr T'_\bullet)$ fournissent un isomorphisme de suites
spectrales
\[
\begin{split}
  {}^{(f)}\mathscr E
  (f_1,\ldots,f_m;
   \mathscr P_\bullet^{(1)},\ldots,\mathscr P_\bullet^{(m)},
   \mathscr Q'_\bullet)
  \xrightarrow{\sim}\;&
  {}^{(f)}\mathscr E
  (f_1,\ldots,f_m;\\
  &\mathscr R_\bullet^{(1)},\ldots,\mathscr R_\bullet^{(m)},
   \mathscr T'_\bullet);
\end{split}
\]
la démonstration est la même que pour (6.7.6) en tenant compte du résultat
de (6.7.6) et de la régularité de la suite $(f)$.
\end{remark}

\begin{corollary}[6.9.6]
Sous les conditions de (6.9.1), supposons que :
\begin{enumerate}
  \item[$1^\circ$] Les complexes $\mathscr P_\bullet^{(i)}$ sont formés de
  Modules plats sur $S$, et les $\mathscr O_Y$-Modules
  \[
    \mathfrak{Tor}_n^S
    (f_1,\ldots,f_m;
     \mathscr P_\bullet^{(1)},\ldots,\mathscr P_\bullet^{(m)})
  \]
  sont plats sur $S$.
  \item[$2^\circ$] Les $\mathscr P_\bullet^{(i)}$ et
  $\mathscr Q'_\bullet$ sont limités inférieurement.
\end{enumerate}

\oldpage[III]{169}
On a alors, en posant
$\mathscr P_\bullet'^{(i)}=
\mathscr P_\bullet^{(i)}\otimes_{\mathscr O_S}\mathscr O_{S'}$, des
isomorphismes canoniques fonctoriels
\[
\begin{split}
  &\mathfrak{Tor}_n^{S'}
  (f'_1,\ldots,f'_m,1_{S'};
   \mathscr P_\bullet'^{(1)},\ldots,
   \mathscr P_\bullet'^{(m)},\mathscr Q'_\bullet)
  \xrightarrow{\sim}\\
  &\qquad
  \bigoplus_{n'+n''=n}
  \mathfrak{Tor}_{n'}^S
  (f_1,\ldots,f_m;
   \mathscr P_\bullet^{(1)},\ldots,
   \mathscr P_\bullet^{(m)})
  \otimes_{\mathscr O_S}\mathscr H_{n''}(\mathscr Q'_\bullet).
\end{split}
\tag{6.9.6.1}\label{III.6.9.6.1.fr}
\]
En particulier, pour $\mathscr Q'_\bullet$ réduit à un seul terme
$\mathscr F'$ de degré $0$, on a des isomorphismes canoniques fonctoriels
\[
\begin{split}
  &\mathfrak{Tor}_n^{S'}
  (f'_1,\ldots,f'_m,1_{S'};
   \mathscr P_\bullet'^{(1)},\ldots,
   \mathscr P_\bullet'^{(m)},\mathscr F')
  \xrightarrow{\sim}\\
  &\qquad
  \mathfrak{Tor}_n^S
  (f_1,\ldots,f_m;
   \mathscr P_\bullet^{(1)},\ldots,
   \mathscr P_\bullet^{(m)})
  \otimes_{\mathscr O_S}\mathscr F'
\end{split}
\tag{6.9.6.2}\label{III.6.9.6.2.fr}
\]
et plus particulièrement, pour $\mathscr F'=\mathscr O_{S'}$,
\[
\begin{split}
  &\mathfrak{Tor}_n^{S'}
  (f'_1,\ldots,f'_m;
   \mathscr P_\bullet'^{(1)},\ldots,
   \mathscr P_\bullet'^{(m)})
  \xrightarrow{\sim}\\
  &\qquad
  \mathfrak{Tor}_n^S
  (f_1,\ldots,f_m;
   \mathscr P_\bullet^{(1)},\ldots,
   \mathscr P_\bullet^{(m)})
  \otimes_{\mathscr O_S}\mathscr O_{S'}.
\end{split}
\tag{6.9.6.3}\label{III.6.9.6.3.fr}
\]
\end{corollary}

\begin{proof}
L'hypothèse de platitude sur les Modules composant les
$\mathscr P_\bullet^{(i)}$ entraîne que les complexes
$\mathscr{T}\!or_{q_i}^S(\mathscr P_\bullet^{(i)},\mathscr O_{S'})$ sont nuls
pour $q_i\ne0$ (6.5.8). La suite $(f')$ est donc dégénérée; l'hypothèse
$2^\circ$ entraîne d'ailleurs qu'elle est birégulière (6.9.3), donc le
edge-homomorphisme
\[
\begin{split}
  &\mathfrak{Tor}_n^S
  (f_1,\ldots,f_m,1_{S'};
   \mathscr P_\bullet^{(1)},\ldots,
   \mathscr P_\bullet^{(m)},\mathscr Q'_\bullet)
  \longrightarrow {}^{(f')}\mathscr E_{n0}^2\\
  &\qquad=
  \mathfrak{Tor}_n^{S'}
  (f'_1,\ldots,f'_m,1_{S'};
   \mathscr P_\bullet'^{(1)},\ldots,
   \mathscr P_\bullet'^{(m)},\mathscr Q'_\bullet)
\end{split}
\tag{6.9.6.4}\label{III.6.9.6.4.fr}
\]
est bijectif (0, 11.1.6). L'hypothèse de platitude sur les Modules
\[
  \mathfrak{Tor}_n^S
  (f_1,\ldots,f_m;
   \mathscr P_\bullet^{(1)},\ldots,\mathscr P_\bullet^{(m)})
\]
entraîne que ${}^{(e)}\mathscr E_{pq}^2=0$ pour $p\ne0$ (6.5.8). La suite
$(e)$ est donc aussi dégénérée, et comme elle est birégulière, le
edge-homomorphisme
\[
\begin{split}
  &\mathfrak{Tor}_n^S
  (f_1,\ldots,f_m,1_{S'};
   \mathscr P_\bullet^{(1)},\ldots,
   \mathscr P_\bullet^{(m)},\mathscr Q'_\bullet)
  \longrightarrow {}^{(e)}\mathscr E_{0n}^2\\
  &\qquad=
  \bigoplus_{n'+n''=n}
  \mathfrak{Tor}_{n'}^S
  (f_1,\ldots,f_m;
   \mathscr P_\bullet^{(1)},\ldots,
   \mathscr P_\bullet^{(m)})
  \otimes_{\mathscr O_S}\mathscr H_{n''}(\mathscr Q'_\bullet)
\end{split}
\tag{6.9.6.5}\label{III.6.9.6.5.fr}
\]
est bijectif (0, 11.1.6); d'où, en combinant les deux isomorphismes
précédents, l'isomorphisme (6.9.6.1). L'isomorphisme (6.9.6.2) s'en déduit
trivialement, puisque l'on a alors
$\mathscr H_q(\mathscr Q'_\bullet)=0$ si $q\ne0$ et
$\mathscr H_0(\mathscr Q'_\bullet)=\mathscr F'$. Enfin, le cas
$\mathscr F'=\mathscr O_{S'}$ dans (6.9.6.2) donne l'isomorphisme (6.9.6.3),
compte tenu de (6.7.12).
\end{proof}

\begin{corollary}[6.9.7]
Sous les conditions de (6.9.1), supposons $S$ et $S'$ affines, et supposons
donné pour chaque $i$ un entier $d_i$ ($1\leq i\leq m$). Il existe alors un
entier $N$ ne dépendant que de $S$, des $X^{(i)}$ et des $d_i$, ayant la
propriété suivante : pour tout entier $n_0$, on a des isomorphismes canoniques
(6.9.6.3) pour $n\leq n_0$ et pour tout système de complexes
$\mathscr P_\bullet^{(i)}$ vérifiant les conditions suivantes :
\begin{enumerate}
  \item[$1^\circ$] $\mathscr P_k^{(i)}=0$ pour $k<d_i$;
  \item[$2^\circ$] $\mathscr P_k^{(i)}$ est plat sur $S$ pour $k<n_0+N$;
  \item[$3^\circ$]
  $\mathfrak{Tor}_q^S(f_1,\ldots,f_m;
  \mathscr P_\bullet^{(1)},\ldots,\mathscr P_\bullet^{(m)})$ est plat sur
  $S$ pour $q<n_0+N$.
\end{enumerate}
\end{corollary}

\begin{proof}
Supposons $\mathscr P_k^{(i)}$ plat sur $S$ pour $k<r$; alors
$\mathscr{T}\!or_{q_i}^S(\mathscr P_k^{(i)},\mathscr O_{S'})=0$ pour $k<r$
et $q_i\ne0$; calculons
\[
\begin{split}
  \mathfrak{Tor}_p^{S'}\Bigl(
  &f'_1,\ldots,f'_m;
  \mathscr{T}\!or_{q_1}^S(\mathscr P_\bullet^{(1)},\mathscr O_{S'}),
  \ldots,\\
  &\mathscr{T}\!or_{q_m}^S(\mathscr P_\bullet^{(m)},\mathscr O_{S'})
  \Bigr).
\end{split}
\tag{6.9.7.1}\label{III.6.9.7.1.fr}
\]

\oldpage[III]{170}
par la méthode de (6.6.2), à l'aide de l'image réciproque d'un recouvrement
affine fixe $\mathfrak U^{(i)}$ de $X^{(i)}$ ($1\leq i\leq m$)
(indépendant de $S'$ et des $\mathscr P_\bullet^{(i)}$); les termes
$L_{jk}^{(i)}$ sont nuls pour $j<-N_i$ (ne dépendant que de
$\mathfrak U^{(i)}$); si l'un des $q_i$ n'est pas nul, le complexe simple dont
l'homologie de degré $p$ est (6.9.7.1) a ses termes nuls pour tous les degrés
\[
  <r+\sum_{j\ne i}d_j-\sum_{i=1}^m N_i,
\]
donc (6.9.7.1) est nul pour $p<r-N$, en désignant par $N$ le plus grand des
nombres
\[
  \sum_{i=1}^m N_i-\sum_{j\ne i}d_j.
\]
On en conclut que l'on a ${}^{(f')}\mathscr E_{pq}^2=0$ pour $q\ne0$ et
$p<r-N$; comme d'autre part ${}^{(f')}\mathscr E_{pq}^2=0$ pour $q<0$, on
voit que le edge-homomorphisme (6.9.6.4) est bijectif pour $n<r-N$
(M, XV, 5.6) (pour $\mathscr Q'_\bullet=\mathscr O_{S'}$). En second lieu, si
\[
  \mathfrak{Tor}_q^S
  (f_1,\ldots,f_m;
   \mathscr P_\bullet^{(1)},\ldots,\mathscr P_\bullet^{(m)})
\]
est plat sur $S$ pour $q<r$, on a ${}^{(e)}\mathscr E_{pq}^2=0$ pour $p\ne0$
et $q<r$; par ailleurs ${}^{(e)}\mathscr E_{pq}^2=0$ pour $p<0$, donc le
edge-homomorphisme (6.9.6.5) est bijectif pour $n<r$, ce qui achève la
démonstration.
\end{proof}

Le cas le plus important de (6.9.3) dans les applications est celui où $m=1$,
$\mathscr Q'_\bullet$ étant réduit à un seul terme $\mathscr F'$ de degré
$0$; nous l'énoncerons à nouveau dans ce cas en vue de références
ultérieures\footnote{Le cas traité dans (6.9.8), et en particulier les suites
spectrales (6.9.8.3) et (6.9.8.4), nous avaient été signalés en 1957 par
J.-P. Serre.} :

\begin{proposition}[6.9.8]
Soient $S$ un préschéma, $g:S'\to S$ un morphisme,
$f:X\to Y$ un $S$-morphisme séparé et quasi-compact de $S$-préschémas,
$\mathscr P_\bullet$ un complexe de $\mathscr O_X$-Modules quasi-cohérents
limité inférieurement, $\mathscr F'$ un $\mathscr O_{S'}$-Module
quasi-cohérent. Il existe deux foncteurs spectraux biréguliers en
$\mathscr P_\bullet$ et $\mathscr F'$, à valeurs dans la catégorie des
$\mathscr O_{Y_{(S')}}$-Modules quasi-cohérents, ayant même aboutissement
$\mathfrak{Tor}_\bullet^S(f,1_{S'};\mathscr P_\bullet,\mathscr F')$, et dont
les termes $E_2$ sont
\[
  {}'\mathscr E_{pq}^2
  =\mathscr{T}\!or_p^S
  (\mathscr H^{-q}(f,\mathscr P_\bullet),\mathscr F')
  \tag{6.9.8.1}\label{III.6.9.8.1.fr}
\]
et
\[
  {}''\mathscr E_{pq}^2
  =\mathscr H^{-p}
  (f',\mathscr{T}\!or_q^S(\mathscr P_\bullet,\mathscr F')),
  \tag{6.9.8.2}\label{III.6.9.8.2.fr}
\]
où $f'=f_{(S')}:X_{(S')}\to Y_{(S')}$. 
\end{proposition}

\begin{proof}
Les suites en question peuvent aussi s'obtenir, non en partant de (6.9.3),
mais des suites $(a)$ et $(b')$ de (6.7.3) pour
$X^{(1)}=X$, $Y^{(1)}=Y$, $X^{(2)}=Y^{(2)}=S'$, $f_1=f$,
$f_2=1_{S'}$. Lorsque $S=S'=Y$, $Y$ étant affine, on obtient deux suites
spectrales de termes $E_2$ égaux à
\[
  {}'\mathscr E_{pq}^2
  =\mathscr{T}\!or_p^Y
  (\mathscr H^{-q}(f,\mathscr P_\bullet),\mathscr F)
  \tag{6.9.8.3}\label{III.6.9.8.3.fr}
\]
\[
  {}''\mathscr E_{pq}^2
  =\mathscr H^{-p}
  (f,\mathscr{T}\!or_q^Y(\mathscr P_\bullet,\mathscr F))
  \tag{6.9.8.4}\label{III.6.9.8.4.fr}
\]
aboutissant (en vertu de (6.7.6)) à l'hypercohomologie
$\mathscr H^\bullet(f,\mathscr P_\bullet\otimes_Y\mathscr F)$ du foncteur
$f_*$ par rapport au complexe
$\mathscr P_\bullet\otimes_Y\mathscr F$ de $\mathscr O_X$-Modules, pour tout
$\mathscr O_Y$-Module quasi-cohérent et $Y$-plat $\mathscr F$ (ou pour tout
$\mathscr O_Y$-Module quasi-cohérent $\mathscr F$ lorsque
$\mathscr P_\bullet$ est formé de $\mathscr O_X$-Modules $Y$-plats), qui sont
distinctes de celles de (6.2.1).
\end{proof}

\begin{corollary}[6.9.9]
Sous les conditions de (6.9.8), supposons que le complexe
$\mathscr P_\bullet$ soit limité inférieurement, formé de Modules plats sur
$S$, et que les $\mathscr O_Y$-Modules
$\mathscr H^n(f,\mathscr P_\bullet)$ soient plats sur $S$.

\oldpage[III]{171}
On a alors des isomorphismes canoniques fonctoriels
\[
  \mathfrak{Tor}_n^{S'}
  (f',1_{S'};\mathscr P'_\bullet,\mathscr F')
  \xrightarrow{\sim}
  \mathscr H^{-n}(f,\mathscr P_\bullet)
  \otimes_{\mathscr O_S}\mathscr F'
  \tag{6.9.9.1}\label{III.6.9.9.1.fr}
\]
où
$\mathscr P'_\bullet=\mathscr P_\bullet\otimes_{\mathscr O_S}
\mathscr O_{S'}$; en particulier, pour $\mathscr F'=\mathscr O_{S'}$, on a
des isomorphismes canoniques fonctoriels
\[
  \mathscr H^n(f',\mathscr P'_\bullet)
  \xrightarrow{\sim}
  \mathscr H^n(f,\mathscr P_\bullet)
  \otimes_{\mathscr O_S}\mathscr O_{S'}.
  \tag{6.9.9.2}\label{III.6.9.9.2.fr}
\]
C'est le cas particulier $m=1$ de (6.9.6). Plus particulièrement :
\end{corollary}

\begin{corollary}[6.9.10]
Soient $S$ un préschéma, $f:X\to Y$ un $S$-morphisme séparé et
quasi-compact de $S$-préschémas, $\mathscr P_\bullet$ un complexe limité
inférieurement, formé de $\mathscr O_X$-Modules quasi-cohérents, plats sur
$S$. On suppose en outre que les $\mathscr O_Y$-Modules
$\mathscr H^n(f,\mathscr P_\bullet)$ soient plats sur $S$. Pour tout $s\in S$,
notons $X_s$ et $Y_s$ les fibres
$X\otimes_S k(s)$, $Y\otimes_S k(s)$,
$f_s:X_s\to Y_s$ le morphisme $f\times_S1$,
$\mathscr P_\bullet^s$ le complexe
$\mathscr P_\bullet\otimes_{\mathscr O_S}k(s)$ de
$\mathscr O_{X_s}$-Modules. On a alors des isomorphismes canoniques
fonctoriels
\[
  \mathscr H^n(f_s,\mathscr P_\bullet^s)
  \xrightarrow{\sim}
  \mathscr H^n(f,\mathscr P_\bullet)\otimes_{\mathscr O_S}k(s).
  \tag{6.9.10.1}\label{III.6.9.10.1.fr}
\]
\end{corollary}

On a donc, moyennant des hypothèses de platitude convenables, un cas où la
formation des foncteurs dérivés $R^nf_*(\mathscr F)$ « commute au passage aux
fibres », que nous retrouverons par une autre méthode au §~7.

\subsection{Structure locale de certains foncteurs cohomologiques.}

\begin{proposition}[6.10.1]
Soient $S=\operatorname{Spec}(A)$ un schéma affine,
$Y^{(i)}$ ($1\leq i\leq n$) une famille finie de $S$-schémas affines, plats
sur $S$; pour chaque $i$, soit $f_i:X^{(i)}\to Y^{(i)}$ un $S$-morphisme
séparé et quasi-compact, et soit $\mathscr P_\bullet^{(i)}$ un complexe de
$\mathscr O_{X^{(i)}}$-Modules quasi-cohérents limité inférieurement. Soit
$Y$ le produit des $S$-schémas $Y^{(i)}$. Il existe un complexe
$\mathscr R_\bullet$ de $\mathscr O_Y$-Modules quasi-cohérents et plats sur
$S$, ayant la propriété suivante : pour tout $S$-schéma affine $S'$ et tout
complexe de $\mathscr O_{S'}$-Modules quasi-cohérents
$\mathscr Q'_\bullet$ limité inférieurement, il y a un isomorphisme
\[
  \mathbf{Tor}_\bullet^S
  (f_1,\ldots,f_n,1_{S'};
   \mathscr P_\bullet^{(1)},\ldots,\mathscr P_\bullet^{(n)},
   \mathscr Q'_\bullet)
  \xrightarrow{\sim}
  \mathscr H_\bullet
  (\mathscr R_\bullet\otimes_{\mathscr O_S}\mathscr Q'_\bullet)
  \tag{6.10.1.1}\label{III.6.10.1.1.fr}
\]
qui est un isomorphisme de $\partial$-foncteurs en $\mathscr Q'_\bullet$. En
outre, pour tout $S$-morphisme $u:S''\to S'$ de $S$-schémas affines, le
diagramme
\[
\xymatrix@C=12pt@R=30pt{
  {\begin{gathered}
    \mathbf{Tor}_\bullet^S
    (f_1,\ldots,f_n,1_{S'};\\
    \mathscr P_\bullet^{(1)},\ldots,\mathscr P_\bullet^{(n)},
    \mathscr Q'_\bullet)
  \end{gathered}}
  \ar[r]^-{\sim}\ar[d]
  &
  {\mathscr H_\bullet
    (\mathscr R_\bullet\otimes_{\mathscr O_S}\mathscr Q'_\bullet)}
  \ar[d]
  \\
  {\begin{gathered}
    \mathbf{Tor}_\bullet^S
    (f_1,\ldots,f_n,1_{S''};\\
    \mathscr P_\bullet^{(1)},\ldots,\mathscr P_\bullet^{(n)},
    u^*(\mathscr Q'_\bullet))
  \end{gathered}}
  \ar[r]^-{\sim}
  &
  {\mathscr H_\bullet
    (\mathscr R_\bullet\otimes_{\mathscr O_S}u^*(\mathscr Q'_\bullet))}
}
\tag{6.10.1.2}\label{III.6.10.1.2.fr}
\]
(où les flèches verticales sont les $(1_Y\times_Su)$-morphismes canoniques
définis dans (6.7.10)) est commutatif.
\end{proposition}

\begin{proof}
\oldpage[III]{172}
Calculons les hypertor par la méthode de (6.6.2), en tenant compte de la remarque
(6.6.6); avec les notations de (6.6.2), chaque
$L_{\bullet\bullet}^{(i)}$ est un bicomplexe de $A_i$-modules, en désignant par
$A_i$ l'anneau de $Y^{(i)}$; il admet donc une résolution de Cartan--Eilenberg
projective $M_{\bullet\bullet\bullet}^{(i)}$ (au sens de (0, 11.7.1)) formée de
$A_i$-modules, et en vertu de (6.6.6), le premier membre de (6.10.1.1) est
canoniquement isomorphe à
\[
  \mathscr H_\bullet
  (M_{\bullet\bullet\bullet}\otimes_A Q'_\bullet),
\]
où
\[
  M_{\bullet\bullet\bullet}
  =
  M_{\bullet\bullet\bullet}^{(1)}
  \otimes_A M_{\bullet\bullet\bullet}^{(2)}
  \otimes_A\cdots\otimes_A M_{\bullet\bullet\bullet}^{(n)}
\]
et $\mathscr Q'_\bullet=\widetilde{Q'_\bullet}$. Comme, par hypothèse, les
anneaux $A_i$ sont des $A$-modules plats, les
$M_{\bullet\bullet\bullet}^{(i)}$ sont des tricomplexes de $A$-modules plats
($0_{\mathrm I}$, 6.2.1), et il en est de même de
$M_{\bullet\bullet\bullet}$; en outre, si $B$ est l'anneau de $Y$, produit
tensoriel des $A_i$, $M_{\bullet\bullet\bullet}$ est un tricomplexe de
$B$-modules; le complexe
\[
  \mathscr R_\bullet=(M_{\bullet\bullet\bullet})^\sim
\]
de $\mathscr O_Y$-Modules (où $M_{\bullet\bullet\bullet}$ est considéré comme
complexe simple) répond donc à la question, comme il résulte aisément de
(6.7.10).
\end{proof}

\begin{corollary}[6.10.2]
Dans l'énoncé de (6.10.1), on peut supposer $\mathscr R_\bullet$ limité
inférieurement. Lorsque les $\mathscr P_\bullet^{(i)}$ sont limités
supérieurement et les $Y^{(i)}$ de dimension cohomologique finie, on peut
supposer $\mathscr R_\bullet$ limité supérieurement.
\end{corollary}

\begin{proof}
La première assertion résulte de ce que les trois degrés de chacun des
$M_{\bullet\bullet\bullet}^{(i)}$ sont limités inférieurement; d'autre part, si
les anneaux $A_i$ sont de dimension cohomologique finie, le troisième degré de
chacun des $M_{\bullet\bullet\bullet}^{(i)}$ ne prend qu'un nombre fini de
valeurs, et il en est de même par construction de son premier degré (6.6.2);
comme son second degré est limité supérieurement s'il en est ainsi du degré de
$\mathscr P_\bullet^{(i)}$ (6.6.2), la seconde assertion en résulte aussitôt.
\end{proof}

\begin{remarks}[6.10.3]
\textnormal{(i)} Avec les notations de (6.10.1),
$\mathscr H_\bullet(\mathscr R_\bullet\otimes_S\mathscr Q'_\bullet)$ est
isomorphe à
$\mathbf{Tor}_\bullet(\mathscr R_\bullet,\mathscr Q'_\bullet)$ puisque
$\mathscr R_\bullet$ est formé de $\mathscr O_Y$-Modules $S$-plats (6.5.9); il
est donc (6.5.4) l'aboutissement d'une suite spectrale régulière de terme
$E_2$ donné par
\[
  \mathscr E_{pq}^2
  =
  \bigoplus_{q'+q''=q}
  \mathbf{Tor}_p^S
  \bigl(
    \mathscr H_{q'}(\mathscr R_\bullet),
    \mathscr H_{q''}(\mathscr Q'_\bullet)
  \bigr)
  \tag{6.10.3.1}\label{III.6.10.3.1.fr}
\]
qui n'est autre que la suite spectrale $(e)$ du changement de base (6.9.3).

\medskip
\textnormal{(ii)} Soit $\mathscr R'_\bullet$ un second complexe de
$\mathscr O_Y$-Modules quasi-cohérents, plat sur $S$, et soit
$g:\mathscr R'_\bullet\to\mathscr R_\bullet$ un homomorphisme de complexes tel
que
\[
  \mathscr H_\bullet(g):
  \mathscr H_\bullet(\mathscr R'_\bullet)
  \longrightarrow
  \mathscr H_\bullet(\mathscr R_\bullet)
\]
soit un isomorphisme. Alors, en vertu de (6.3.3) et (6.5.9), on déduit de $g$
un isomorphisme de $\partial$-foncteurs en $\mathscr Q'_\bullet$
\[
  \mathscr H_\bullet
  (\mathscr R'_\bullet\otimes_S\mathscr Q'_\bullet)
  \xrightarrow{\sim}
  \mathscr H_\bullet
  (\mathscr R_\bullet\otimes_S\mathscr Q'_\bullet)
\]
tel que le diagramme
\[
\xymatrix@C=28pt@R=30pt{
  \mathscr H_\bullet
  (\mathscr R'_\bullet\otimes_S\mathscr Q'_\bullet)
  \ar[r]^-{\sim}\ar[d]
  &
  \mathscr H_\bullet
  (\mathscr R_\bullet\otimes_S\mathscr Q'_\bullet)
  \ar[d]
  \\
  \mathscr H_\bullet
  (\mathscr R'_\bullet\otimes_Su^*(\mathscr Q'_\bullet))
  \ar[r]^-{\sim}
  &
  \mathscr H_\bullet
  (\mathscr R_\bullet\otimes_Su^*(\mathscr Q'_\bullet))
}
\]
\par\medskip\noindent
soit commutatif. Cela prouve donc que le complexe $\mathscr R_\bullet$ n'est
pas entièrement déterminé par les propriétés de (6.10.1).

\medskip
\oldpage[III]{173}
\textnormal{(iii)} Dans la démonstration de (6.10.1), on peut supposer les
$M_{\bullet\bullet\bullet}^{(i)}$ formés de $A_i$-modules libres (comme il
résulte aisément de la démonstration de (0, 11.5.2.1) « dualisée »); les
\[
  M_{\bullet\bullet\bullet}^{(i)}\otimes_{A_i}B
\]
sont alors formés de $B$-modules libres, et comme
$M_{\bullet\bullet\bullet}$ est égal à leur produit tensoriel sur $B$, on voit
qu'on peut supposer en outre dans (6.10.1) que $\mathscr R_\bullet$ est associé
à un complexe de $B$-modules libres. En outre, en vertu de (M, XVII, 1.2), le
tricomplexe $M_{\bullet\bullet\bullet}$ dépend fonctoriellement de chacun des
bicomplexes $L_{\bullet\bullet}^{(i)}$ (donc de chacun des
$\mathscr P_\bullet^{(i)}$, lorsqu'on fixe un recouvrement fini de chacun des
$X^{(i)}$), les « morphismes » de tricomplexes devant être ici entendus comme
les classes d'homomorphismes pour la relation d'homotopie; d'ailleurs, le
remplacement d'un recouvrement de $X^{(i)}$ par un recouvrement plus fin
donnant lieu pour les $L_{\bullet\bullet}^{(i)}$ à des homomorphismes définis
précisément à homotopie près (6.6.8), on voit finalement qu'avec la convention
précédente pour les morphismes, le tricomplexe
$M_{\bullet\bullet\bullet}$ est un foncteur en chacun des
$\mathscr P_\bullet^{(i)}$. Nous préciserons cette dépendance fonctorielle, et
notamment le comportement de $\mathscr R_\bullet$ relativement à des suites
exactes
\[
  0\longrightarrow\mathscr P_\bullet^{\prime(i)}
  \longrightarrow\mathscr P_\bullet^{(i)}
  \longrightarrow\mathscr P_\bullet^{\prime\prime(i)}
  \longrightarrow0
\]
de complexes, dans le chapitre consacré à une algèbre générale de foncteurs
cohomologiques, mentionné dans (6.1.3).
\end{remarks}

\begin{scholium}[6.10.4]
Le fait que $\mathscr R_\bullet$ est formé de $\mathscr O_Y$-Modules
$S$-plats entraîne aisément que
\[
  \mathscr H_\bullet
  (\mathscr R_\bullet\otimes_S\mathscr Q'_\bullet)
\]
est un foncteur homologique en $\mathscr Q'_\bullet$ (voir le raisonnement de
(7.7.1)). C'est cette propriété qui, ainsi qu'on l'a mentionné en (6.1.1), est
la motivation de l'introduction de l'hypertor. Posons en effet
\[
  X=X^{(1)}\times_SX^{(2)}\times_S\cdots\times_SX^{(n)},
  \qquad
  f=f_1\times_Sf_2\times_S\cdots\times_Sf_n,
\]
\[
  \mathscr P_\bullet
  =
  \mathscr P_\bullet^{(1)}
  \otimes_S\mathscr P_\bullet^{(2)}
  \otimes_S\cdots\otimes_S\mathscr P_\bullet^{(n)},
\]
\[
  X'=X\times_SS',
  \qquad
  Y'=Y\times_SS',
  \qquad
  f'=f\times_S1_{S'}.
\]
Les problèmes de changement de base amènent à étudier l'hypercohomologie
\[
  \mathscr H^\bullet_{f'}
  (\mathscr P_\bullet\otimes_S\mathscr N')
\]
en tant que foncteur par rapport au $\mathscr O_{S'}$-Module quasi-cohérent
$\mathscr N'$, ou encore l'hypercohomologie
\[
  \mathscr H^\bullet_f
  (\mathscr P_\bullet\otimes_S\mathscr N)
\]
comme foncteur en le $\mathscr O_S$-Module quasi-cohérent $\mathscr N$. Lorsque
les $\mathscr P_\bullet^{(i)}$ (donc aussi $\mathscr P_\bullet$) sont
$S$-plats, il résulte de ce qui précède et de (6.7.6) que ce foncteur est bien
un foncteur cohomologique en $\mathscr N$; mais il n'en est plus de même
lorsqu'on ne fait plus l'hypothèse de platitude sur les
$\mathscr P_\bullet^{(i)}$, et l'on ne peut plus alors aborder l'étude de
\[
  \mathscr H^\bullet_f
  (\mathscr P_\bullet\otimes_S\mathscr N)
\]
par les méthodes usuelles de l'Algèbre homologique.

Nous aurons toutefois surtout à utiliser le cas où $n=1$, $Y=S$ et où
$\mathscr P_\bullet$ est formé de $\mathscr O_X$-Modules $Y$-plats. On a dans
ce cas le
\end{scholium}

\begin{theorem}[6.10.5]
Soient $Y=\operatorname{Spec}(A)$ un schéma affine noethérien,
$f:X\to Y$ un morphisme propre, $\mathscr P_\bullet$ un complexe de
$\mathscr O_X$-Modules cohérents, plats sur $Y$, limité inférieurement. Il
existe alors un complexe $\mathscr L_\bullet$ de $\mathscr O_Y$-Modules,
limité inférieurement, dont les termes $\mathscr L_i$ sont des
$\mathscr O_Y$-Modules de la forme $\mathscr O_Y^{m_i}$, et un isomorphisme
\[
  \mathscr H^\bullet
  (f,\mathscr P_\bullet\otimes_Y\mathscr Q_\bullet)
  \xrightarrow{\sim}
  \mathscr H_\bullet
  (\mathscr L_\bullet\otimes_Y\mathscr Q_\bullet)
  \tag{6.10.5.1}\label{III.6.10.5.1.fr}
\]
de $\partial$-foncteurs en le complexe $\mathscr Q_\bullet$ de
$\mathscr O_Y$-Modules quasi-cohérents, limité inférieurement. En outre, pour
tout morphisme $u:Y'\to Y$, on a, en posant
\[
  X'=X_{(Y')},
  \qquad
  f'=f_{(Y')},
  \qquad
  \mathscr P'_\bullet
  =
  \mathscr P_\bullet\otimes_Y\mathscr O_{Y'},
  \qquad
  \mathscr L'_\bullet=u^*(\mathscr L_\bullet),
\]
\oldpage[III]{174}
(qui est un complexe de $\mathscr O_{Y'}$-Modules localement libres de type
fini), un isomorphisme
\[
  \mathscr H^\bullet
  (f',\mathscr P'_\bullet\otimes_{Y'}\mathscr Q'_\bullet)
  \xrightarrow{\sim}
  \mathscr H_\bullet
  (\mathscr L'_\bullet\otimes_{Y'}\mathscr Q'_\bullet)
  \tag{6.10.5.2}\label{III.6.10.5.2.fr}
\]
de $\partial$-foncteurs en le complexe $\mathscr Q'_\bullet$ de
$\mathscr O_{Y'}$-Modules quasi-cohérents, limité inférieurement, de sorte que
le diagramme
\[
\xymatrix@C=25pt@R=30pt{
  \mathscr H^\bullet
  (f,\mathscr P_\bullet\otimes_Y\mathscr Q_\bullet)
  \ar[r]^-{\sim}\ar[d]
  &
  \mathscr H_\bullet
  (\mathscr L_\bullet\otimes_Y\mathscr Q_\bullet)
  \ar[d]
  \\
  \mathscr H^\bullet
  (f',\mathscr P'_\bullet\otimes_{Y'}u^*(\mathscr Q_\bullet))
  \ar[r]^-{\sim}
  &
  \mathscr H_\bullet
  (\mathscr L'_\bullet\otimes_{Y'}u^*(\mathscr Q_\bullet))
}
\tag{6.10.5.3}\label{III.6.10.5.3.fr}
\]
soit commutatif.
\end{theorem}

\begin{proof}
L'application de (6.10.1) donne d'abord un complexe limité inférieurement
(6.10.2) $\mathscr R_\bullet$ de $\mathscr O_Y$-Modules quasi-cohérents et
$Y$-libres (6.10.3, (iii)) et un isomorphisme
\[
  \mathscr H^\bullet
  (f,\mathscr P_\bullet\otimes_Y\mathscr Q_\bullet)
  \xrightarrow{\sim}
  \mathscr H_\bullet
  (\mathscr R_\bullet\otimes_Y\mathscr Q_\bullet)
  \tag{6.10.5.4}\label{III.6.10.5.4.fr}
\]
de $\partial$-foncteurs en $\mathscr Q_\bullet$, mais \emph{a priori} les
termes de $\mathscr R_\bullet$ ne sont pas nécessairement des
$\mathscr O_Y$-Modules de type fini. Mais si l'on applique (6.10.5.4) au cas
où $\mathscr Q_\bullet$ est un complexe réduit à un seul terme
$\mathscr O_Y$, on voit que $\mathscr H_\bullet(\mathscr R_\bullet)$ est
isomorphe à $\mathscr H^\bullet(f,\mathscr P_\bullet)$, et est par suite formé
de $\mathscr O_Y$-Modules cohérents (6.2.5). On sait alors (0, 11.9.2) qu'il
existe un complexe $\mathscr L_\bullet$ limité inférieurement, formé de
$\mathscr O_Y$-Modules associés à des $A$-modules libres de type fini, et un
homomorphisme
\[
  \mathscr L_\bullet\longrightarrow\mathscr R_\bullet,
\]
tels que l'homomorphisme correspondant pour l'homologie
\[
  \mathscr H_\bullet(\mathscr L_\bullet)
  \longrightarrow
  \mathscr H_\bullet(\mathscr R_\bullet)
\]
soit bijectif; d'où l'isomorphisme (6.10.5.1), en vertu de (6.10.3, (ii)).
Les autres assertions de (6.10.5) résultent de (6.10.1) et (6.10.3, (ii))
lorsque $Y'$ est affine; dans le cas général, il suffit de vérifier que
lorsqu'on considère un recouvrement $(V_\alpha)$ de $Y'$ par des ouverts
affines, et l'isomorphisme correspondant (6.10.5.2) relatif à chacun des
$V_\alpha$, les restrictions à un ouvert affine
$W\subset V_\alpha\cap V_\beta$ des isomorphismes correspondant à
$V_\alpha$ et à $V_\beta$ coïncident avec l'isomorphisme correspondant à
$W$, ce qui résulte de la commutativité du diagramme (6.10.1.2) appliqué aux
injections canoniques $W\to V_\alpha$ et $W\to V_\beta$.
\end{proof}

\begin{remark}[6.10.6]
Dans les chapitres suivants, nous appliquerons surtout (6.10.5) au cas où
$\mathscr P_\bullet$ est réduit à un seul $\mathscr O_X$-Module cohérent
$\mathscr F$, plat sur $Y$. Comme on a alors
\[
  \mathscr H_n(\mathscr L_\bullet)
  =
  \mathscr H^{-n}(f,\mathscr F)
  =
  R^{-n}f_*(\mathscr F)
  \qquad(6.2.1),
\]
on voit que les $\mathscr H_n(\mathscr L_\bullet)$ sont nuls pour $n>0$;
nous verrons plus loin (7.7.12, (i)) qu'on peut alors supposer que
$\mathscr L_\bullet$ n'a que des termes de degrés $\leqslant0$ (donc en
nombre fini), à condition de remplacer l'hypothèse que les $\mathscr L_i$
sont associés à des $A$-modules libres de type fini par celle que les
$\mathscr L_i$ sont localement libres de type fini.

Le complexe $\mathscr L_\bullet$ correspondant à un tel
$\mathscr O_X$-Module $\mathscr F$ ne paraît posséder aucune
\oldpage[III]{175}
propriété particulière, en dehors de la restriction précédente sur les
degrés. On peut alors se demander si inversement, étant donné un complexe
$\mathscr L_\bullet$ formé de $\mathscr O_Y$-Modules associés à des
$A$-modules projectifs de type fini, limité inférieurement et dont les termes
de degré $>0$ sont nuls, il existe un $Y$-schéma $X$, projectif et plat sur
$Y$ et un $\mathscr O_X$-Module localement libre $\mathscr F$, tels qu'il y
ait un isomorphisme
\[
  \mathscr H^\bullet
  (f,\mathscr F\otimes_Y\mathscr Q_\bullet)
  \xrightarrow{\sim}
  \mathscr H_\bullet
  (\mathscr L_\bullet\otimes_Y\mathscr Q_\bullet)
\]
fonctoriel en $\mathscr Q_\bullet$. L'intérêt d'un tel résultat serait de
réduire complètement la théorie cohomologique des Modules cohérents et
$Y$-plats sur des $Y$-schémas propres, à la théorie « à homotopie près » des
complexes de $A$-modules projectifs de type fini sur un anneau noethérien
$A$.
\end{remark}
